\documentclass[11pt,a4paper]{article}
\usepackage[T1]{fontenc}
\usepackage[utf8]{inputenc}
\usepackage{lmodern}
\ifdefined\pdfmapfile
  \pdfmapfile{=lm.map}
  \pdfmapfile{+cm.map}
  \pdfmapfile{+symbols.map}
  \pdfmapfile{+euler.map}
  \pdfmapfile{+rsfs.map}
\fi
\usepackage[a4paper,margin=27mm]{geometry}
\usepackage{amsmath,amssymb,amsthm,mathtools,mathrsfs}
\usepackage{booktabs,tabularx,longtable,array,enumitem}
\usepackage{microtype,xcolor,etoolbox}
\usepackage{hyperref}
\hypersetup{colorlinks=true,linkcolor=blue!45!black,citecolor=blue!45!black,
urlcolor=blue!45!black,pdfauthor={Research record prepared with ChatGPT},
pdftitle={Bounded classical curvature kernels and arbitrary-background potential quotients}}
\usepackage[capitalise,noabbrev]{cleveref}
\numberwithin{equation}{section}
\newtheorem{theorem}{Theorem}[section]
\newtheorem{lemma}[theorem]{Lemma}
\newtheorem{proposition}[theorem]{Proposition}
\newtheorem{corollary}[theorem]{Corollary}
\theoremstyle{definition}
\newtheorem{definition}[theorem]{Definition}
\newtheorem{example}[theorem]{Example}
\theoremstyle{remark}
\newtheorem{remark}[theorem]{Remark}
\newcommand{\Cb}{\mathcal C_b}
\newcommand{\R}{\mathbf R}
\newcommand{\Q}{\mathbf Q}
\newcommand{\Z}{\mathbf Z}
\newcommand{\PSH}{\operatorname{PSH}}
\newcommand{\BTcap}{\operatorname{Cap}}
\newcommand{\supp}{\operatorname{supp}}
\newcommand{\dd}{d'd''}
\newcommand{\an}{\mathrm{an}}
\setlist{topsep=4pt,itemsep=2pt,parsep=0pt}
\setlength{\parindent}{1.3em}
\setlength{\parskip}{2pt}
\setlength{\emergencystretch}{2em}
\makeatletter
\patchcmd{\l@section}{1.5em}{2.1em}{}{}
\renewcommand*\l@subsection{\@dottedtocline{2}{1.5em}{3.6em}}
\makeatother

\title{Bounded classical curvature kernels\\
and arbitrary-background potential quotients\\[5pt]
\large Fifth research record in non-Archimedean Bedford--Taylor theory}
\author{Research record prepared with ChatGPT\\
\small Relative to the authorized manuscripts R, B, and A2}
\date{10 November 2026 (UTC)}
\begin{document}
\maketitle
\begin{abstract}
Let $K/\mathbf Q_p$ be finite with its prescribed absolute value.
On an ordinary open in a good Hausdorff boundaryless quasi-smooth
$K$-analytic space, we prove that the actual difference of two
locally bounded classical plurisubharmonic functions with equal
original CLD curvature is continuous and classically pluriharmonic.
It follows that their relative curvatures agree for every fixed
symmetric original-positive separately $d'$- and $d''$-closed
background current. The mixed bounded Bedford--Taylor calculus
therefore descends in its differentiated slots to currents with
bounded classical local potentials.

The proof supplies the two gates left open in the fourth record:
kernel preservation after every complete scalar extension and
ordinary continuity of the actual curvewise harmonic difference.
The first uses prescribed smooth rows, finite-etale traces, full
analytic fiber potentials including type IV, and a simultaneous
product induction. The second uses a Shilov bound proved without
ambient semicontinuity, polynomial Noether coordinates, and exact
logarithmic interpolation on strictly semistable alteration charts.
All equalities are tested on the full original CLD test sheaf.
Explicit scope tests exclude weaker closedness, arbitrary bounded
upper-semicontinuous representatives, and an unrestricted passage
to relative energy or capacitary completions. The previously
proved seven fixed-background BT properties are retained as inputs.
An exact global finite-dimensional curvature hypothesis yields
local uniform compactness of normalized classical potentials.
Without that hypothesis, a two-intersection example shows that
the relative classical image can fail to be closed in total
variation even when no boundary values are prescribed.
\end{abstract}

\section*{Reading guide}
\addcontentsline{toc}{section}{Reading guide}
The exact main theorem and its short final proof are in
Section~\ref{sec:r5-arbitrary-background-quotient}. The new analytic
and geometric arguments are in
Sections~\ref{sec:r5-scalar-extension-kernel}--\ref{sec:r5-algebraic-continuity-geometry}.
Section~\ref{sec:r5-inputs} declares the inputs and the point-value
and test-form conventions. Section~\ref{sec:r5-status} compares the
current status with R4. The later examples distinguish false
extensions from the stated theorem.

This is a focused continuation. It contains the new proofs rather
than reprinting all 304 pages of R4. The complete archive retains
the exact frozen R4 archive, its original internal inputs, all
modular and stand-alone R5 LaTeX sources, source-location records,
independent reading records, and the continuation handoff.
References beginning ``R4--'' are numbered results in that frozen
final record. Historical open statements in R4 are superseded only
in the precise scopes stated here.

\tableofcontents
\clearpage

% BEGIN canonical R5 source: sections/round5_inputs.tex
\section{The question, conventions, and frozen inputs}
\label{sec:r5-inputs}

The question left open in the fourth record was
\begin{equation}
 D[u]=D[v]\quad\Longrightarrow\quad D((u-v)S)=0,
 \qquad u,v\in\mathcal C_b(U),
 \label{eq:r5-main-question}
\end{equation}
for an arbitrary fixed symmetric positive separately closed
background current $S$. The equality on the left is an equality
of original CLD currents with background $[1]$. It must not be
silently interpreted as equality after multiplication by a
singular coefficient measure. The actual point values of $u$
and $v$ matter on the right.

The original ground field throughout the main theorem is a finite
extension $K/\mathbf Q_p$ with its given absolute value. An
\emph{allowed space} is good, Hausdorff, boundaryless, quasi-smooth,
and of pure dimension. All opens are ordinary topological opens.
We set $D=d'd''$. A $(p,p)$-current on a space of dimension $n$
acts on original compact smooth $(n-p,n-p)$-forms. Unless a
different degree is explicitly stated, positivity is the original
positive cone of R, namely the symmetric quadratic-form cone of
\cite[Definition 5.1.1(c) and \S5.4]{CLD12}. It is not identified
with either strong or weak positivity.

The coefficient-Radon arguments use this middle cone directly.
They do not invoke the strong-test order-zero assertion as printed
in \cite[Proposition 7.2.9]{CLD25}. Section
\ref{sec:r5-weak-order-zero-tropical} gives an explicit closed
real-tropical counterexample to that broader assertion and checks
the distinct 2012 hypotheses. This does not change R's stipulated
cone or any of the authorized internal inputs.

The word \emph{actual} means that no function is replaced on a
null, polar, or coefficient-null set. For a locally bounded Borel
function $a$ and a coefficient-Radon current $T$, multiplication
means
\[
             \langle aT,\eta\rangle
                  =\int a\,d\mu_{T,\eta}
\]
on the actual Berkovich space. A measure specified only on one
tropical image does not give this operation.

\subsection{Authorized internal results}

\begin{description}[style=nextline]
\item[\texttt{INTERNAL\_INPUT}: R.]
R provides the continuous original-positive local calculus over
the original field, its continuous semipositive smoothing,
the original positive-cone coefficient-Radon controls, and
compatibility of weak convergence of original-positive currents
with a fixed continuous multiplier. Its statements use the original test
category. Its local algebraic tube construction is used only
over the original field, after B's algebraic localization.

\item[\texttt{INTERNAL\_INPUT}: B.]
B, Definition 2.2, defines $\mathcal C_b(U)$ by actual local
boundedness, upper semicontinuity, and subharmonic pullbacks
along every analytic curve after every complete isometric field
extension. Theorem 2.5 gives decreasing continuous rational
semipositive PL rows on one fixed ordinary neighborhood.
Proposition 4.5 and the pure-Fubini--Study smoothing used in
its construction give the prescribed smooth positive rows
needed here. Explicitly, if $p_j\downarrow u$ and
$|s_j-p_j|\leq\varepsilon_j$ with
$\varepsilon_{j+1}\leq\varepsilon_j/4$, then
$s_j+2\varepsilon_j\downarrow u$ because its values lie
between $p_j+\varepsilon_j$ and $p_j+3\varepsilon_j$.
The original smooth curvatures stay positive. The fixed domain,
rather than a separate shrinking domain for every $j$, is
essential. B, Lemma D.4, supplies ordinary affine algebraic
neighborhoods. B, Appendix A.2, supplies the positive Radon
Laplacian and its original scalar-test identity on smooth
algebraic curves over every complete nontrivially valued field
of characteristic zero; that particular curve statement has
the larger field range explicitly stated in B.

\item[\texttt{INTERNAL\_INPUT}: A2.]
The same-model geometric adelic nefness result is already
complete under its stated arithmetic and boundary assumptions.
It is not used to prove the local quotient and is not counted
as an open problem or a result of this round.
\end{description}

No comprehensive proof audit of R, B, or A2 has been undertaken.
The checks concern their statements, fields, positivity cones,
test categories, and dependencies.

\subsection{Results proved in the earlier project records}

The frozen R4 record contains the full R1--R3 sources and has
304 pages in its final stand-alone version. References of the
form ``R4--$a.b$'' below refer to that final pagination and
numbering, not to the older intermediate modular build. These
results retain their original project provenance; they are not
new theorems of R5.

\begin{longtable}{@{}p{0.22\textwidth}p{0.51\textwidth}p{0.20\textwidth}@{}}
\toprule
Input & Exact content used & Location in final R4\\
\midrule
\endhead
Fixed-background BT & All seven bounded-potential properties,
including weighted joint decreasing convergence, for one
arbitrary original-positive separately closed $S$. &
Theorem 3.7 and preceding construction; p.~18 onward\\
Common smooth rows & Prescribed smaller norm boxes admit
simultaneous smooth decreasing positive rows for finitely
many original-$K$ classical potentials. & Lemma 7.1, p.~24\\
Compatible piecewise tests & Tangentially compatible original
piecewise smooth tests admit uniform pairing approximation
against mass-bounded positive closed first-order currents. &
Theorem 71.4, p.~240\\
Original-field analytic kernel & $D[u]=D[v]$ remains true
after every analytic map of allowed original-$K$ spaces;
in particular after every finite polydisc projection. &
Theorem 75.3, p.~253\\
Exact fiber test spreading & Every compact original test on
an algebraic parameter fiber is the exact restriction of
a compact original test on an ambient parameter neighborhood. &
Lemma 77.2, p.~256\\
Finite-coefficient descent & Every finite family of original
tests after $L/K$ descends exactly to
$\overline{K(c_1,\ldots,c_m)}\subset L$ for a finite tuple. &
Lemma 78.1, p.~260\\
Energy and completions & The faithful form, complete-pole
splitting, and bounded relative capacitary closure are used
in the energy consequences and scope tests, not in the proof of (Q). &
Theorems 17.6, 18.5, 84.2--84.3, 86.1--86.2
in the frozen record\\
\bottomrule
\end{longtable}

The initial statement $D[u]=D[v]$ is never pulled back as an
arbitrary current. Each transfer theorem proves equality for
the actual composed potentials, on every original test on the
new space. The distinction is indispensable.

\subsection{Logical structure of the new argument}

The analytic part extends the prescribed-row finite-etale
kernel argument to auxiliary complete fields. A continuous
potential for the \emph{full} analytic fiber of every nonrigid
point of the affine line then closes the scalar-extension
gate, including type IV. The product kernel is propagated
simultaneously in every finite product dimension. The
resulting analytic-map kernel gives actual harmonicity on
all curve pullbacks, including their type-I and type-IV values.

The geometric part does not assume that a curvewise harmonic
function is upper semicontinuous. It proves its Shilov bound
by finite flat coordinates and normalized algebraic curve
slices. Strictly semistable alteration charts supply continuous
logarithmic interpolation of their finite Shilov data. The
Shilov bound identifies the actual function with those
interpolants, and proper compact descent proves continuity.

Only after that argument are both signs of $u-v$ known to be
classical psh. Applying the existing first-order positive
construction to the arbitrary background $S$ makes
$D((u-v)S)$ and its negative positive, which proves
\eqref{eq:r5-main-question}. The quotient is therefore
downstream of all the kernel and continuity arguments.

\paragraph{Status convention.}
\texttt{NEW\_PROVED} means that a complete argument is supplied
in this project relative to the declared inputs. It does not
claim worldwide priority, external refereeing, or formal
verification. It is distinct from \texttt{DIRECT\_COROLLARY},
\texttt{INTERNAL\_INPUT}, and \texttt{EXTERNAL\_KNOWN}.
The logical statuses are \texttt{PROVED},
\texttt{CONDITIONAL}, \texttt{OPEN}, and \texttt{FALSE};
an invalid proof mechanism is labeled \texttt{FAILED\_ATTEMPT}.


% END source: sections/round5_inputs.tex


% BEGIN canonical R5 source: sections/prescribed_rows_field_package.tex
\section{Field-independent tools for prescribed smooth positive rows}
\label{sec:r5-prescribed-rows-package}

Let $F$ be a complete non-Archimedean field with its given absolute
value.  Let $U$ be a good Hausdorff boundaryless $F$-analytic space of
pure dimension $n\geq1$.  We use only the original CLD smooth forms,
their canonical integration, and $D=d'd''$.  A current is positive
in the original positive cone used in R, and all currents called
positive below are symmetric.  This is not a convention identifying
the positive, strongly positive and weakly positive cones.

A \emph{prescribed smooth positive row} for a function $u$ on an
ordinary open $O$ means original smooth functions $u_j$ on that
same $O$, with positive $(1,1)$-forms $Du_j$, such that $u_j$
decreases pointwise to the actual finite locally bounded function
$u$.  Whenever a continuous potential is used below, its smooth
positive approximation is likewise part of the hypotheses.  No
regularization theorem over $F$ is asserted.  In particular, these
results may be applied to rows already obtained over the original
finite extension of $\mathbf Q_p$ and then pulled back by an
analytic map or by scalar extension.

\subsection{Positive coefficient measures over the given field}

\begin{lemma}[Finite positive controls and actual coefficient measures]
\label{lem:r5-field-independent-radon}
\label{lem:r5-prescribed-coefficient-control}
Let $T$ be an original positive current of bidegree $(p,p)$ on $O$,
and put $q=n-p$.  For every original compact smooth complementary
test $\eta$, there is a signed Radon measure $\mu_{T,\eta}$ on $O$
such that
\[
 \langle T,a\eta\rangle=\int a\,d\mu_{T,\eta}
\]
for every original smooth scalar $a$.  The assignment is compatible
with finite localization, relations between original forms, and
common original moment refinements.

There are finitely many original positive compact smooth tests
$\Theta_\ell$ and a constant $C_\eta$, depending only on a fixed
presentation of $\eta$ and its compact support stage, such that
\[
 |\mu_{T,\eta}|\leq
 C_\eta\sum_\ell\mu_{T,\Theta_\ell}.
 \tag{R}
\]
The same controls work after adjoining additional scalar moments.
Consequently a fixed locally bounded Borel multiplier $a$ defines
the original current $aT$.  If $T_j\to T$ weakly and all these
currents are positive, then $aT_j\to aT$ for every fixed continuous
multiplier $a$.  Weak convergence alone is not claimed to allow
an arbitrary fixed discontinuous multiplier.
\end{lemma}

\noindent\textbf{Status: \texttt{DIRECT\_COROLLARY}} of the original
finite-dimensional positive-cone controls and Radon representation.
The proof below records why this part has no restriction to a
discretely valued or locally compact ground field.

\begin{proof}
Localize the test into finitely many ordinary moment presentations
with compact support.  On one such presentation write
$t=(\log|g_1|,\ldots,\log|g_r|)$, where the $g_i$ are analytic
units, and put
\[
 \beta=\sum_{\nu=1}^r d't_\nu\wedge d''t_\nu.
\]
In the usual signed identification of symmetric $(q,q)$-tensors
with symmetric quadratic forms on $\bigwedge^q\mathbf R^r$,
$\beta^q/q!$ is the identity matrix.  On a fixed compact real
presentation stage it therefore dominates, up to a finite
constant, both signs of the symmetric coefficient tensor of
$\eta$.  With an original compact smooth cutoff $\chi$ equal
to one near that stage, this gives
\[
 C\chi\beta^q\pm\eta\geq0.
\]
The control $\chi\beta^q$ is strongly positive.  The inequalities
are in the original positive cone; they do not assume that $\eta$
is strongly positive.  Antisymmetric parts have zero pairing
with the symmetric current.

For every nonnegative original smooth scalar $b$, positivity gives
\[
 |\langle T,b\eta\rangle|
       \leq C\langle T,b\chi\beta^q\rangle .
\]
Original smooth scalar functions separate points on each good
affinoid neighborhood: the absolute values of analytic functions
separate Berkovich points, and smooth real profiles constant near
zero turn these into original smooth scalars.  Their compact
versions are uniformly dense in continuous compactly supported
functions on a fixed compact neighborhood.  This is the original
scalar Stone--Weierstrass argument, valid over every complete $F$;
see also \cite[Propositions 15.2.6--15.2.9]{CLD25}.
Positivity and the displayed control extend the functionals
uniquely to that continuous scalar space.  Riesz representation
gives the positive measure of the control and the signed measure
of $\eta$, with the displayed domination.  The finitely many
pieces give (R).

Relations between forms and changes of presentation are compatible
because the corresponding original scalar-tested functionals
are identical; uniqueness of Radon measures preserves every such
identity.  The matrix inequalities pull back under common
refinements, so adding scalar coordinates does not change the
controls.  Applying this argument to the finitely many coefficient
basis forms gives the actual coefficient-Radon description.
It also proves continuity of $aT$ in the original test topology
for locally bounded Borel $a$: on each fixed compact presentation,
the coefficient norms of a test bound the pairing by its fixed
control masses times $\sup|a|$.

For the last assertion, weak convergence bounds the finitely many
positive control masses.  Approximate the fixed continuous scalar
$a$ uniformly on their compact support by original smooth scalars.
For each smooth scalar the claimed convergence is the definition
of weak convergence.  Inequality (R) bounds the two approximation
errors uniformly in $j$.  Taking the approximation error to zero
proves the assertion.  No countable atlas or measure separability
has entered the argument.
\end{proof}

\subsection{A direct first-order estimate near an analytic zero set}

Let $H$ be analytic on $O$ and not identically zero on any
irreducible component, and put $Z=\{H=0\}$.  A fact from the
original calculus used below is
\begin{equation}
 \operatorname{supp}\omega\cap Z=\varnothing
 \quad\hbox{for every original smooth top-degree form }\omega.
 \label{eq:r5-top-form-avoids-divisor}
\end{equation}
Indeed $\dim Z<n$, and the germ of an original form of bidegree
$(n,n)$ is zero at any point whose residual transcendence dimension
$d_F(x)$ is less than $n$; see
\cite[Proposition 15.2.1]{CLD25}.  In particular $\log|H|$ is
bounded on the compact actual support of every such compact form.
This also implies its usual local integrability against original
smooth top-degree trace measures.

For smooth real functions put
\[
 \Gamma(a,b)=\tfrac12(d'a\wedge d''b+d'b\wedge d''a),
 \qquad \Gamma(a)=d'a\wedge d''a.
\]
All first-order pairings in the following proof involve smooth
functions.  No energy object for an unknown nonsmooth function
is being assumed.

\begin{lemma}[Uniform analytic-zero-set estimate for a prescribed row]
\label{lem:r5-prescribed-row-divisor-estimate}
Let $u_j\downarrow u$ be a prescribed smooth positive row on $O$.
Fix one original moment presentation near a compact set $K$,
put $\beta=\sum d't_\nu\wedge d''t_\nu$, and set
$\theta=\beta^{n-1}$, with $\theta=1$ when $n=1$.
Choose original smooth cutoffs $0\leq\chi,\chi_1\leq1$ with
$\chi=1$ near $K$ and $\chi_1=1$ near $\operatorname{supp}\chi$,
both compactly supported in that presentation.

After shrinking to a relatively compact ordinary neighborhood of
these supports, there are smooth psh functions $q_N$, $N\geq N_0$,
such that
\[
 -1\leq q_N\leq0,\qquad
 q_N\longrightarrow-\mathbf1_Z
\]
pointwise, and a constant $C$, independent of $j,N$, such that
\begin{equation}
 0\leq\int\chi(-q_N)\,Du_j\wedge\theta
                       \leq C N^{-1/2}.
 \label{eq:r5-uniform-divisor-estimate}
\end{equation}
The constant depends on the two cutoffs, the fixed moment frame,
$H$, and one bound for the row on their compact neighborhood.
\end{lemma}

\begin{proof}
Choose a constant $A$ so that
$\varphi=\log|H|-A\leq-2$ on this neighborhood.  Let $\psi$ be
a smooth convex nondecreasing real function which is zero on
$(-\infty,-1]$, equals the identity on $[1,\infty)$, and satisfies
\[
 \max(0,s)\leq\psi(s)\leq\max(0,s)+1.
\]
Set
\[
 \ell_N=-N+\psi(\varphi+N),\qquad q_N=\ell_N/N,
\]
extending $\ell_N$ by $-N$ near every zero of $H$.
These extensions are original smooth: the nonconstant part is
supported in the unit locus of $H$.  Convexity gives $Dq_N\geq0$.
For $N\geq2$ the displayed bounds and pointwise limit hold.

Whenever $\omega$ is a fixed compact original smooth top form,
\eqref{eq:r5-top-form-avoids-divisor} shows that $\varphi$ is
bounded on $\operatorname{supp}\omega$.  Hence, for all sufficiently
large $N$ depending on $\omega$,
\begin{equation}
 q_N=\varphi/N\quad\hbox{on }\operatorname{supp}\omega.
 \label{eq:r5-exact-truncation-on-top-support}
\end{equation}
Stokes against the smooth closed form $\theta$ gives
\[
 0\leq\int\chi Dq_N\wedge\theta
       =\int q_N D\chi\wedge\theta=O(N^{-1}).
\]
The constant is finite by
\eqref{eq:r5-exact-truncation-on-top-support} applied to
$D\chi\wedge\theta$.

Write $E_\chi(a)=\int\chi\Gamma(a)\wedge\theta$.
The smooth square identity and compact Stokes yield
\[
 2E_\chi(q_N)
  =\int q_N^2D\chi\wedge\theta
           -2\int\chi q_NDq_N\wedge\theta .
\]
Since $-1\leq q_N\leq0$, the second term is at most
$2\int\chi Dq_N\wedge\theta=O(N^{-1})$.
The absolute value of the first term is $O(N^{-2})$, again by
\eqref{eq:r5-exact-truncation-on-top-support}.  Consequently
\begin{equation}
 E_\chi(q_N)=O(N^{-1}),\qquad
 \int q_N^2\Gamma(\chi)\wedge\theta=O(N^{-2}).
 \label{eq:r5-log-cutoff-energies}
\end{equation}
The second assertion uses the same support argument for the fixed
top form $\Gamma(\chi)\wedge\theta$.

The row is uniformly bounded on the compact support stage.
Indeed $u\leq u_j\leq u_1$, $u$ is locally bounded, and $u_1$
is continuous.  Write $|u_j|\leq M$ there.  For any fixed
nonnegative compact smooth cutoff $\xi$ in that stage,
\[
 \int\xi Du_j\wedge\theta
       =\int u_j D\xi\wedge\theta
       \leq M\,|\mu_{D\xi\wedge\theta}|(O).
\]
Combining this with
\[
 2E_\xi(u_j)=\int u_j^2D\xi\wedge\theta
                    -2\int\xi u_jDu_j\wedge\theta
\]
gives a uniform bound for $E_\xi(u_j)$.  Apply it to $\chi$
and $\chi_1$.

One integration by parts, and symmetry of the remaining form,
now give the exact identity
\[
 \int\chi(-q_N)Du_j\wedge\theta
   =\int\chi\Gamma(q_N,u_j)\wedge\theta
          +\int q_N\Gamma(\chi,u_j)\wedge\theta .
\]
The positive-square Cauchy--Schwarz inequality bounds its absolute
value by
\[
 E_\chi(q_N)^{1/2}E_\chi(u_j)^{1/2}
 +\left(\int q_N^2\Gamma(\chi)\wedge\theta\right)^{1/2}
                              E_{\chi_1}(u_j)^{1/2}.
\]
Here $\chi_1=1$ near $\operatorname{supp}d\chi$ justifies the
last energy factor.  Equation \eqref{eq:r5-log-cutoff-energies}
proves \eqref{eq:r5-uniform-divisor-estimate}.  All integrations
occur inside one original moment presentation with compact
cutoffs, so no boundary term or unproved global exhaustion is used.
\end{proof}

\begin{theorem}[First-order nonconcentration for prescribed rows]
\label{thm:r5-prescribed-row-analytic-nonmass}
Let $u$ have a prescribed smooth positive row locally on $U$.
Then $D[u]$ is an original positive separately closed
coefficient-Radon current.  If an analytic function $H$ does not
vanish identically on any local irreducible component, every
coefficient measure of $D[u]$ gives zero mass to $\{H=0\}$.
In particular this holds on every proper analytic divisor.
All assertions are on the full original test sheaf over $F$.
\end{theorem}

\noindent\textbf{Status: \texttt{NEW\_PROVED}} for the stated
prescribed-row first-order extension.  The conclusion is not an
assertion of decreasing regularization for all classical functions
over $F$.

\begin{proof}
Bounded convergence on each compact support against the fixed
smooth top-form measures gives $[u_j]\to[u]$.  Hence
$Du_j\to D[u]$.  Positivity and separate closedness pass to
this limit, and Lemma~\ref{lem:r5-field-independent-radon}
gives its coefficient measures.

In the notation of the preceding lemma, fix $N$.  The form
$\chi(-q_N)\theta$ is an original smooth compact test.  Weak
convergence passes \eqref{eq:r5-uniform-divisor-estimate} to
\[
 \int\chi(-q_N)\,D[u]\wedge\theta\leq C N^{-1/2}.
\]
Now apply bounded convergence to the ONE fixed finite positive
measure $\chi D[u]\wedge\theta$.  Since
$-q_N\to\mathbf1_Z$, it follows that
\[
 \bigl(\chi D[u]\wedge\theta\bigr)(Z)=0.
\]
For an arbitrary complementary test on this stage, the finite
positive control construction dominates the variation of its
coefficient measure by a constant times
$\chi D[u]\wedge\beta^{n-1}$, with a cutoff equal to one
on its support.  Thus every coefficient measure vanishes on $Z$.
Finite localization handles an arbitrary compact test.  First
$j$ and then $N$ were sent to infinity; no joint discontinuous
limit is used.
\end{proof}

\subsection{One relative slot and exchange with a continuous slot}

\begin{proposition}[Relative first-order construction from a given row]
\label{prop:r5-prescribed-relative-first-order}
Let $S$ be any original positive separately closed $(p,p)$-current,
$p<n$, and let $u$ have a prescribed smooth positive row.
Then
\[
 D(uS)=\lim_j Du_j\wedge S
\]
is an original positive separately closed coefficient-Radon
current.  Its definition is independent of the prescribed row.
\end{proposition}

\begin{proof}
Lemma~\ref{lem:r5-field-independent-radon} defines the actual
Borel product $uS$.  Since $u_j$ is locally uniformly bounded
and converges pointwise, bounded convergence for the fixed
coefficient measures of $S$ gives $u_jS\to uS$.
Distributional differentiation gives the displayed convergence.
Every $Du_j\wedge S$ is positive and separately closed.
Apply Lemma~\ref{lem:r5-field-independent-radon} to its limit.
The defining expression $D(uS)$ depends only on the fixed actual
function $u$, so it is independent of the row.
\end{proof}

\begin{proposition}[Exchange with a prescribed continuous potential]
\label{prop:r5-prescribed-continuous-exchange}
Let $S$ be as above, with $p\leq n-2$.  Let $u$ have a
prescribed smooth positive row.  Let $v$ be finite continuous
and admit locally either a uniformly convergent smooth positive
approximation or a prescribed smooth positive decreasing row.
Then, on every original compact complementary test,
\[
 D\bigl(uD(vS)\bigr)=D\bigl(vD(uS)\bigr).
\]
\end{proposition}

\noindent\textbf{Status: \texttt{DIRECT\_COROLLARY}} for both
propositions.  They use no general regularization over $F$ and
no joint decreasing continuity theorem for two discontinuous slots.

\begin{proof}
The relative first-order currents exist by the preceding proposition;
for a uniform approximation of $v$, the same proof uses uniform
convergence instead of bounded convergence.  For every smooth
$u_j$, the distributional product rule and separate closedness give
\[
 D\bigl(u_jD(vS)\bigr)
    =Du_j\wedge D(vS)
    =D\bigl(v(Du_j\wedge S)\bigr).
\]
On the left, bounded convergence uses the fixed coefficient measures
of $D(vS)$ and yields $D(uD(vS))$.  On the right,
$Du_j\wedge S\to D(uS)$ through positive currents, and the
fixed multiplier $v$ is continuous.  The last assertion of
Lemma~\ref{lem:r5-field-independent-radon} therefore gives
the limit $D(vD(uS))$.  These are two justified one-parameter
passages, not an interchange of discontinuous limits.
\end{proof}

\paragraph{Scope of the package.}
The field-independent ingredients are finite positive matrix
controls, compact smooth scalar density, original Stokes, and the
vanishing of top-degree form germs on proper analytic subsets.
The only approximations of the potentials are the prescribed rows.
These results supply the first-order analytic-divisor removal and
continuous-slot exchange used in an original-test projection
argument.  They do not establish a bounded mixed calculus for
arbitrary unknown potentials over $F$, the arbitrary-background
potential quotient, or a new pullback operation for arbitrary
currents.

% END source: sections/prescribed_rows_field_package.tex


% BEGIN canonical R5 source: sections/prescribed_finite_etale_kernel.tex
\section{Finite etale kernel transfer for prescribed smooth rows}
\label{sec:r5-prescribed-finite-etale}

Let $F$ be a complete nontrivially valued field of characteristic zero.
Its absolute value is fixed. All analytic spaces in this section are
good, Hausdorff, boundaryless, and quasi-smooth of pure dimension.
All forms and currents are in the original CLD sheaf. In particular,
``positive'' denotes the positive semidefinite cone, with no
replacement by the weakly smooth or delta-form sheaf.

\begin{definition}[A prescribed smooth row]
\label{def:r5-prescribed-row}
An actual locally bounded classical psh function $a$ on an ordinary
open $U$ is said to have a prescribed smooth row if every point has
one ordinary neighborhood $V$ and original smooth functions $a_j$
on $V$ such that $a_j\downarrow a$ pointwise and $Da_j$ is a positive
$(1,1)$-form. The neighborhoods do not depend on $j$.
The word ``classical'' retains the upper-semicontinuity and
all-complete-extension curve conditions. Those conditions are not
deduced from the current equation in this definition.
\end{definition}

The functions occurring in the original problem have these rows by
the authorized regularization in B, followed by smoothing its pure
Fubini--Study presentations on a common smaller domain. Explicitly,
if $p_j\downarrow a$ are the continuous semipositive weights, choose
smooth positive weights $s_j$ with $|s_j-p_j|\leq\epsilon_j$,
where $\epsilon_{j+1}\leq\epsilon_j/4$. Then
$s_j+3\epsilon_j$ decrease and have limit $a$. The same fixed rows
pull back to every complete extension and every analytic morphism:
smooth pullback commutes with $D$, preserves positive $(1,1)$-forms,
and preserves pointwise decrease. Classicality of the limit is
already supplied by the original all-extension definition. This is
a \texttt{DIRECT\_COROLLARY} about specified rows; it is not a
regularization theorem for all classical psh functions over $F$.

\begin{lemma}[Piecewise tests over an arbitrary complete field]
\label{lem:r5-prescribed-piecewise-tests}
Let $\eta$ be a compactly supported finite original piecewise smooth
$(q,q)$-test: in a finite unit-moment presentation it is piecewise
smooth on a finite affine hyperplane arrangement, smooth up to each
closed chamber, and its tangential restrictions agree on all common
faces. If symmetric positive separately closed currents $T_j$ of complementary
bidegree converge to a symmetric positive separately closed current $T$ on all
original smooth tests, then
\[
 \langle T_j,\eta\rangle\longrightarrow\langle T,\eta\rangle.
\]
The piecewise pairing uses the actual coefficient Radon measures and
is independent of the chamber representatives chosen on the walls.
\end{lemma}

\begin{proof}
The proof is finite dimensional once the field-independent
coefficient control of Lemma~\ref{lem:r5-prescribed-coefficient-control}
has been supplied. We give the quantitative steps, so that no
arithmetic hypothesis from R4 is carried over implicitly.

Write the unit logarithms as $t_1,\ldots,t_m$, put
$\beta=\sum_i d't_i\wedge d''t_i$, and enlarge the presentation
to include the finitely many cutoffs. For an affine wall equation
$\ell$, the smooth convex function
$g_\epsilon=(\ell^2+\epsilon^2)^{1/2}$ satisfies
\[
 Dg_\epsilon\geq 2^{-3/2}\epsilon^{-1}
       \mathbf1_{\{|\ell|\leq\epsilon\}}d'\ell\wedge d''\ell.
\]
Compact Stokes, separate closedness, and the fixed coefficient
controls imply, on each compact stage $H$,
\[
 (d'\ell\wedge d''\ell\wedge\beta^{q-1}\wedge T_j)
       (H\cap\{|\ell|\leq\epsilon\})
       \leq C\epsilon M_j,
 \qquad M_j=\sum_b\langle T_j,\Theta_b\rangle.
\]
Here the $\Theta_b$ are finitely many fixed positive compact smooth
tests, and $\sup_jM_j<\infty$ by weak convergence. The same bound
holds for $T$. Matrix Cauchy--Schwarz for the original positive cone
therefore bounds every coefficient with one primed or double-primed
normal factor on that tube by $C\epsilon^{1/2}M_j$.

Convolve the ambient chamber tensors by a radius-$\epsilon$ real
smooth kernel, then multiply by the fixed compact cutoff. This gives
an original smooth $\eta_\epsilon$. On any one wall, tangential
matching writes the branch difference as
\[
 \ell\theta+d'\ell\wedge\xi+d''\ell\wedge\zeta
\]
with uniformly bounded coefficients. For finitely many walls,
subdivide a segment joining two points at distance at most
$\epsilon$ by its wall crossings. The smooth variations along its
chambers sum to $O(\epsilon)$; each jump is in the corresponding
normal ideal, and its wall is at distance $O(\epsilon)$ from the
initial point. A fixed linear right inverse for each normal ideal
gives uniform coefficient bounds. Degenerate segments follow by
limits. Thus
\[
 |\langle T_j,\eta_\epsilon-\eta\rangle|
       \leq C\epsilon^{1/2}M_j.
\]
The normal-factor estimate with decreasing tubes also makes every
change of chamber representative on a wall invisible to the pairing.
First pass $j\to\infty$ with the smooth $\eta_\epsilon$ fixed,
then let $\epsilon\downarrow0$. This proves the assertion for
$q\geq1$. For $q=0$, matching means continuity and ordinary uniform
convolution approximation suffices. A finite original partition
handles all presentations. Each inequality is preserved on further
original refinements, because it is an inequality of pulled-back
positive tensors. No valuation-group or residue-field property was
used.
\end{proof}

\begin{lemma}[Finite etale trace in the original piecewise class]
\label{lem:r5-prescribed-finite-trace}
Let $p:Z\to U$ be finite etale between the spaces above. The
pushforward current of any compact original smooth form on $Z$ is
represented locally on $U$ by finite original compatible piecewise
smooth forms. The identity holds on every original complementary
test and every further original moment refinement.
\end{lemma}

\begin{proof}
The construction of Theorem~\textup{R4--73.3}
in R4 uses only the following
operations, all valid over the present $F$. We recall the construction
and the calibration rather than assuming a field-extension rule for
currents.

Properness and finiteness allow one to shrink an ordinary base
neighborhood until the source is a finite disjoint union of pieces
carrying single moment presentations. Treat one such piece of rank
$d$, with source units $h_1,\ldots,h_m$ and form
$\omega(\log|h_1|,\ldots,\log|h_m|)$. For each
$a\in\{0,\ldots,d\}^m$, take the characteristic polynomial of
$h^a=\prod_i h_i^{a_i}$ in the finite locally free analytic algebra:
\[
 Q_a(T)=T^d+c_{a,1}T^{d-1}+\cdots+c_{a,d}.
\]
The roots over an algebraic closure of each completed residue field
give the $d$ geometric sheets with their full multiplicities. Their
ordered logarithmic norms are the Newton-polygon slopes of $Q_a$.
On a fixed compact neighborhood all these norms are bounded by
some $B$. A coefficient vanishing at the chosen base point can be
omitted after shrinking until its norm is less than
$\exp(-dB-1)$: it cannot lie on the lower Newton hull. The other
coefficients are original analytic units after a further shrinking.
Consequently the ordered root norms are finite continuous piecewise
affine functions of finitely many original base unit logarithms.

These records determine the unordered vectors
$z_l=(\log|h_i|)_i$, including multiplicities. Indeed, the records
determine
$\prod_{l=1}^d(X-a\cdot z_l)$ on the grid
$\{0,\ldots,d\}^m$; polynomial interpolation determines this
polynomial in all its variables, and its monic linear factorization
determines the vectors. All possible joint listings are obtained by
permuting coordinate lists. Consistency is a finite system of linear
equalities and inequalities. Thus on finitely many polyhedral cells
one may list the vectors as affine maps $z_l(t)$ and set
\[
 \tau=\sum_{l=1}^d z_l^*\omega.
\]
On common faces the unordered lists agree. A finite further
subdivision fixes their matching permutations, so the forms have
equal tangential restrictions. This construction initially gives a
form on a finite closed polyhedral locus in real moment space.
Triangulate this locus together with a larger box and barycentrically
subdivide. On the neighborhood where the sum of hat functions of
its vertices is greater than $1/2$, their normalized barycentric
combination retracts piecewise smoothly onto the locus. Pullback
by this retraction extends $\tau$ to a full real neighborhood of
each compact moment image, preserving tangential compatibility.
The resulting form is therefore an original piecewise test, not
merely a form on an abstract skeleton.

For the current identity, use the finite-flat degree formula for
calibrated skeletons: \cite[Lemma~3.15 and Proposition~3.19,
especially equation~(3.20)]{Mih24}. Those results are
stated over an arbitrary complete non-Archimedean field. Their
local weights are finite algebra degrees; they are not replaced
by a product of value-group and residue-field degrees. Hence
possible valued-field defect causes no omitted factor. Adjoining
the source units does not increase graded transcendence degree,
by their characteristic equations. Omitted coefficients are still
strictly below the maximal terms of each initial equation, so the
retained unit coefficients suffice for this assertion. The degree
formula therefore gives, for every original complementary test
$\gamma$ on the base,
\[
 \int_Z\omega(\log|h|)\wedge p^*\gamma
      =\int_U\tau\wedge\gamma.
\]
It remains valid after every extra base moment refinement. The
extension off the consistent real locus contributes nothing on
the actual calibrated cells. There is no additional wall current.
Finite partitions and Lemma~\ref{lem:r5-prescribed-piecewise-tests}
complete the construction and its independence statements.
\end{proof}

\begin{theorem}[Finite etale preservation for prescribed rows]
\label{thm:r5-prescribed-finite-etale-kernel}
Let $p:Z\to U$ be finite etale, and let $u,v$ be actual locally
bounded classical psh functions on $U$, each with prescribed smooth
rows. If $D[u]=D[v]$ on all original tests on $U$, then
\[
 D[u\circ p]=D[v\circ p]
\]
on all original tests on $Z$.
\end{theorem}

\noindent\textbf{Status: \texttt{NEW\_PROVED} field extension of
the R4 finite-etale argument, relative to the stated coefficient
lemma.} No regularization of new functions over $F$ is asserted.

\begin{proof}
Let $\alpha\in A_c^{n-1,n-1}(Z)$. Choose a finite original smooth
partition on a neighborhood of its projected compact support,
subordinate to ordinary neighborhoods carrying rows for both
$u$ and $v$. Replace $\alpha$ in turn by
$(p^*\chi)\alpha$ for the members $\chi$ of this partition.
It suffices by linearity to treat each resulting test, whose
projected support now lies in one common row neighborhood.
For such a test, still denoted by $\alpha$, put
$P=p_*[\alpha]$ and $\nu=DP=p_*[D\alpha]$. The latter is one
fixed signed Radon measure of compact support. By the preceding
lemma, $P$ is represented on a finite compact cover by compatible
original piecewise tests. Take the given bounded smooth decreasing
rows $u_j,v_j$ on the chosen neighborhood of the projected
support. Their positive separately closed curvatures converge to
the same current, by fixed smooth-test dominated convergence.
Lemma~\ref{lem:r5-prescribed-piecewise-tests} gives
\[
 \langle P,D(u_j-v_j)\rangle\longrightarrow0.
\]
Compact smooth Stokes identifies this expression with
$\int(u_j-v_j)\,d\nu$. Bounded convergence against this fixed
measure identifies its limit with
\[
 \int(u-v)\,d\nu
      =\langle D[(u-v)\circ p],\alpha\rangle.
\]
Sum the finitely many localized identities. The derivatives in
$D\alpha$ include those of the fixed partitioned test; no derivative
of a partitioned potential is taken. This proves the equality on
the arbitrary original test without assuming a global row.
\end{proof}

\begin{corollary}[Finite scalar extensions]
\label{cor:r5-prescribed-finite-field}
If $F'/F$ is finite, the same prescribed pair satisfies
$D[u_{F'}]=D[v_{F'}]$ on all original $F'$-tests.
\end{corollary}

\begin{proof}
Characteristic zero makes the scalar extension map finite etale.
Apply the theorem with $U_{F'}$ viewed as an $F$-space. Its
original analytic functions, moment tests, and their topology are
the same in both interpretations. Its integration current in the
$F$-calibration is $[F':F]$ times its integration current in the
$F'$-calibration. This is the finite-constant formula of CLD,
also checked directly in Mihatsch, Lemma~3.16. Dividing by this
single positive degree proves the assertion. The different scalar
pushforward formula has factor one; it is not being substituted
for the present same-space calibration.
\end{proof}

% END source: sections/prescribed_finite_etale_kernel.tex


% BEGIN canonical R5 source: sections/scalar_extension_kernel.tex
\section{Complete scalar extensions through one-variable fiber potentials}
\label{sec:r5-scalar-extension-kernel}

The original finite extension $K/\mathbf Q_p$ keeps its given absolute
value.  Auxiliary complete fields below always carry extending
absolute values.  Tests are original CLD forms.  The new-field input
is the prescribed-row finite-etale kernel theorem
\ref{thm:r5-prescribed-finite-etale-kernel}, with the coefficient
control of Lemma~\ref{lem:r5-prescribed-coefficient-control}.
Neither a regularization theorem for arbitrary psh functions over an
auxiliary field nor the arbitrary-background potential quotient is
used in this section.

Here a \emph{prescribed smooth row} for $u$ means a sequence of original
smooth functions $u_j\downarrow u$, locally uniformly bounded, with
positive $(1,1)$-curvature forms.  When such a row is pulled back by
scalar extension, the smooth forms remain positive: in an original
moment presentation their pullbacks are the pullbacks of the same
positive quadratic forms under linear maps of tangent spaces.
The extra moment directions introduced by a refinement cause no
change to this argument.  Thus the rows supplied over the original
$K$ are used as specified functions over later fields; no new
regularization assertion over those fields is inferred.  All row
arguments are local on fixed ordinary algebraic neighborhoods.

\subsection{Exchange with a continuous potential needs only a prescribed row}

\begin{lemma}[Continuous-slot exchange over an auxiliary field]
\label{lem:r5-prescribed-continuous-exchange}
Let $F$ be a complete non-Archimedean field, let $W$ be a good
Hausdorff boundaryless quasi-smooth $F$-analytic space, and suppose
$u$ has a prescribed smooth row on $W$.  Let $b$ be continuous and
suppose $D[b]$ is a symmetric original-positive current.  Then
\[
 D\bigl(uD[b]\bigr)=D\bigl(bD[u]\bigr).
\]
All products here are multiplication of coefficient-Radon currents
by the displayed actual Borel or continuous function.  In particular,
if $u,v$ have prescribed rows and $D[u]=D[v]$, then
\[
 D\bigl((u-v)D[b]\bigr)=0.
\]
\end{lemma}

\noindent\textbf{Status: \texttt{DIRECT\_COROLLARY}} of smooth
Leibniz, fixed-measure dominated convergence, and the prescribed-row
coefficient control.  No joint BT convergence is needed.

\begin{proof}
For a smooth row member $u_j$, separate closedness and smooth
Leibniz give
\[
 D\bigl(u_jD[b]\bigr)=Du_j\wedge D[b]
                    =D\bigl(bDu_j\bigr).
\]
For a fixed compact original test, its coefficient measures against
$D[b]$ are finite signed Radon measures.  The row has a common bound
near that support and decreases to the actual function $u$.
Dominated convergence therefore gives $u_jD[b]\to uD[b]$ on the
fixed test and on its $D$-derivative.  On the other hand
$D[u_j]\to D[u]$ distributionally.  Positivity and the coefficient
control permit multiplication of this convergence by the fixed
continuous function $b$.  Passing to the limit proves the first
identity.  Subtract the corresponding identity for $v$ to obtain
the last assertion.
\end{proof}

\subsection{A bounded continuous potential for a nonrigid parameter point}

\begin{lemma}[Full-fiber integration is continuous in the parameter]
\label{lem:r5-continuous-full-fiber-integral}
Let $F$ be complete, let $Z$ be a good Hausdorff boundaryless smooth
$F$-analytic space of pure dimension $n$, and let $A$ be an ordinary
open in $\mathbf A_F^{1,\rm an}$.  For an original compact smooth
$(n,n)$-form $\eta$ on $Z\times_F A$, put
\[
 I_\eta(a)=\int_{Z_{\mathscr H(a)}}\eta_a.
\]
Then $I_\eta$ is an ordinary continuous function on $A$.  The fiber
is the whole completed-residue-field analytic fiber and its
integration uses that field's original calibration.
\end{lemma}

\noindent\textbf{Status: \texttt{EXTERNAL\_KNOWN}} in the stated
product setting.

\begin{proof}
The projection is separated, flat, boundaryless and purely of
relative dimension $n$.  The compact support is in particular
relatively compact over $A$.  The assertion follows from
\cite[Theorem 16.3.3(d), pp.~447--450]{CLD25}, whose conclusion is
that $I_\eta$ is $G$-tropically $C^\infty$.
For clarity about the ordinary topology, on an affinoid neighborhood
of a base point a finite compact analytic $G$-cover represents such
a function by continuous functions of finitely many unit norms.
These representatives agree on intersections.  Finite closed-set
pasting on that neighborhood gives ordinary continuity.  This
argument concerns a smooth unweighted form; no continuity of an
integral with a discontinuous psh multiplier is asserted here.
\end{proof}

\begin{proposition}[A continuous full-fiber potential at every nonrigid point]
\label{prop:r5-type-four-full-fiber-potential}
Let $C$ be an algebraically closed complete non-Archimedean field,
and let $a\in\mathbf A_C^{1,\rm an}$ be a point other than a
$C$-rational point.  There is a finite continuous function $b_a$ on
$\mathbf A_C^{1,\rm an}$, locally bounded, such that for every $Z$
as in Lemma~\ref{lem:r5-continuous-full-fiber-integral},
\begin{equation}\label{eq:r5-all-nonrigid-full-fiber}
 D[b_a\circ p]=j_{a*}[1]_{Z_{\mathscr H(a)}}
\end{equation}
on all original tests of $Z\times_C\mathbf A_C^{1,\rm an}$.
The right-hand side denotes the full analytic fiber, with coefficient
one.  In particular the left-hand side is original-positive and
separately closed.  The equality remains valid after multiplication
by any fixed locally bounded Borel function on the product.
\end{proposition}

\noindent\textbf{Status: \texttt{DIRECT\_COROLLARY}.}
The type-II and type-III formula is \texttt{EXTERNAL\_KNOWN},
namely the CLD Gauss-fiber formula. The type-IV passage uses a
uniform potential limit and unweighted fiber-integral continuity.
The limiting construction is included
because using the full fiber, and only then adding a fixed Borel
multiplier, is essential in the subsequent kernel argument.

\begin{proof}
Over $C$, the classification of points of the affine line says that
a type-II or type-III point is $\zeta_{c,r}$ with $c\in C$ and $r>0$.
In that case take
\[
 b_a(t)=\log\max(r,|t-c|).
\]
Equation~\eqref{eq:r5-all-nonrigid-full-fiber} is exactly
\cite[Corollary 18.6.2, pp.~533--534]{CLD25} for the analytic function
$t-c$ and one parameter.  The additive constant $\log r$ has no
curvature.

For type IV, choose nested closed discs
$D(c_j,r_j)$ with empty intersection in $C$ representing $a$,
and put $a_j=\zeta_{c_j,r_j}$.  The radii can be chosen decreasing
to $r>0$, and $a_j\to a$ in the ordinary Berkovich topology.
If their infimum were zero, the centers would be Cauchy and
completeness would give a point in their intersection.  These
classification statements are recorded, with their proofs, in
\cite[Lemma I.2.11, Proposition I.2.12, and Remark I.2.13,
pp.~8--9]{PT20}.

Set $b_j(t)=\log\max(r_j,|t-c_j|)$.  For $\ell\geq j$, nesting
gives $|c_\ell-c_j|\leq r_j$ and $r_\ell\leq r_j$.  Outside
$D(c_j,r_j)$, the two displayed maxima for $b_j$ and $b_\ell$
agree.  Inside that disc, $b_j=\log r_j$ and
$\log r_\ell\leq b_\ell\leq\log r_j$.  Thus on the entire
analytic line,
\[
 0\leq b_j-b_\ell\leq\log(r_j/r_\ell).
\]
Consequently $b_j$ decreases uniformly to a finite continuous
function $b_a$, with
\begin{equation}\label{eq:r5-type-four-uniform-error}
 0\leq b_j-b_a\leq\log(r_j/r)\longrightarrow0.
\end{equation}
Although the individual functions grow at infinity, this uniform
error bound is global and their values are bounded on every compact
set.

For an original compact $(n,n)$-test $\eta$ on the product,
distributional differentiation and uniform convergence give
\[
 \begin{aligned}
 \langle D[b_a\circ p],\eta\rangle
 &=\lim_j\langle D[b_j\circ p],\eta\rangle\\
 &=\lim_j I_\eta(a_j)=I_\eta(a).
 \end{aligned}
\]
The middle equality is the already known Gauss-fiber identity;
the last one is Lemma~\ref{lem:r5-continuous-full-fiber-integral}.
This proves \eqref{eq:r5-all-nonrigid-full-fiber}.  Positivity and
separate closedness follow either from the positive closed currents
$D[b_j\circ p]$ or from the displayed full-fiber equality.

To add a Borel multiplier, fix $\eta$ and test the unweighted
identity also on $\chi\eta$ for every original smooth scalar
$\chi$.  This identifies the fixed signed coefficient Radon
measure of $D[b_a\circ p]$ with the pushforward of that of
$\eta_a$.  Integrate the fixed locally bounded Borel function in
this measure identity.  In particular no limit of Borel-weighted
$a_j$-fiber measures is used.
\end{proof}

\subsection{Spreading and algebraic closure with prescribed tests}

\begin{lemma}[Spreading over an auxiliary complete field]
\label{lem:r5-auxiliary-field-test-spreading}
Let $F$ be complete, let $X$ be an ordinary open in a smooth affine
$F$-variety, and let $a\in\mathbf A_F^{1,\rm an}$.  Every original
compact test $\beta$ on $X_{\mathscr H(a)}$ is exactly the restriction
of an original compact test $\alpha$ on $X\times_F A$ for some
ordinary neighborhood $A$ of $a$.  The equality commutes with $D$.
\end{lemma}

\begin{proof}
The affine-case proof of
Lemma~\textup{R4--77.2} uses only affinoid
density and original smooth cutoffs, and works over the present
$F$.  Its essential steps can be stated directly.  On finitely many
compact rational norm boxes, replace the analytic unit moments of
$\beta$ by rational moments with the same norms, using errors
smaller than their positive minimum norms.  Localize with norm
profiles constant below a positive threshold at every possible
zero; impose separate lower cutoffs on the unit moments.  There
are then finitely many coefficients in $\mathscr H(a)$.
Approximate them by values of rational functions in $F(t)$ with
denominators nonzero at $a$.  Such values are dense by the
definition of the completed residue field.  Uniform coefficient
bounds on a fixed compact affine stage preserve exactly every
active unit norm, and preserve every zero-protected profile on the
whole stage.  Thus the localized forms remain exactly equal on
the fiber.  Remove the finite denominator zero loci in an ordinary
parameter neighborhood and add compact cutoffs in $X$ and in that
neighborhood.  Their supports lie in a finite union of compact
analytic products of affinoid domains.  This gives $\alpha$ with
the stated exact restriction.  All operations were constructions
of smooth test forms, with no psh regularization over $F$.
\end{proof}

\begin{lemma}[Transfer to the completed algebraic closure]
\label{lem:r5-completed-algebraic-kernel}
Let $F$ be complete of characteristic zero, let $X$ be an ordinary
open in a smooth affine $F$-variety, and let $u,v$ be actual locally
bounded classical psh functions with prescribed rows.  Suppose
$D[u]=D[v]$.  Under the prescribed-row finite-etale
kernel theorem, the equality also holds on all original tests over
$C=\widehat{F^{\rm alg}}$.
\end{lemma}

\begin{proof}
For each finite $F'/F$, the scalar map $X_{F'}\to X$ is finite
etale when viewed over $F$.  Apply
Theorem~\ref{thm:r5-prescribed-finite-etale-kernel}; then divide by
the positive same-space calibration factor $[F':F]$ to obtain the
equality on original $F'$-tests.

Now take an original compact test $\beta$ on $X_C$.  Rationalize
and localize its finitely many unit moments as in the preceding
lemma.  Only finitely many coefficients in $C$ remain.  Approximate
all of them by elements of $F^{\rm alg}$, with errors below all
the fixed unit and profile thresholds.  They lie in a common finite
extension $F'/F$.  The resulting compact $F'$-test $\eta$ pulls
back to exactly $\beta$: active logarithms and zero-protected
profiles, hence the differential forms, are unchanged.  For
$f=u-v$, the factor-one scalar integration identity gives
\[
 \int_{X_C}f_C D\beta=\int_{X_{F'}}f_{F'}D\eta=0.
\]
This uses the fixed signed-measure form of
\cite[Proposition 15.3.1, p.~425]{CLD25}.  No limit of the actual
values of $u$ or $v$ at moving points has been taken.
\end{proof}

\subsection{One parameter and then arbitrary complete extensions}

\begin{proposition}[One-parameter transfer from a product kernel]
\label{prop:r5-one-parameter-kernel}
Let $F$ be a complete characteristic-zero field with a countable
dense subfield, and let $X$ be an ordinary open in a smooth affine
$F$-variety.  Let $u,v$ be actual locally bounded classical psh
functions with prescribed rows.  Suppose both
\[
 D[u]=D[v]\quad\hbox{on }X,
 \qquad
 D[u\circ q]=D[v\circ q]\quad\hbox{on }X\times_F\mathbf A_F^{1,\rm an},
\]
where $q$ is the projection.  Then for every
$a\in\mathbf A_F^{1,\rm an}$,
\[
 D[u_{\mathscr H(a)}]=D[v_{\mathscr H(a)}]
\]
on all original tests.  Its new-field analytic inputs are exactly
the prescribed-row coefficient control and finite-etale kernel.
\end{proposition}

\begin{proof}
Write $E=\mathscr H(a)$ and take a completed algebraic closure
$\Omega$ of $E$.  Embed an algebraic closure of $F$ in $\Omega$;
its closure is an isometric copy of $C=\widehat{F^{\rm alg}}$.
The coordinate value $a(T)\in E\subset\Omega$ defines a point
$y\in\mathbf A_C^{1,\rm an}$.  Its completed residue field
$H=\mathscr H(y)$ is the closure of $C(a(T))$ in $\Omega$ and
contains $E$ isometrically.

Apply Lemma~\ref{lem:r5-completed-algebraic-kernel} to the two
assumed kernel equalities, including the product equality.  Thus
they hold over $C$.  This field still has a countable dense
subfield: approximate finite polynomial coefficients in a
countable dense subfield of $F$, and use the separable root
approximation argument to see that its algebraic closure is dense
in $F^{\rm alg}$ and then in $C$.

If $y$ is $C$-rational, then $H=C$ and there is nothing further
to prove over $H$.  Otherwise take the continuous parameter
potential $b_y$ of
Proposition~\ref{prop:r5-type-four-full-fiber-potential}.  On
$X_C\times_C\mathbf A_C^{1,\rm an}$ the product equality and
Lemma~\ref{lem:r5-prescribed-continuous-exchange} give
\[
 D\bigl((u_C-v_C)\circ q\;D[b_y\circ p]\bigr)=0.
\]
For an arbitrary original compact complementary test $\beta$ on
$X_H$, spread it exactly to a compact original $C$-test $\alpha$
on $X_C\times_C A$, using
Lemma~\ref{lem:r5-auxiliary-field-test-spreading}.  Restrict the
last equality to this ordinary neighborhood.  The full-fiber
identity, with its fixed Borel multiplier, yields
\[
 0=\int_{X_H}(u_H-v_H)D\beta.
\]
This proves the kernel equality over $H$ in all cases.

Finally pull an arbitrary original $E$-test back to $X_H$.  The
factor-one scalar integration identity gives the same pairing on
$X_E$ as on $X_H$, which is zero.  This descends the equality to
$E$.  None of the spaces over $C$, $E$ or $H$ was regarded as a
finite-type $F$-analytic space.
\end{proof}

\begin{theorem}[Scalar kernel functoriality for every complete extension]
\label{thm:r5-all-complete-scalar-kernel}
Let $X$ be an allowed space over the original finite field $K$,
let $u,v\in\mathcal C_b(X)$, and assume $D[u]=D[v]$ on every
original $K$-test.  Then for every complete extension $L/K$,
\[
 D[u_L]=D[v_L]
\]
on every original $L$-test, with the inherited absolute value and
the actual pullback functions.
\end{theorem}

\noindent\textbf{Status: \texttt{NEW\_PROVED}}, relative to the
prescribed-row finite-etale theorem and coefficient control stated
at the beginning of this section.  The proof does not apply B's
regularization theorem over an auxiliary infinite field.

\begin{proof}
The assertion is local on $X$.  Work in an ordinary affine
algebraic neighborhood over $K$, and choose the local prescribed
smooth decreasing rows supplied by the accepted original-$K$
inputs R and B.  They and their scalar pullbacks have common bounds
on each fixed compact stage.  The curvature-form positivity noted
at the beginning ensures that these are prescribed rows over every
auxiliary field used below.

For a complete field $F/K$ and an integer $r\geq0$, let
$\mathsf P(F,r)$ denote the original-current kernel equality for
the pullbacks of $u,v$ to
$X_F\times_F\mathbf A_F^{r,\rm an}$.  The R4 same-field product
result proves $\mathsf P(K,r)$ for every $r$.  We propagate all
these statements simultaneously; this avoids presuming a new
same-field product theorem at each infinite intermediate field.

Choose finitely many elements $c_1,\ldots,c_m$ in a complete
extension of $K$, and set
\[
 F_i=\overline{K(c_1,\ldots,c_i)},\qquad F_0=K.
\]
Each $F_i$ has a countable dense subfield.  It is the completion
of $F_{i-1}(c_i)$ and therefore is the residue field of a point of
$\mathbf A_{F_{i-1}}^{1,\rm an}$.  Assume
$\mathsf P(F_{i-1},r)$ for all $r$.  For each $r$, apply
Proposition~\ref{prop:r5-one-parameter-kernel} to
$X_{F_{i-1}}\times\mathbf A_{F_{i-1}}^r$.  Its product-kernel
hypothesis is precisely $\mathsf P(F_{i-1},r+1)$.
It gives $\mathsf P(F_i,r)$.  Induction proves all $r$ statements
over every such finite-parameter field.

For an arbitrary complete $L/K$ and an arbitrary original compact
test $\beta$ on $X_L$, apply the exact finite-coefficient descent
of Lemma~\textup{R4--78.1}.  It gives
$E=\overline{K(c_1,\ldots,c_m)}\subset L$ and a compact original
$E$-test $\eta$ with $\beta=p^*\eta$.  The just-proved
$\mathsf P(E,0)$ and the fixed scalar integration identity give
\[
 \langle D[u_L-v_L],\beta\rangle
       =\langle D[u_E-v_E],\eta\rangle=0.
\]
The test was arbitrary.  Ordinary localization proves the
assertion on the original allowed space $X$.
\end{proof}

This theorem closes the scalar-extension gate isolated in R4.
It does not by itself prove continuity of the actual difference
$u-v$, or its functoriality under every analytic map over an
arbitrary extension.  Those are separate assertions.  In particular
the arbitrary-background quotient is not invoked anywhere in the
proof and is not claimed as a consequence without the remaining
analytic argument.

% END source: sections/scalar_extension_kernel.tex


% BEGIN canonical R5 source: sections/extension_analytic_kernel.tex
\section{Analytic pullback of the bounded kernel after every complete extension}
\label{sec:r5-extension-analytic-kernel}

The ground field of the original potentials is the fixed finite extension
$K/\mathbf Q_p$, with its original absolute value.  An allowed space over a
complete extension $L/K$ is good, Hausdorff, boundaryless, and quasi-smooth.
In particular, its structural morphism is smooth in the sense of
Berkovich.  We use the original CLD forms and currents over the indicated
field.  The functions are actual locally bounded classical psh functions,
with the all-complete-extension curve condition.

The result of this section is a pullback statement for differences of
specified functions.  No pullback of an arbitrary current is defined.
Algebraic neighborhoods and tubular neighborhoods are used only over the
original finite field $K$, where B, Lemma D.4, and R, Lemma 5.2, supply
them.  The new source over $L$ is handled by genuinely etale coordinates
and a field-independent normal trace.

\begin{lemma}[Ordinary finite etale coordinates]
\label{lem:r5-ordinary-finite-etale-coordinates}
Let $L$ be a complete nontrivially valued field and let $Z$ be a good,
boundaryless, quasi-smooth $L$-analytic space.  For every point $z\in Z$
of local dimension $d$, there are ordinary open neighborhoods $Z_0$ of
$z$ and $B$ in $\mathbf A_L^{d,\mathrm{an}}$ and a finite etale
morphism
\[
                         g:Z_0\longrightarrow B.
\]
The neighborhoods can be made small enough to meet any prescribed
ordinary neighborhood condition at $z$.
\end{lemma}

\noindent\textbf{Status: \texttt{EXTERNAL\_KNOWN}.}
The precise geometry is the definition of smooth morphisms in
Berkovich~\cite[\S3.5]{Ber93}, together with the definition of
quasifinite morphisms in \cite[Definition~3.1.1 and Lemma~3.1.2]{Ber93}
and the fact that etale morphisms are quasifinite in
\cite[3.3.4]{Ber93}.  These conventions and
their ordinary-open meaning are also stated explicitly on printed
page~9 of that paper.

\begin{proof}
Smoothness gives, after an ordinary shrinking at $z$, an etale map to
$\mathbf A_L^{d,\mathrm{an}}$.  In Berkovich's terminology an etale
map is quasifinite.  Quasifiniteness means that at the chosen source
point it restricts to a finite map between ordinary open neighborhoods
of that point and its image.  The restriction remains etale.  Starting
with the prescribed smaller ordinary neighborhood proves the final
assertion.  This argument does not replace etale by quasi-etale: an
analytic-domain immersion with relative boundary would not justify the
assertion.  Boundarylessness in the hypotheses is used precisely to
pass from quasi-smoothness to smoothness.
\end{proof}

\begin{lemma}[Normal trace over an arbitrary complete field]
\label{lem:r5-arbitrary-field-normal-trace}
Let $L$ be complete, let $B$ be an allowed pure $d$-dimensional
$L$-analytic space, and put
\[
 W=B\times_L\mathbb D_s,\qquad
 \pi:W\to B,\qquad i:B\to W,
 \qquad \mathbb D_s=\{|t|<s\}.
\]
For every actual locally bounded classical psh function $a$ on $W$
and every original $\eta\in A_c^{d,d}(B)$, one has
\begin{equation}
 \lim_{N\to\infty}\int_W a\,\beta_N\wedge\pi^*\eta
       =\int_B(a\circ i)\,\eta,
 \label{eq:r5-arbitrary-field-normal-trace}
\end{equation}
where
\[
 \beta_N=D[-N+\psi(\log|t|+N)]
\]
and $\psi$ is any fixed smooth convex function equal to zero for
$x\leq-\epsilon$ and to $x$ for $x\geq\epsilon$.
The function in brackets is extended constantly at $t=0$; for large
$N$, $\beta_N$ is an original smooth positive closed form supported
on the compact normal annulus
$e^{-N-\epsilon}\leq |t|\leq e^{-N+\epsilon}<s$.
All coefficients are evaluated on their actual Borel representatives.

Consequently, for actual locally bounded classical psh functions
$a,b$ on $W$,
\[
 D[a]=D[b]\quad\Longrightarrow\quad
 D[a\circ i]=D[b\circ i]
\]
on every original complementary test on $B$.
\end{lemma}

\noindent\textbf{Status: \texttt{NEW\_PROVED} extension in the
ground field of the normal-trace argument from R3.}
The proof requires the given functions to be classical; it does not
obtain classicality from a distributional inequality or assume a
regularization theorem over $L$.

\begin{proof}
For $x\in B$ and $r>0$, let $\sigma_r(x)$ denote the image in
$W$ of the radius-$r$ Gauss point in the normal disc over
$\mathscr H(x)$.  Let $\mu_\eta$ be the signed Radon measure
associated with the original smooth top form $\eta$.  Original
product calibration gives, for every bounded Borel coefficient $c$
on a neighborhood of the relevant compact support,
\begin{equation}
 \int_W c\,\beta_N\wedge\pi^*\eta
 =\int_B\int_{\mathbf R}
 c(\sigma_{e^t}(x))\psi''(t+N)\,dt\,d\mu_\eta(x).
 \label{eq:r5-arbitrary-field-normal-coarea}
\end{equation}
Here the inner measure is supported below $\log s$ when $N$ is
large.  To specify the calibration, take a finite original moment
presentation for the base top form, retaining the individual
calibrated top-dimensional cells.  Adjoining the unit $t$ on the
normal annulus gives the product of each such cell with the radial
interval.  The relative cell measure is $dt$ with multiplicity one;
the base multiplicities are already included in $\mu_\eta$.
Thus smooth scalar coefficients satisfy
\eqref{eq:r5-arbitrary-field-normal-coarea} by integration on these
product cells.  The same calculation identifies the coefficient
Radon measure itself, so Borel multiplication and ordinary Fubini
give the displayed formula for bounded Borel $c$.  Finite partitions
and signed linearity recover all original $\eta$.
This is the original product calibration; it is not an assertion
about arbitrary points of a set-theoretic product or about products
of arbitrary currents.  The argument is unchanged over $L$.

The Gauss lift $(x,r)\mapsto\sigma_r(x)$ is continuous on every
compact affinoid stage and compact radial interval.  Indeed, the norm
of a polynomial $\sum_j f_jt^j$ there is
$\max_j|f_j(x)|r^j$; uniform analytic approximation proves the
claim for analytic functions.  The integrand in
\eqref{eq:r5-arbitrary-field-normal-coarea} is therefore Borel.

Fix $x\in B$.  Classicality, including complete extension to
$\mathscr H(x)$, makes the restriction of $a$ to the normal
$\mathscr H(x)$-disc a locally bounded subharmonic function.
The maximum principle on a closed disc, whose Shilov boundary is its
Gauss point, gives
\[
 a(i(x))\leq a(\sigma_r(x)).
\]
As $r\downarrow0$ the Gauss points tend to the zero section.  Upper
semicontinuity gives the reverse limiting upper bound, so
\begin{equation}
             a(\sigma_r(x))\longrightarrow a(i(x)).
 \label{eq:r5-normal-actual-radial-limit}
\end{equation}
This concerns a pullback of a function to a curve.  It assumes no
extension or pullback of a current equality.

Cover $\operatorname{supp}\eta$ by finitely many compact affinoid
neighborhoods in $B$ and choose a fixed closed normal disc of radius
smaller than $s$.  The union of their actual analytic products is
compact.  Local boundedness gives a
uniform bound for $a$ there.  For large $N$ all the normal annuli in
\eqref{eq:r5-arbitrary-field-normal-coarea} lie in this compact
product, and their radial measures have mass one and radii tending
uniformly to zero.  Apply \eqref{eq:r5-normal-actual-radial-limit}
and dominated convergence against the fixed finite measure
$|\mu_\eta|$.  This proves
\eqref{eq:r5-arbitrary-field-normal-trace}.

For the last assertion assume $d\geq1$, take an arbitrary original
$\alpha\in A_c^{d-1,d-1}(B)$, and test the given equality against
$\beta_N\wedge\pi^*\alpha$.  Separate closedness of $\beta_N$
gives
\[
 0=\int_W(a-b)\,\beta_N\wedge\pi^*(D\alpha).
\]
Apply the trace limit to $a$ and $b$ separately.  The limit is
$\langle D[a\circ i]-D[b\circ i],\alpha\rangle$.
When $d=0$ there are no curvature tests and the assertion is
immediate.  Iterating the argument gives the same restriction
statement for a zero section of any finite-dimensional normal
polydisc, with arbitrary allowed analytic base over $L$.
\end{proof}

\begin{theorem}[Analytic functoriality over every complete extension]
\label{thm:r5-all-extension-analytic-kernel}
Let $X$ be an allowed $K$-analytic space, let
$u,v\in\mathcal C_b(X)$ satisfy $D[u]=D[v]$ on every original
compact smooth complementary test, let $L/K$ be any complete
extension with the induced absolute value, and let
\[
                       F:Z\longrightarrow X_L
\]
be an arbitrary analytic morphism from an allowed $L$-analytic
space.  Then
\begin{equation}
               D[u_L\circ F]=D[v_L\circ F]
 \label{eq:r5-all-extension-analytic-kernel}
\end{equation}
on every original compact smooth complementary test on $Z$.
Both functions are their actual locally bounded classical
pullbacks.
\end{theorem}

\noindent\textbf{Status: \texttt{NEW\_PROVED}}, relative to
Theorem~\ref{thm:r5-all-complete-scalar-kernel}, the prescribed-row
finite-etale theorem, and the accepted original-$K$ geometric
inputs.  In particular the theorem includes analytic curve maps
over arbitrary complete extensions, not only maps defined over a
finite extension of $K$.

\begin{proof}
The assertion is local at a point $z\in Z$.  A zero-dimensional
source has no curvature tests.  A zero-dimensional target is
locally constant, so its pullbacks also have zero curvature.
We may assume that the relevant target and source dimensions are
positive.

Let $x$ be the image in $X$ of $F(z)$.  B, Lemma D.4, gives an
ordinary algebraic neighborhood of $x$ over the original $K$.
Choose a smooth affine $K$-variety $Y$ containing this neighborhood
and a closed algebraic embedding $Y\hookrightarrow\mathbf A_K^N$.
After shrinking at $x$, R, Lemma 5.2, gives an ordinary open
$V\subset X$ containing $x$, an ordinary open
$W\subset\mathbf A_K^{N,\mathrm{an}}$, and an analytic isomorphism
\[
        \Phi:V\times_K\mathbb D_s^c\xrightarrow{\sim}W
\]
whose zero section is the given embedding.  Let
$R_W=\operatorname{pr}_V\circ\Phi^{-1}$ and put
\[
            \widetilde u=u\circ R_W,
            \widetilde v=v\circ R_W.
\]
These are actual classical functions.  They have prescribed smooth
positive rows by pulling back the original-$K$ rows.  The R4
polydisc kernel theorem and invariance of original forms and
calibrated integrals under an analytic isomorphism show that
$D[\widetilde u]=D[\widetilde v]$ on $W$.
For every integer $d\geq0$, the same R4 product theorem gives the
kernel equality for the functions depending only on $W$ on
$W\times_K\mathbf A_K^{d,\mathrm{an}}$.
Theorem~\ref{thm:r5-all-complete-scalar-kernel} then gives that
equality on
\begin{equation}
                 W_L\times_L\mathbf A_L^{d,\mathrm{an}}.
 \label{eq:r5-ambient-product-after-extension}
\end{equation}
Only the originally defined functions were extended in this step.

Let $d=\dim_z Z$.  Shrink at $z$ until $F$ takes values in $V_L$.
By Lemma~\ref{lem:r5-ordinary-finite-etale-coordinates}, a further
ordinary neighborhood $Z_0$ admits a finite etale map
$g:Z_0\to B$, with $B$ ordinary open in
$\mathbf A_L^{d,\mathrm{an}}$.  The morphism
\[
 \operatorname{id}_{W_L}\times g:
             W_L\times_L Z_0\longrightarrow W_L\times_L B
\]
is finite etale.  Apply
Theorem~\ref{thm:r5-prescribed-finite-etale-kernel} to the
restriction of \eqref{eq:r5-ambient-product-after-extension}.
The needed rows are the scalar pullbacks of the original rows of
$\widetilde u,\widetilde v$, followed by the projection.
This proves the kernel equality for $\widetilde u_L(x)$ and
$\widetilde v_L(x)$ on $Z_0\times_L W_L$.
No algebraization or regularization theorem over $L$ is used.

Regard $F|_{Z_0}$ as an analytic map to $\mathbf A_L^N$ and use
the analytic translation
\[
 A:Z_0\times_L\mathbf A_L^{N,\mathrm{an}}\longrightarrow
          Z_0\times_L\mathbf A_L^{N,\mathrm{an}},\qquad
 A(c,t)=(c,F(c)+t).
\]
Its inverse is $(c,x)\mapsto(c,x-F(c))$.  The inverse image
$O=A^{-1}(Z_0\times_L W_L)$ is an ordinary open containing the
zero section.  On it the two actual classical functions
\[
 a(c,t)=\widetilde u_L(F(c)+t),\qquad
 b(c,t)=\widetilde v_L(F(c)+t)
\]
have equal original curvature currents, because $A$ is an analytic
isomorphism.

The compact-tube proof of Lemma~\textup{R4--75.2}
works over $L$ verbatim.  Indeed, on a compact affinoid base
neighborhood the nested compact analytic tubes of radii tending
to zero intersect in the zero section; compactness forces a
sufficiently small entire tube into $O$.  Hence there is an
ordinary neighborhood $Z_1$ of $z$ and $r>0$ such that
$Z_1\times_L\mathbb D_r^N\subset O$.  This is an assertion
about actual analytic products, not about Cartesian products of
their points.  Restrict the equality to this tube and apply
Lemma~\ref{lem:r5-arbitrary-field-normal-trace} successively in
the $N$ normal coordinates.  It gives
\[
                       D[a|_{t=0}]=D[b|_{t=0}]
                              \quad\hbox{on }Z_1.
\]
The zero-section values are precisely $u_L\circ F$ and
$v_L\circ F$.  Locality proves
\eqref{eq:r5-all-extension-analytic-kernel} on all of $Z$.
\end{proof}

\begin{corollary}[All complete-extension curve kernels]
\label{cor:r5-all-extension-curve-kernel}
Under the hypotheses on $X,u,v$ in the theorem, every complete
$L/K$ and every analytic map from a smooth boundaryless
$L$-analytic curve $C$ to $X_L$ satisfy
\[
                 D[(u_L-v_L)\circ F]=0
\]
on every original compact smooth scalar test on $C$.
\end{corollary}

\noindent\textbf{Status: \texttt{DIRECT\_COROLLARY}.}
The assertion includes the actual representatives of both
subharmonic pullbacks.  It does not by itself assert that
$u-v$ or $v-u$ is upper semicontinuous on $X$.

\paragraph{Scope of the remaining quotient question.}
This section closes the analytic curve-test gate after all complete
extensions.  It does not infer continuity on a higher-dimensional
space merely from the existence of those curve tests.  In
particular, the implication
\[
       D[u]=D[v]\quad\Longrightarrow\quad D((u-v)S)=0
\]
for every allowed positive closed background $S$ still requires
the separate actual-regularity or arbitrary-background argument.
The all-extension functoriality established here does not define
the pullback of such an $S$.

% END source: sections/extension_analytic_kernel.tex


% BEGIN canonical R5 source: sections/curve_kernel_actual.tex
\section{The actual harmonic representative on complete-extension curves}
\label{sec:r5-curve-kernel-actual}

The distinction between equality of curve distributions and equality
of actual functions matters here.  Thuillier's degree-zero
distributions are determined by the values at points of type II or
III.  An arbitrary Borel modification at a rigid point would not be
detected by those values.  The following argument recovers the missing
values from the separate subharmonicity of the two given functions.
It does not assume that their difference is semicontinuous.

\begin{lemma}[The algebraic scope of the curve input]
\label{lem:r5-curve-local-algebraization}
Over any complete nontrivially valued field $F$, a good Hausdorff
boundaryless quasi-smooth analytic curve is, locally for the ordinary
topology, isomorphic to an ordinary open in the analytification of a
smooth affine algebraic $F$-curve.  No extension of $F$ is necessary.
Consequently the explicitly complete-field statement of B, Lemma A.2,
applies locally to actual locally bounded subharmonic functions on such
an analytic curve, when $\operatorname{char}F=0$.
\end{lemma}

\noindent\textbf{Status: \texttt{EXTERNAL\_KNOWN} geometry and
\texttt{DIRECT\_COROLLARY} for the use of the internal input.}

\begin{proof}
Apply Ducros~\cite[Lemma~(0.21)]{Duc12} with the full radius group
$\Gamma=\mathbf R_{>0}$.  The lemma is stated for an arbitrary
complete non-Archimedean field.  At the chosen point it gives an
affinoid neighborhood $V$ identified with an affinoid domain in a
smooth affine algebraic $F$-curve $Y$.
Shrink the initial neighborhood to be separated.  The relative
interior composition formula, with its locally separated intermediate
spaces, identifies
\[
 \operatorname{Int}(V/Z)=\operatorname{Int}(V/F)
                       =\operatorname{Int}(V/Y^{\mathrm{an}}),
\]
because $Z$ and $Y^{\mathrm{an}}$ are boundaryless.
For an analytic-domain embedding the relative interior is the
ordinary topological interior.  The common interior contains the
chosen point and is ordinary open in both ambient spaces.

This uses the external algebraization lemma in its stated field
generality, not an unstated extension of B, Lemma D.4.  The invoked
part of B, Lemma A.2, already assumes an arbitrary complete
nontrivially valued characteristic-zero field and an ordinary open
in a smooth algebraic curve.  Its canonical-calibration identity is
\begin{equation}
 \int_C\alpha\,d(\Delta_C a)=\int_C a\,D\alpha,
       \qquad \alpha\in A_c^{0,0}(C),
 \label{eq:r5-curve-laplacian-compatibility}
\end{equation}
where $\Delta_C a$ is the positive Radon Laplacian of the given
subharmonic function.  Thus its hypotheses hold on the obtained
ordinary algebraic neighborhood.
\end{proof}

\begin{lemma}[Recovering the actual curve function]
\label{lem:r5-actual-curve-harmonic-representative}
Let $F$ be complete, algebraically closed, nontrivially valued, and
of characteristic zero.  Let $C$ be a good Hausdorff smooth
boundaryless $F$-analytic curve.  Let $a,b$ be actual locally bounded
subharmonic functions on $C$.  If
\[
                         D[a]=D[b]
\]
on every original compact smooth scalar test, then $a-b$ is an
actual continuous harmonic function on $C$ in the curve-potential
sense.  In particular both signs of $a-b$ are subharmonic.
\end{lemma}

\noindent\textbf{Status: \texttt{DIRECT\_COROLLARY}} of the
complete-field curve Laplacian input and the classical curve kernel
theorem, with the actual-value passage proved explicitly below.
The external distributional kernel is Thuillier's
Lemma~3.3.12; Wanner~\cite[Remark~7.2, p.~22, and Remark~7.12,
p.~27]{Wan22} explicitly applies it before identifying the endpoint
values.  We also give its finite-graph verification in the present
bounded setting.

\begin{proof}
Work on an ordinary algebraic neighborhood supplied by
Lemma~\ref{lem:r5-curve-local-algebraization}.  B, Lemma A.2, and
\eqref{eq:r5-curve-laplacian-compatibility} show that the positive
Radon measures $\Delta_C a$ and $\Delta_C b$ agree against every
original compact smooth scalar function.  Original smooth functions
are uniformly dense on each compact stage, with a cutoff on a fixed
larger stage.  The measures have finite mass on that larger stage,
so they agree against all compact continuous functions and hence as
Radon measures.  Therefore the degree-zero curve distribution
defined by $a-b$ has Laplacian zero.  This distribution records the
actual values of $a-b$ on the subset $I(C)$ of type-II and type-III
points.

Here is a local graph argument for its harmonic representative.
Take a connected strictly affinoid domain $Y$ compactly contained
in the chosen curve open, with nonempty boundary.  On a finite
semistable skeleton, solve the finite linear Dirichlet problem with
boundary values $(a-b)|_{\partial Y}$, and extend the solution by
the skeleton retraction.  The resulting function $h$ is continuous
on $Y$ and harmonic on its interior.  Existence and uniqueness on
the finite graph follow by minimizing the strictly positive
Dirichlet quadratic form with fixed boundary values; equivalently,
the weighted graph Laplacian on the nonboundary vertices is
invertible.

For $x\in I(\operatorname{Int}Y)$ use the curve Green function
$g_{Y,x}$, extended by zero outside $Y$.  It is lisse and compactly
supported; its Laplacian is supported on $\partial Y\cup\{x\}$
and has coefficient $-1$ at $x$.  These are the Green functions of
Thuillier~\cite[Proposition~3.3.7 and Corollary~3.3.9]{Thu05},
also stated in Wanner~\cite[Proposition~2.12]{Wan19}.
The equation $\Delta(a-b)=0$, evaluated on $g_{Y,x}$, and the
finite-graph Green identity for the harmonic extension $h$, give
\[
              \int (a-b-h)\,d(\Delta g_{Y,x})=0.
\]
The boundary values of $a-b-h$ vanish, so $a(x)-b(x)=h(x)$.
Such interiors cover $C$ locally, and their harmonic functions agree
on overlaps by continuity and density of type-II and type-III
points.  They therefore give a continuous harmonic $h$ agreeing
with $a-b$ on $I(C)$.

It remains to identify actual values at points of types I and IV.
Each such point $x$ has an ordinary open-ball neighborhood; see
\cite[Theorem~2.11]{Wan22}.  In that ball choose nested closed
discs $D_j$ containing $x$, with Gauss boundary points
$\xi_j\in I(C)$ tending to $x$.  For type I their radii tend to
zero about the corresponding rational center.  For type IV the
discs come from its nested-disc description, with empty
intersection of classical points and positive limiting radius.
They may be chosen sufficiently far along the system to lie in
the given open-ball neighborhood.  Their radii may also be chosen
in $|F^\times|$, which is dense in $\mathbf R_{>0}$ after taking
logarithms, so the disc maximum principle can be used on strictly
affinoid domains.

Apply the subharmonic maximum principle separately to $a$ and $b$:
\[
                  a(x)\leq a(\xi_j),\qquad
                  b(x)\leq b(\xi_j).
\]
Upper semicontinuity of each given function and $\xi_j\to x$
give the reverse limiting upper bounds.  Thus
\[
            a(\xi_j)\to a(x),\qquad b(\xi_j)\to b(x).
\]
Since $a(\xi_j)-b(\xi_j)=h(\xi_j)$ and $h$ is continuous,
we conclude $a(x)-b(x)=h(x)$.  This proves equality at every point,
including the endpoint values invisible to the degree-zero
distribution by itself.
\end{proof}

\begin{theorem}[Actual harmonicity on every complete-extension curve]
\label{thm:r5-all-extension-curve-harmonicity}
Let $u,v\in\mathcal C_b(X)$ on an allowed original $K$-space
satisfy $D[u]=D[v]$.  For every complete extension $L/K$ and every
analytic map $F:C\to X_L$ from a good Hausdorff smooth boundaryless
$L$-analytic curve, the actual function
\[
                         h_C=(u_L-v_L)\circ F
\]
is continuous and harmonic: both $h_C$ and $-h_C$ are subharmonic
on $C$.
\end{theorem}

\noindent\textbf{Status: \texttt{DIRECT\_COROLLARY}} of
Theorem~\ref{thm:r5-all-extension-analytic-kernel}, the preceding
curve lemma, and scalar descent for subharmonic functions.

\begin{proof}
Let $\Omega$ be a completed algebraic closure of $L$.
Theorem~\ref{thm:r5-all-extension-analytic-kernel}, applied to
$F_\Omega:C_\Omega\to X_\Omega$, gives the original-current
equality for the two actual subharmonic pullbacks on $C_\Omega$.
Lemma~\ref{lem:r5-actual-curve-harmonic-representative}, applied
locally, says that $h_C\circ p$ is continuous and harmonic on
$C_\Omega$, where $p:C_\Omega\to C$ is scalar projection.

Continuity descends without a countability assumption.  On every
compact affinoid neighborhood $A\subset C$, the map
$A_\Omega\to A$ is a continuous surjection from a compact space
onto a Hausdorff space, hence a quotient map.  The continuous
function $h_C\circ p$ is constant on its fibers, so $h_C|_A$ is
continuous.  Affinoid neighborhoods give continuity locally on
$C$.  Finally the scalar-extension equivalence for subharmonic
functions, Gubler--Rabinoff~\cite[Proposition~3.5(7)]{GR25},
applied separately to the now continuous functions $h_C$ and
$-h_C$, gives their subharmonicity on $C$.
\end{proof}

\paragraph{A boundary of this conclusion.}
The theorem establishes actual continuity and harmonicity on each
of the permitted source curves, after all complete extensions.
It does not replace a proof of ordinary continuity of $u-v$ on
the higher-dimensional target.  In the definition of classical
pluriharmonicity that additional topological condition still has
to be established; no such inference is made from density of
curve images alone.

% END source: sections/curve_kernel_actual.tex


% BEGIN canonical R5 source: sections/shilov_maximum_without_usc.tex
\section{A Shilov maximum principle without ambient upper semicontinuity}
\label{sec:r5-shilov-without-usc}

Throughout this section the absolute values are fixed, and scalar
extensions carry the extending absolute value.  Harmonicity on a
smooth analytic curve means harmonicity of the actual function; in
particular that function is continuous on the curve.  We do not
assume that a function which has harmonic curve pullbacks is upper
semicontinuous on its higher-dimensional ambient space.

\begin{lemma}[The Gauss fiber of a finite flat affinoid map]
\label{lem:r5-finite-flat-gauss-shilov}
Let $F$ be a complete non-Archimedean field, let
$A=F\langle T_1,\ldots,T_n\rangle$, and let $B$ be a nonzero
finite flat $A$-algebra.  Write $V=\mathcal M(B)$,
$\mathbb B_F^n=\mathcal M(A)$, and let
$g:V\to\mathbb B_F^n$ be the resulting map.  If $\zeta$ is the
unit-polydisc Gauss point, then
\[
                         \Gamma(V)=g^{-1}(\zeta).
\]
In particular this is a finite nonempty set.  Nilpotents in $B$
do not change the assertion.
\end{lemma}

\noindent\textbf{Status: \texttt{DIRECT\_COROLLARY}} of the
standard finite-algebra norm and Shilov-boundary facts.  We give the
argument to fix the exact boundary needed below.

\begin{proof}
The finite locally free $A$-module $B$ has a positive constant rank
$d$.  For $b\in B$, let
\[
 \det(T-m_b)=T^d+c_1T^{d-1}+\cdots+c_d\in A[T]
\]
be the characteristic polynomial of multiplication by $b$.  For
each $y\in\mathbb B_F^n$, its roots after extension to an
algebraic closure of $\mathscr H(y)$ are the values of $b$ on the
finite geometric fiber, with their algebraic multiplicities.
The maximum of their absolute values is
\[
 \max_{x\in g^{-1}(y)}|b(x)|
                  =\max_{1\leq i\leq d}|c_i(y)|^{1/i}.
\]
Since $|c_i(y)|\leq|c_i(\zeta)|$, this maximum is at most the
corresponding maximum over $g^{-1}(\zeta)$.  Thus that finite
fiber is a boundary for $B$, and the Shilov boundary is contained
in it.

Put $H=\mathscr H(\zeta)$ and
$D=B\otimes_A H$.  This finite-dimensional $H$-algebra is a
product of Artinian local algebras, one for each point of
$g^{-1}(\zeta)$.  Fix a point $z$ in that fiber and let $e_z$
be the idempotent of its factor.  The field $\operatorname{Frac}(A)$
is dense in $H$, and consequently
$B\otimes_A\operatorname{Frac}(A)$ is dense in $D$.
Approximate $e_z$ by an element $b/a$, with $b\in B$ and
$0\ne a\in A$, closely enough that
\[
 |(b/a)(z)|=1,
 \qquad |(b/a)(z')|<1\quad(z'\ne z, z'\in g^{-1}(\zeta)).
\]
The denominator has the same nonzero absolute value $|a(\zeta)|$
at every point of the fiber.  Hence $|b|$ has a strict unique
maximum at $z$ within this finite boundary.  Since the Shilov
boundary is contained in the fiber and is itself a boundary, it
must contain $z$.  This proves the reverse inclusion.  The proof
uses the full finite algebra and therefore also applies in the
presence of nilpotents or valued-field defect.
\end{proof}

\begin{theorem}[Curvewise harmonic functions satisfy the affinoid Shilov bound]
\label{thm:r5-curvewise-shilov-without-usc}
Let $F$ be a complete nontrivially valued field of characteristic zero, let $Y$ be a
smooth affine $F$-variety of pure dimension $n\geq1$, and let
$U\subset Y^{\rm an}$ be ordinary open.  Let $f:U\to\mathbf R$
be an actual real-valued function.  Assume the following
curve property:

for every complete extension $E/F$, every smooth algebraic
$E$-curve $C$, and every algebraic morphism $h:C\to Y_E$, the
actual pullback of $f$ is harmonic on each ordinary open subset
of $C^{\rm an}$ mapped into $U_E$.

Suppose $V\subset U$ is a nonempty strict affinoid domain and
$g_1,\ldots,g_n$ are regular functions on $Y$ such that
\[
 g=(g_1,\ldots,g_n):V\longrightarrow\mathbb B_F^n
\]
is finite flat.  Then
\begin{equation}
                  |f(x)|\leq\max_{z\in\Gamma(V)}|f(z)|
                  \qquad(x\in V).
 \label{eq:r5-shilov-bound-without-usc}
\end{equation}
No ordinary upper semicontinuity, local boundedness, or Borel
measurability of $f$ on $U$ is assumed.
\end{theorem}

\noindent\textbf{Status: \texttt{NEW\_PROVED}.}
Only the one-dimensional maximum principle is used; higher-dimensional
ambient upper semicontinuity is not inserted into its hypotheses.
The finite Shilov bound proves boundedness on the displayed domain.

\begin{proof}
We first prove a one-coordinate inequality.  Fix $x\in V$, write
$b=g(x)$, and fix an index $i$.  Let $y$ be the image of $b$ in
the polydisc of the other $n-1$ coordinates, and put
$E=\mathscr H(y)$.  The fiber
\[
 W=V\times_{\mathbb B_F^{n-1}}\mathcal M(E)
\]
is an affinoid $E$-space, finite flat over the unit disc with
coordinate $T_i$.  It is purely one-dimensional.  The point $x$
defines a point of $W$, using the canonical embedding
$E\hookrightarrow\mathscr H(x)$.

Here the full analytic fiber is used.  It must not be replaced by
a selected independent section.  Its points, including the chosen
lift of $x$, retain their actual pullback values of $f$.

Let $a_j\in E$, $j\ne i$, be the coordinate values of $y$.
Inside $Y_E$, form the algebraic slice
\[
                 Z=\{g_j=a_j\text{ for all }j\ne i\}.
\]
The space $W$ is an analytic domain of $Z^{\rm an}$ lying in
$U_E$.  Local analytic dimension is unchanged by restriction to
an analytic domain, including at its ordinary boundary; see
\cite[\S1.4.9]{DucrosFamilies}.  Thus $Z^{\rm an}$ has local
dimension one at every point of $W$.  Dimension and irreducible
components commute with algebraic analytification
\cite[Propositions 2.7.7(2) and 2.7.16]{DucrosFamilies}.
Remove the finitely many algebraic components of $Z$ of dimension
different from one, all disjoint from $W$, and give the resulting
algebraic curve its reduced structure.  Its normalization $C$ is
finite and, since $E$ has characteristic zero, is a smooth
algebraic $E$-curve.  Its algebraic map to $Y_E$ will be denoted
by $h$.

The inverse image $W^\nu$ of $W$ is a compact affinoid analytic
domain in $C^{\rm an}$.  Normalization commutes with restriction
to analytic domains; equivalently one may use finite algebraic
normalization followed by analytification, as in
\cite[Lemma 2.7.15]{DucrosFamilies}.  The map
\[
                     W^\nu\longrightarrow\mathbb B_E^1
\]
is finite.  Each of its irreducible components dominates the
disc: $W$ is finite flat and purely one-dimensional, and finite
normalization preserves that dimension.  Its coordinate algebra
is therefore torsion-free over the regular one-dimensional Tate
algebra $E\langle T_i\rangle$, hence is finite flat.  By
Lemma~\ref{lem:r5-finite-flat-gauss-shilov},
$\Gamma(W^\nu)$ is its fiber over the unit-disc Gauss point.

By hypothesis $f_E\circ h$ is harmonic on the ordinary open
subset $h^{-1}(U_E)$ of $C^{\rm an}$, which contains all of
$W^\nu$.  It is consequently continuous on a genuine
neighborhood of this compact domain, including its boundary.
The curve maximum principle, applied to both signs, gives
\[
 |f(x)|\leq\max_{w\in\Gamma(W^\nu)}|f_E(h(w))|.
\]
This use of continuity is solely on a smooth curve.  The
higher-dimensional function $f$ has not been assumed continuous
or upper semicontinuous.

Define $G_i(b)$ as follows: retain the other coordinate point $y$,
take the unit-radius Gauss point in the $T_i$-disc over
$\mathscr H(y)$, and map it into $\mathbb B_F^n$.  Every point
$h(w)$ in the last maximum projects to a point of $V$ lying
over $G_i(b)$.  Since $g$ is finite, the function
\[
 M(b)=\max\{|f(z)|:z\in V,\ g(z)=b\}
\]
is finite and defined at every $b\in\mathbb B_F^n$: each fiber
is nonempty and finite, and $f$ is real-valued. No common bound
over different fibers is assumed. The one-coordinate
argument proves
\begin{equation}
                              M(b)\leq M(G_i(b)).
 \label{eq:r5-coordinate-gauss-maximum}
\end{equation}

For completeness, these coordinate Gauss operations commute.
For a polynomial $P=\sum_k P_k(T_1,\ldots,\widehat T_i,\ldots,T_n)T_i^k$
one has
\[
                         |P(G_i(b))|=\max_k|P_k(y)|.
\]
Applying this formula successively to two different coordinates
gives the maximum of the absolute values of the corresponding
double coefficient polynomials, independently of their order.
After all $n$ coordinates, the value of every polynomial is its
Gauss norm over $F$.  By uniform density of polynomials in the
Tate algebra, the resulting point is precisely the full polydisc
Gauss point $\zeta$.

Iterating \eqref{eq:r5-coordinate-gauss-maximum} therefore gives
\[
                   |f(x)|\leq M(g(x))\leq M(\zeta).
\]
Lemma~\ref{lem:r5-finite-flat-gauss-shilov} identifies the last
maximum with $\max_{\Gamma(V)}|f|$, proving
\eqref{eq:r5-shilov-bound-without-usc}.
\end{proof}

\paragraph{Scope and use of the ambient neighborhood.}
The regularity of the $g_i$ on $Y$ ensures that every sliced
affinoid curve is contained in the analytification of an algebraic
curve extending across its compact boundary.  Merely knowing a
harmonic function on the interior of an abstract affinoid curve
would not justify its boundary maximum with arbitrary assigned
boundary values.  The theorem states the needed ambient condition
explicitly.  It requires no global psh exhaustion and no strictly
positive ambient form.  Producing the indicated affinoids and
ambient regular normalization coordinates is a separate geometric
input; it is not inferred from the maximum principle itself.

% END source: sections/shilov_maximum_without_usc.tex


% BEGIN canonical R5 source: sections/vertical_shilov_interpolation.tex
\section{Vertical divisors and continuous Shilov interpolation}
\label{sec:r5-vertical-shilov-interpolation}

This elementary geometric input replaces any need to identify an
entire skeleton over a non-algebraically closed field.  The base field
$F$ is complete and discretely valued, with valuation ring $F^\circ$
and uniformizer $\varpi$.  Its given absolute value is retained.

\begin{lemma}[Vertical components are the exact Shilov boundary]
\label{lem:r5-vertical-shilov-boundary}
Let $R$ be an integral, regular, flat, finite-type $F^\circ$-algebra.
Suppose its nonempty special fiber is reduced and connected, with
irreducible components $D_0,\ldots,D_r$.  Write
\[
 R^\wedge=\varprojlim_m R/\varpi^mR,\qquad
 A=R^\wedge[1/\varpi],\qquad W=\mathcal M(A).
\]
For each $i$, the height-one prime of $R^\wedge$ corresponding to
the generic point of $D_i$ defines a divisorial point $\xi_i\in W$.
Then
\begin{equation}
 \Gamma(W)=\{\xi_0,\ldots,\xi_r\},\qquad
 \sup_{x\in W}|a(x)|=\max_i|a(\xi_i)|\quad(a\in A).
 \label{eq:r5-vertical-shilov-norm}
\end{equation}
If, in addition, $D_i$ is the principal effective Cartier divisor
$(x_i)$ on $\operatorname{Spec}R$ for each $i$, and
$x_0\cdots x_r=\varpi$, then
\begin{equation}
 |x_i(\xi_j)|=
 \begin{cases}|\varpi|,&i=j,\\1,&i\ne j.\end{cases}
 \label{eq:r5-vertical-log-coordinate-values}
\end{equation}
\end{lemma}

\noindent\textbf{Status: \texttt{DIRECT\_COROLLARY}} of regular
completion and the height-one intersection theorem for normal rings.
The proof is included to specify the original-field points and norms.

\begin{proof}
First, $R^\wedge$ is a Noetherian regular ring.  Here is a reduction
to local completion, avoiding an unstated formal-model normality
assumption.  Every maximal ideal $\mathfrak n$ of $R^\wedge$ contains
$\varpi$: for every $b\in R^\wedge$, the geometric series makes
$1-\varpi b$ invertible.  Put $\mathfrak m=\mathfrak n\cap R$.
Since $R^\wedge/\varpi R^\wedge=R/\varpi R$, this is a maximal
ideal and $\mathfrak n=\mathfrak mR^\wedge$.  The local map
\[
 R_{\mathfrak m}\longrightarrow (R^\wedge)_{\mathfrak n}
\]
is flat, has the same residue field, and extends the maximal ideal
to the maximal ideal.  It therefore induces an isomorphism of
maximal-ideal completions \cite[Tag 0AGX]{Stacks}.  Regularity is
equivalent to regularity of that completion
\cite[Tag 07NY]{Stacks}.  Thus every maximal localization of
$R^\wedge$ is regular, and so is $R^\wedge$.

The ring $R^\wedge$ is normal.  It is also connected: any idempotent
reduces to $0$ or $1$ modulo $\varpi$ because $R/\varpi R$ is
connected, and that reduction determines the idempotent in the
complete ring.  A connected Noetherian normal ring is a domain.
Flatness of completion preserves the fact that $\varpi$ is a
nonzero divisor.  The height-one primes containing $\varpi$ are
exactly the generic points of the components of
$R^\wedge/\varpi R^\wedge=R/\varpi R$.  At each such prime the
local ring is a discrete valuation ring.  Reducedness of the special
fiber says that its order $v_i$ satisfies $v_i(\varpi)=1$.

The formula $|a|_i=|\varpi|^{v_i(a)}$ on
$\operatorname{Frac}(R^\wedge)$ restricts to a bounded
multiplicative norm on $A$ extending the prescribed absolute value
of $F$.  Boundedness follows because it is at most one on
$R^\wedge$, which is the ring of definition of the affinoid algebra.
This gives the point $\xi_i$.

For $a\ne0$ in $A$, set $m=\min_i v_i(a)\in\mathbf Z$ and
$b=\varpi^{-m}a$.  At every height-one prime above $\varpi$, the
order of $b$ is nonnegative.  At every other height-one prime it is
also nonnegative, since $b\in R^\wedge[1/\varpi]$ and $\varpi$
is a unit there.  The height-one intersection theorem for the
Noetherian normal domain $R^\wedge$
\cite[Tag 0AVB]{Stacks}, applied to the module $R^\wedge$ itself,
gives $b\in R^\wedge$.  Hence $|b(x)|\leq1$ for every $x\in W$,
and
\[
 |a(x)|\leq|\varpi|^m=\max_i|a(\xi_i)|.
\]
The case $a=0$ is immediate.  Thus the finite set of $\xi_i$ is a
boundary and contains $\Gamma(W)$.

To see that no $\xi_i$ can be omitted, choose
$b_i\in R^\wedge$ belonging to all the generic primes
$\mathfrak p_j$ for $j\ne i$ and not belonging to
$\mathfrak p_i$.  Such a choice exists: for each $j\ne i$ choose
an element of $\mathfrak p_j\setminus\mathfrak p_i$ and multiply
them; use $b_i=1$ when there is only one component.  Then
$|b_i(\xi_i)|=1$ and $|b_i(\xi_j)|<1$ for $j\ne i$.
Since the Shilov boundary is contained in the displayed finite
boundary, it must contain $\xi_i$.  This proves
\eqref{eq:r5-vertical-shilov-norm}.

Under the additional principal-divisor assumptions, the order of
$x_i$ along $D_j$ is $\delta_{ij}$.  This proves
\eqref{eq:r5-vertical-log-coordinate-values}.  In that case the
particularly simple peaking function $\varpi/x_i$ takes absolute
value $1$ at $\xi_i$ and $|\varpi|$ at the other displayed points.
\end{proof}

\begin{lemma}[Small algebraic charts with logarithmic interpolation]
\label{lem:r5-local-vertical-interpolation}
Let $\mathcal N$ be a strictly semistable finite-type
$F^\circ$-scheme in the sense of de Jong
\cite[\S2.16]{dJ96}: it is regular and flat, and its special
fiber is a reduced strict normal crossings divisor.  Let
$s\in\mathcal N_s$.  There is an
algebraic affine neighborhood $\mathcal U=\operatorname{Spec}R$
of $s$ with the following properties.  Its nonempty special fiber
is reduced and connected, with irreducible components
$D_0,\ldots,D_r$ all containing $s$.  There are regular functions
$x_0,\ldots,x_r$ defining these components as effective Cartier
divisors and satisfying $x_0\cdots x_r=\varpi$.  Its formal
generic fiber $W$ is a strict affinoid domain in the smooth
algebraic generic fiber $\mathcal U_F^{\rm an}$, and its Shilov
boundary is the set $\{\xi_0,\ldots,\xi_r\}$ in
Lemma~\ref{lem:r5-vertical-shilov-boundary}.

For any prescribed real numbers $c_0,\ldots,c_r$, the function
\begin{equation}
 h=-\frac{1}{-\log|\varpi|}
       \sum_{i=0}^r c_i\log|x_i|
 \label{eq:r5-vertical-log-interpolant}
\end{equation}
is continuous on the whole algebraic generic fiber, is harmonic
after every pullback to a smooth analytic curve over every complete
extension, and satisfies $h(\xi_i)=c_i$.  Every topological
affinoid generator of $W$ can be chosen as the restriction of a
regular function on $\mathcal U_F$.
\end{lemma}

\noindent\textbf{Status: \texttt{DIRECT\_COROLLARY}} of the
strictly semistable divisor description and the preceding lemma.
No algebraic closure or splitting of all strata is needed.

\begin{proof}
Strict semistability implies regularity of $\mathcal N$, reducedness
of its special fiber, and that each irreducible component of the
special fiber is an effective Cartier divisor; these are the
properties in \cite[\S2.16]{dJ96}.  Remove all special-fiber
components not containing $s$.  Shrink to an algebraic affine
neighborhood of $s$ on which the remaining finitely many Cartier
divisors are principal, say with equations $x_i$.  Their
restrictions remain irreducible: each is a nonempty Zariski-open
subset of its original irreducible component.  They all meet at
$s$, so their union is connected.  The ambient regular scheme can
be restricted to its component containing $s$, giving the integral
ring $R$.

The effective Cartier divisor of $\varpi$ is exactly
$D_0+\cdots+D_r$, each with coefficient one.  Therefore
$x_0\cdots x_r/\varpi$ is a regular unit after a further shrinking
around $s$.  Absorb its inverse into $x_0$ to obtain the stated
product equality.  Each $x_i$ is a unit on $\mathcal U_F$.

Apply Lemma~\ref{lem:r5-vertical-shilov-boundary}.  Its coordinate
values give $h(\xi_i)=c_i$.  The function $h$ is a real linear
combination of logarithms of analytic units on a genuine
boundaryless algebraic neighborhood of $W$.  Such logarithms are
continuous and have harmonic curve pullbacks, also after arbitrary
complete scalar extension.  Finally, finite-type algebraic
generators of $R$ give topological generators after completion and
inverting $\varpi$, and are regular on $\mathcal U_F$.
\end{proof}

\paragraph{Use in the continuity argument.}
For an actual curvewise harmonic function $f$, choose
$c_i=f(\xi_i)$.  If the separately proved Shilov principle applies
on $W$ with its ambient regular Noether coordinates, then it applies
to $f-h$.  The latter vanishes on every Shilov point, and hence
$f=h$ everywhere on $W$.  This deduction uses no ambient
semicontinuity of $f$ and no pre-existing pluriharmonicity of $f$.

% END source: sections/vertical_shilov_interpolation.tex


% BEGIN canonical R5 source: sections/algebraic_continuity_geometry.tex
\section{Algebraic neighborhoods for the continuity argument}
\label{sec:r5-algebraic-continuity-geometry}

This section isolates the geometry needed for a Shilov maximum
argument. The original field is a fixed finite extension
$K/\mathbf Q_p$, with its given absolute value. No function in this
section is assumed upper semicontinuous merely because its
restrictions to curves are harmonic.

\begin{lemma}[Polynomial Noether coordinates]
\label{lem:r5-polynomial-noether-coordinates}
Let $F$ be a complete nontrivially valued field of characteristic
zero. Let $Y$ be a smooth affine $F$-variety of pure dimension $n$,
and let $V\subset Y^{\mathrm{an}}$ be a strictly affinoid domain.
Assume that the affinoid algebra $A=\mathscr O(V)$ has a finite
Tate presentation $A=F\langle T_1,\ldots,T_m\rangle/I$ whose
generators $a_1,\ldots,a_m$ are restrictions of regular
functions on $Y$. Then there are
polynomials $P_1,\ldots,P_n\in F[T_1,\ldots,T_m]$ for which
\[
 g_i=P_i(a_1,\ldots,a_m),\qquad
 F\langle Z_1,\ldots,Z_n\rangle\longrightarrow A,\quad Z_i\mapsto g_i,
\]
is finite and faithfully flat. Thus $g=(g_1,\ldots,g_n)$ gives a
finite flat surjection $V\to\mathbb B_F^n$, and each $g_i$
extends regularly to the boundaryless algebraic ambient space.
\end{lemma}

\noindent\textbf{Status: \texttt{DIRECT\_COROLLARY}} of polynomial
Noether normalization and miracle flatness. The extension of the
coordinates is part of the conclusion; it is not an assumed
extension theorem for arbitrary affinoid functions.

\begin{proof}
Write $A=F\langle T_1,\ldots,T_m\rangle/I$ using the given tuple.
The proof of affinoid Noether normalization uses triangular
polynomial automorphisms
\[
 T_i\longmapsto T_i+T_m^{N_i}\quad(i<m),\qquad T_m\longmapsto T_m,
\]
their polynomial inverses, and coordinate inclusions. One chooses
a nonzero element of the current ideal, makes it distinguished in
the last variable, and uses Weierstrass division to obtain a finite
map from the algebra in one fewer variables. Induction on $m$
gives an injective finite map from a Tate algebra of dimension
$\dim A=n$. Since every change of variables is polynomial, the
images of its variables are polynomials in the original tuple.
This is the construction in Bosch's April 2008 manuscript,
\S1.2, Lemma 7 and Corollary 10, printed pages 11 and 14
\cite{Bosch08}. It does not replace the given tuple by arbitrary
analytic functions on $V$.

The algebra $A$ is regular and of pure dimension $n$, since $V$
is an analytic domain in the smooth $Y^{\mathrm{an}}$. At a
maximal ideal of the Tate algebra, its local ring is regular of
dimension $n$. Every maximal localization of the finite algebra
above it is Cohen--Macaulay of dimension $n$, and its fiber is
zero dimensional. Miracle flatness, in the form of
\cite[Tag 00R4]{Stacks}, gives flatness at these localizations.
A finite module is flat if and only if its localizations at
maximal ideals are flat, so the map is finite flat. Its injection
and integrality make the map on spectra, hence the finite map
of affinoid spaces, surjective. The regular extension of $g_i$
is the displayed polynomial expression.
\end{proof}

\begin{lemma}[Algebraic norm boxes]
\label{lem:r5-algebraic-norm-boxes}
Let $Y$ be a smooth affine $K$-variety and let $U$ be an ordinary
open in $Y^{\mathrm{an}}$. Every $x\in U$ has a strictly affinoid
neighborhood $V\subset U$ which is the generic fiber of the
$\pi$-adic completion of an integral affine finite-type flat
$K^\circ$-scheme $\mathcal M$, after restricting to one irreducible
component at $x$. Its algebraic generic fiber is an affine
Zariski open $Y_0\subset Y$, and $\mathscr O(V)$ admits a
finite Tate-presentation generating tuple whose entries are
restrictions of regular functions on $Y_0$.
The point $x$ is in the ordinary interior of $V$.
\end{lemma}

\begin{proof}
Use the topology of the analytification of an affine variety:
absolute values of finitely many regular functions give basic
ordinary neighborhoods. Choose strict upper and, where needed,
strict positive lower bounds at $x$, within the prescribed
neighborhood. The bounds can be chosen in
$|K^\times|^{\mathbf Q}$, since its logarithm is a dense
subgroup of $\mathbf R$. Add sufficiently loose upper bounds
for a finite set of affine coordinates to make the closed box
compact. Replacing the bounds by non-strict inequalities gives
an affinoid neighborhood still contained in $U$; the strict
inequalities at $x$ ensure that $x$ is an interior point.

After inverting the functions with positive lower bounds,
we are on an affine Zariski open $Y_0$. Raising a defining
inequality to a positive integral power and multiplying by
an integral power of $\pi$ writes it as $|a_i|\leq1$, with
$a_i$ regular on $Y_0$. Add a scaled finite generating set
of $K[Y_0]$ whose norms are at most one on the box. These
additional inequalities do not alter the box. Put
\[
             M_0=K^\circ[a_1,\ldots,a_m]\subset K[Y_0].
\]
Then $M_0[1/\pi]=K[Y_0]$ and $M_0$ is a finite-type,
torsion-free $K^\circ$-algebra. It is flat over the DVR
$K^\circ$. Its $\pi$-adic completion has generic fiber
the affinoid domain defined by the chosen inequalities:
the universal affinoid presentation adjoins precisely the
displayed power-bounded generators. Thus
$\mathcal M=\operatorname{Spec}M_0$ has the required properties.
Smooth components of $Y$ are disjoint locally; restrict to the
component containing $x$ to obtain integrality. This construction
uses the original absolute value and makes no normalization
$|\pi|=e^{-1}$.
\end{proof}

\begin{proposition}[Semistable covers with an algebraic ambient map]
\label{prop:r5-semistable-overconvergent-cover}
Let $V\subset U\subset Y^{\mathrm{an}}$ be an algebraic norm box
as in Lemma~\ref{lem:r5-algebraic-norm-boxes}. There exist a
finite extension $K'/K$, a smooth $K'$-variety $Z$, a proper
surjective $K$-analytic map
\[
                         q:W\longrightarrow V
\]
from a compact analytic domain $W\subset Z^{\mathrm{an}}$,
and a finite affinoid cover $W=\bigcup_\lambda W_\lambda$
with the following properties.
\begin{enumerate}
\item The map $q$ is the restriction of an algebraic map
$Z\to Y_0$ over $K$. The inverse image of $U$ is therefore
an ordinary boundaryless neighborhood of the whole of $W$.
\item Each $W_\lambda$ is the generic fiber of the completion
of a small affine open in a strictly semistable algebraic
$K'^\circ$-scheme. Its topological affinoid generators extend
to regular functions on a smooth affine neighborhood in $Z$.
\item Each chart has regular functions $x_0,\ldots,x_r$
defining the irreducible components of its reduced special
fiber, with $x_0\cdots x_r=\varpi$ for a uniformizer
$\varpi$ of $K'$. Its Shilov boundary is precisely the
corresponding set $\{\xi_0,\ldots,\xi_r\}$ of divisorial
points. The continuous logarithmic functions
\[
 \lambda_i=-\frac{\log|x_i|}{-\log|\varpi|},\qquad
 \sum_{i=0}^r\lambda_i=1
\]
satisfy $\lambda_i(\xi_j)=\delta_{ij}$. Each $x_i$ is a
unit on the algebraic generic-fiber neighborhood.
\end{enumerate}
In particular, the compact charts occur inside actual ordinary
ambient neighborhoods. They are not treated as boundaryless
merely because they have formal models.
\end{proposition}

\noindent\textbf{Status: \texttt{DIRECT\_COROLLARY}} of de Jong's
Theorem 6.5 \cite{dJ96} and
Lemma~\ref{lem:r5-local-vertical-interpolation}. Strict
semistability here has de Jong's meaning: the total scheme is
regular and its special fiber is a reduced strict normal-crossing
divisor. No assertion about an entire skeleton is needed.

\begin{proof}
Apply \cite[Theorem 6.5]{dJ96} to $\mathcal M$ and its special
fiber. It gives a finite extension of traits
$\operatorname{Spec}K'^\circ\to\operatorname{Spec}K^\circ$,
an alteration $\mathcal N\to\mathcal M$, and an open embedding
of $\mathcal N$ in a projective strictly semistable scheme.
Thus $\mathcal N$ itself is strictly semistable and is proper
over $\mathcal M$. Its generic fiber $Z$ is smooth over $K'$
and has an algebraic proper surjection to $Y_0$.

Let $W$ be the generic fiber of the completion of $\mathcal N$
along its special fiber. It is exactly the inverse image of
$V$ in $Z^{\mathrm{an}}$. To check this at an actual point
$z$, the condition that its image belong to $V$ says that
the corresponding map to $\mathcal M$ extends over
$\mathscr H(z)^\circ$. The valuative criterion for the proper
map $\mathcal N\to\mathcal M$ uniquely extends the map to
$\mathcal N$. Its center belongs to the special fiber, which
is exactly the condition $z\in W$. The reverse inclusion
follows directly from the morphism of completions. The
analytic map $W\to V$ is proper and surjective and $W$
is compact.

Apply Lemma~\ref{lem:r5-local-vertical-interpolation} at every
point of the special fiber of $\mathcal N$. Quasicompactness
gives finitely many of its small algebraic affine charts.
Completing them gives the asserted finite affinoid cover.
Their finite-type coordinate rings supply regular topological
generators on their algebraic generic fibers. The lemma gives
the exact Shilov boundary and the logarithmic coordinate
values over the same field $K'$. In particular no algebraic
closure or splitting of strata is involved. Every chart maps
into $V$, while its ordinary algebraic neighborhood maps into
$U$ after restricting to the inverse image of $U$. This proves
the ambient assertion.
\end{proof}

\begin{theorem}[Continuity from actual curvewise harmonicity]
\label{thm:r5-continuity-from-curves}
Let $f:U\to\mathbf R$ be an actual real-valued function on an
allowed $K$-analytic space. Suppose that after every complete
extension $F/K$, its pullback along every analytic map from
an ordinary open in a smooth algebraic $F$-curve is an actual
continuous harmonic function. Then $f$ is continuous and therefore
locally bounded.
\end{theorem}

\noindent\textbf{Status: \texttt{NEW\_PROVED}}, using
Theorem~\ref{thm:r5-curvewise-shilov-without-usc}.
No semicontinuity, local boundedness, or Borel measurability of
$f$ is among the assumptions. Finiteness of its actual values
and actual harmonicity on the indicated curves are essential.

\begin{proof}
Zero-dimensional smooth spaces are locally discrete, and their
case is immediate. In positive dimension, B, Lemma D.4,
allows us to work in an ordinary open of a smooth affine
$K$-variety. Around a given point choose a norm box
$V\Subset U$ with that point in its ordinary interior.
Use Proposition~\ref{prop:r5-semistable-overconvergent-cover}.
For a chart $W_\lambda$, denote its Shilov points
by $\xi_0,\ldots,\xi_r$ and put
\[
 h_\lambda=\sum_{i=0}^r(f\circ q)(\xi_i)\lambda_i .
\]
This is continuous. It is a real linear combination of logarithms
of analytic units, so its pullback to every smooth curve
is harmonic. The same is true of $f\circ q-h_\lambda$
by the assumed curve property of $f$. The difference vanishes
on every Shilov point of $W_\lambda$.

The affinoid generators of $W_\lambda$ extend to its smooth
algebraic neighborhood. Lemma
\ref{lem:r5-polynomial-noether-coordinates} therefore gives
finite flat Noether coordinates extending to that neighborhood.
Apply Theorem~\ref{thm:r5-curvewise-shilov-without-usc}
over $K'$ to the difference.
Its absolute value is bounded by its maximum on the Shilov
boundary, which is zero. Hence
\[
                      f\circ q=h_\lambda
                       \quad\hbox{on }W_\lambda.
\]
The finitely many affinoid domains $W_\lambda$ are compact,
hence closed in the Hausdorff compact space $W$. Closed-set
pasting shows that $f\circ q$ is continuous on $W$. Finally
a continuous surjection from a compact space onto the Hausdorff
space $V$ is a quotient map. Since $f\circ q$ is constant on
its fibers by its definition, $f|_V$ is continuous. The
ordinary interiors of these $V$ cover $U$, proving the claim.
\end{proof}

\paragraph{Logical boundary.}
This argument has not used upper semicontinuous regularization
of $f$, convergence of moving coefficient measures, pullback
of an arbitrary current, or a strictly psh exhaustion of
every analytic domain. Its indispensable analytic input is
the Shilov principle for the actual curvewise harmonic
function on compact slices with boundaryless ambient
neighborhoods.


% END source: sections/algebraic_continuity_geometry.tex


% BEGIN canonical R5 source: sections/arbitrary_background_quotient.tex
\section{The arbitrary-background potential quotient}
\label{sec:r5-arbitrary-background-quotient}

Fix a finite extension $K/\mathbf Q_p$ with its original absolute
value. Let $U$ be an ordinary open in a good Hausdorff boundaryless
quasi-smooth $K$-analytic space of pure dimension $n$.
Throughout this section $D=d'd''$, the test forms are all original
compactly supported CLD smooth complementary forms, and positivity
is the original positive cone used in R. Every positive current
under consideration is symmetric. Strong and weak positivity are
not substituted for that cone.

Write $\mathscr{DSH}_{\mathrm{lb}}(U)$ for the actual real functions
which are locally differences of two members of $\mathcal C_b$.
A function in $\mathscr{PH}(U)$ is a finite continuous function
whose two signs are classical psh. Such a function is locally
bounded. No quotient by a measure-null set occurs in either
definition.

\begin{theorem}[Regularity of the bounded classical curvature kernel]
\label{thm:r5-bounded-kernel-regularity}
If $u,v\in\mathcal C_b(U)$ and $D[u]=D[v]$ on every original
test, then the actual difference $u-v$ is continuous and belongs
to $\mathscr{PH}(U)$. More generally,
\[
 \ker\bigl(D:\mathscr{DSH}_{\mathrm{lb}}(U)
                   \longrightarrow\mathcal D'^{1,1}(U)\bigr)
                         =\mathscr{PH}(U).
\]
The equality is an equality of spaces of actual functions.
It does not assert regularity for the kernel on all locally
integrable functions or on all abstract distributions.
\end{theorem}

\noindent\textbf{Status: \texttt{NEW\_PROVED}} in the stated
bounded-classical scope, relative to the declared internal
inputs and the preceding proofs.

\begin{proof}
Let $H=u-v$. Theorem~\ref{thm:r5-all-extension-curve-harmonicity}
shows that after every complete extension $F/K$, the pullback
of $H$ along every allowed smooth analytic curve is an actual
continuous harmonic function. In particular it has the
algebraic-curve property required by
Theorem~\ref{thm:r5-continuity-from-curves}. That theorem applies
to the actual locally bounded $H$ and proves its ordinary
continuity on $U$.

Both $H$ and $-H$ are therefore upper semicontinuous. To check
the full curve quantifier, let $C\to U_F$ be an arbitrary
$F$-analytic curve map, including curves with boundary or
singularities. Its two pullbacks are continuous. In the
definition of subharmonicity on $C$, test further over any
complete extension $E/F$ by a smooth connected strict
$E$-curve mapping to $C_E$. The composite maps to $U_E$,
so the preceding harmonicity theorem makes both signs
subharmonic on that test curve. Smooth curves can be
localized in good Hausdorff boundaryless ordinary opens,
and subharmonicity is local. Definition 3.4 and \S4.1.3
of \cite{GR25} therefore give subharmonicity on the original
$C$. Thus the full classical psh condition holds for both
signs, including its boundary and nonsmooth curve tests.
They belong to $\mathcal C_b(U)$, and
$H\in\mathscr{PH}(U)$.

For a function in $\mathscr{DSH}_{\mathrm{lb}}$ with zero
curvature, apply this argument on its local difference
presentations. Continuity and classical pluriharmonicity
are local and hence glue. Conversely, for
$H\in\mathscr{PH}$ the already established first-order
positive-current theorem applied to $H$ and $-H$ shows that
$D[H]$ and $-D[H]$ are positive. A symmetric current and its
negative can both be positive only if it is zero: on each
finite moment presentation every symmetric smooth test is
dominated in both signs by a positive control. Thus $D[H]=0$.
\end{proof}

\begin{theorem}[Potential independence for every fixed positive closed background]
\label{thm:r5-arbitrary-background-potential-quotient}
Let $0\leq p<n$, and let $S$ be any symmetric original-positive
$(p,p)$-current on $U$, separately $d'$-closed and $d''$-closed.
For any $u,v\in\mathcal C_b(U)$,
\[
 D[u]=D[v]\quad\Longrightarrow\quad D(uS)=D(vS).
 \tag{Q}
\]
All products by $u$ and $v$ are the actual coefficient-Radon
Borel products. The conclusion holds on every original compact
smooth test of bidegree $(n-p-1,n-p-1)$, including all original
moment refinements.

No smoothness, model origin, analytic-cycle representation, or
pre-existing BT representation is required of $S$.
For $p=n$ the corresponding differentiated statement is vacuous.
\end{theorem}

\noindent\textbf{Status: \texttt{NEW\_PROVED}}. The final
positivity step is a direct use of the previously established
arbitrary-fixed-background first-order construction; the new
input is Theorem~\ref{thm:r5-bounded-kernel-regularity}.

\begin{proof}
Put $H=u-v$. Theorem~\ref{thm:r5-bounded-kernel-regularity}
gives $H,-H\in\mathcal C_b(U)$ as actual functions.
Apply the existing first-order BT construction, or explicitly
Proposition~\ref{prop:r5-prescribed-relative-first-order} to
the B-prescribed rows of $H$ and $-H$. For the arbitrary
fixed background $S$ it gives two positive separately
closed currents
\[
                         D(HS),\qquad D((-H)S).
\]
Borel multiplication and distributional differentiation are
linear, so the second is the negative of the first.
The positive-cone argument at the end of the preceding proof
makes this current zero on every original complementary test.
Linearity now gives $D(uS)-D(vS)=D(HS)=0$.

Only the already defined multiplication of a fixed positive
current by a locally bounded Borel function was used. In
particular no pullback of $S$, no regularization of $S$, and
no moving-background limit occurred.
\end{proof}

\begin{corollary}[Differentiated-slot descent]
\label{cor:r5-differentiated-slot-quotient}
Let $S$ be as above, let $0\leq k\leq n-p$, and let
$u_1,\ldots,u_k,v_1,\ldots,v_k\in\mathcal C_b(U)$.
If $D[u_i]=D[v_i]$ for every $i$, then
\[
              T_{S,k}(u_1,\ldots,u_k)
                    =T_{S,k}(v_1,\ldots,v_k)
\]
on all original tests.
\end{corollary}

\noindent\textbf{Status: \texttt{DIRECT\_COROLLARY}} of (Q)
and the frozen fixed-background BT calculus.

\begin{proof}
Replace one differentiated slot at a time. By symmetry of
the established mixed product, place the slot being changed
first. The product of the remaining slots with $S$ is one
fixed positive separately closed background. Apply (Q) to
this background and the two potentials in that slot.
After the finitely many replacements the displayed equality
follows. The case $k=0$ is the identity $S=S$.
\end{proof}

\begin{definition}[Positive currents with bounded classical local potentials]
\label{def:r5-classical-potential-current}
An original $(1,1)$-current $\alpha$ on $U$ has bounded
classical local potentials if every point has an ordinary
neighborhood $O$ and $a\in\mathcal C_b(O)$ such that
$\alpha|_O=D[a]$.
\end{definition}

These currents are positive and separately closed by B.
The definition is a local-potential condition, and it is not
asserted for all positive closed $(1,1)$-currents.

\begin{corollary}[A current-valued mixed operation with one arbitrary background]
\label{cor:r5-local-potential-current-products}
Let $\alpha_1,\ldots,\alpha_k$ have bounded classical local
potentials on $U$, with $k\leq n-p$. There is a well-defined
original positive separately closed current
\[
                        \alpha_1\wedge\cdots\wedge\alpha_k\wedge S
\]
whose expression on any ordinary open carrying local potentials
$\alpha_i=D[a_i]$ is $T_{S,k}(a_1,\ldots,a_k)$.
It is independent of all local potentials and is symmetric
in the $\alpha_i$. It has the restriction, gluing, additivity,
and positive homogeneity properties inherited from the
bounded-potential BT calculus.
\end{corollary}

\noindent\textbf{Status: \texttt{DIRECT\_COROLLARY}}.

\begin{proof}
Choose simultaneous local potentials on an ordinary cover.
On an overlap, both local potentials of each $\alpha_i$
have the same ambient original curvature current.
Corollary~\ref{cor:r5-differentiated-slot-quotient} gives equality
of the two local mixed currents on every original test.
The ordinary sheaf property for currents glues them uniquely.
Their positivity and separate closedness are local.
Symmetry and the linearity assertions follow from those
properties for the chosen local potentials; potential
independence removes the dependence on those choices.
\end{proof}

\begin{corollary}[Pullback of bounded-potential curvature currents]
\label{cor:r5-potential-current-pullback}
Let $\alpha$ have bounded classical local potentials on an
allowed original $K$-space $X$. Let $L/K$ be any complete
extension and let $F:Z\to X_L$ be analytic with $Z$ allowed
over $L$. The local rule
\[
                      F^*\alpha=D[a_L\circ F],
                \qquad \alpha=D[a]\ \hbox{locally},
\]
defines an original positive separately closed $(1,1)$-current
on $Z$, independent of the chosen original-$K$ potentials.
Iterating the rule with the pullbacks of the original potentials
agrees with further scalar extension and analytic composition.
For maps between allowed original-$K$ spaces, this is the usual
intrinsic functoriality on the stated local-potential class.
\end{corollary}

\noindent\textbf{Status: \texttt{DIRECT\_COROLLARY}} of
Theorem~\ref{thm:r5-all-extension-analytic-kernel}.
This is a pullback for the specified local-potential class,
not for an arbitrary positive current.

\begin{proof}
On each inverse-image neighborhood, the actual function
$a_L\circ F$ has a prescribed smooth positive row pulled
back from its original-$K$ row. Its original curvature is
therefore positive and separately closed by the fixed-row
first-order argument. Two choices of $a$ on an original
overlap have equal ambient curvature; the all-extension
analytic kernel theorem gives equal displayed currents on
the whole inverse-image overlap. They glue uniquely.
Composition and extension commute with the actual
compositions of functions and with their original smooth
rows. Applying the same independence statement to the
original potentials proves the asserted compatibility.
When all intermediate spaces are over the original $K$,
the theorem also applies to every alternative bounded
classical potential on each intermediate space, giving
intrinsic functoriality there. Over an arbitrary new field,
the assertion about iteration is restricted to the induced
original-potential presentations; it is not an extension
of B to all new functions over that field.
\end{proof}

\paragraph{Undifferentiated weights and limits.}
The quotient is for differentiated slots. An undifferentiated
weight is still an actual function: replacing it by another
function with the same curvature need not preserve a weighted
current. For example $D[0]=D[1]=0$, but $0S=0$ and $1S=S$
for every nonzero $S$. The seven established fixed-background
BT properties retain their original functional hypotheses.
This section does not turn arbitrary weak convergence of
curvature currents or arbitrary varying backgrounds into
weighted joint convergence.

\paragraph{Original tests and the remaining boundaries of scope.}
The operation above is an equality of ordinary original
CLD currents. It is not promoted to equality on all delta-test
forms. It does not define the wedge of two completely arbitrary
positive closed currents. An arbitrary permitted $S$ may still
charge an ambient pluripolar set. Neither the regularity
theorem nor the quotient proof requires that such concentration
be excluded.


% END source: sections/arbitrary_background_quotient.tex


% BEGIN canonical R5 source: sections/signed_potential_current_products.tex
\section{Signed currents with bounded classical difference potentials}
\label{sec:r5-signed-potential-current-products}

The potential quotient also permits a real multilinear operation
on a specified signed local-potential class.  It does not provide
local potentials for arbitrary closed currents.

\begin{definition}[The signed local-potential class]
\label{def:r5-signed-local-potential-current}
Let $U$ be an allowed ordinary open over the original finite
extension $K/\mathbf Q_p$.  Write
$\mathscr V^{1,1}_{\mathrm{cl}}(U)$ for the real original
$(1,1)$-currents $\alpha$ such that every point has an ordinary
neighborhood $O$ on which
\[
             \alpha|_O=D[H],\qquad
             H=a-b,\quad a,b\in\mathcal C_b(O).
\]
Equivalently, locally $\alpha$ is a difference of two positive
currents admitting bounded classical psh potentials.  The
equivalence is local; a global positive-minus-positive
decomposition on $U$ is not required.
\end{definition}

Every current in this class is symmetric, separately closed,
and of coefficient-Radon type.  These properties follow locally
from $\alpha=D[a]-D[b]$.  The class is a real vector space,
because sums and nonnegative multiples of classical psh
functions remain classical, while a negative coefficient can
be absorbed by exchanging the two terms.

\begin{theorem}[Signed differentiated slots over one arbitrary positive background]
\label{thm:r5-signed-potential-current-products}
Let $S$ be any symmetric original-positive $(p,p)$-current,
separately $d'$- and $d''$-closed, and let $0\leq k\leq n-p$.
There is a well-defined coefficient-Radon separately closed
current
\[
        \alpha_1\wedge\cdots\wedge\alpha_k\wedge S,
       \qquad \alpha_i\in\mathscr V^{1,1}_{\mathrm{cl}}(U),
\]
which is symmetric and real multilinear in the differentiated
slots.  It commutes with ordinary restriction and gluing.

On a simultaneous local presentation
$\alpha_i=D[a_i]-D[b_i]$, with $a_i,b_i\in\mathcal C_b$, it
is given by the finite polarization
\begin{equation}
 \sum_{\varepsilon\in\{0,1\}^k}
   (-1)^{\varepsilon_1+\cdots+\varepsilon_k}
   T_{S,k}(c_{1,\varepsilon_1},\ldots,c_{k,\varepsilon_k}),
 \qquad c_{i,0}=a_i,\quad c_{i,1}=b_i.
 \label{eq:r5-signed-potential-polarization}
\end{equation}
The result is independent of the actual difference potentials,
their presentations as differences of classical functions, and
the chosen positive-minus-positive current decompositions
within the displayed local-potential class.  The equalities
hold on every original compact complementary smooth test,
including every original moment refinement.
\end{theorem}

\noindent\textbf{Status: \texttt{DIRECT\_COROLLARY}} of
Theorem~\ref{thm:r5-arbitrary-background-potential-quotient}
and the frozen cone-additive symmetric BT calculus.

\begin{proof}
First choose one simultaneous local presentation.  The finite
sum \eqref{eq:r5-signed-potential-polarization} is a
coefficient-Radon separately closed current, since each of its
terms has these properties.  Symmetry follows by permuting
the slots in the known positive-potential mixed product and
reindexing the finite sum.

We prove independence in one slot, holding the others fixed.
Expand the other slots in their positive and negative choices.
For each resulting term their mixed product with $S$ is one
fixed symmetric positive separately closed current $Q$.
Thus it suffices to prove independence of
\[
                              D(aQ)-D(bQ)
\]
whenever $D[a]-D[b]$ is fixed.  Suppose another allowed
presentation gives the same current:
\[
                   D[a]-D[b]=D[\widetilde a]-D[\widetilde b].
\]
Then the two classical psh functions
$a+\widetilde b$ and $\widetilde a+b$ have equal ambient
original curvature.  Applying the arbitrary-background
quotient to this fixed positive $Q$ gives
\[
                D((a+\widetilde b)Q)
                    =D((\widetilde a+b)Q).
\]
Linearity of coefficient-Radon multiplication and
distributional differentiation yields
\[
                D(aQ)-D(bQ)
                    =D(\widetilde aQ)-D(\widetilde bQ).
\]
This proves the required cancellation.  In particular, it
covers the case of two decompositions of the same actual
difference, and the more general case of distinct difference
potentials with the same current.  It also covers any two
positive-minus-positive decompositions whose terms admit the
specified classical local potentials.  Replacing one slot
at a time proves full independence.  The quotient theorem
was applied only to positive backgrounds $Q$; no quotient
for an arbitrary signed background was assumed.

For addition in a slot, present two signed potentials by
$a-b$ and $c-d$ and present their sum by
$(a+c)-(b+d)$.  Cone additivity of $T_{S,k}$, followed by
the displayed cancellation, gives additivity of the signed
operation.  Nonnegative scalar homogeneity follows in the
same way.  A negative scalar exchanges the positive and
negative choices and contributes its sign, proving real
linearity.  This proves real multilinearity in all slots.

Choose such presentations on an ordinary open cover of $U$.
On an overlap, the given currents $\alpha_i$ have not changed,
so the independence just proved makes the two local finite
sums equal on all original tests.  The sheaf property for
currents glues them uniquely.  Separate closedness and
coefficient-Radon type are local; the algebraic properties
and restriction identities consequently glue as well.
When $k=0$, the operation is simply $S$.
\end{proof}

\paragraph{Positivity and scope.}
If every differentiated slot admits a classical psh local
potential, the operation agrees with the already positive
BT product.  For the general signed class above, positivity
is not part of the assertion.  In particular, this section
does not claim that a positive current of the form $D[H]$
with a bounded classical difference potential necessarily
has a classical psh potential.  It does not claim that all
positive closed currents, or all signed closed currents,
belong to $\mathscr V^{1,1}_{\mathrm{cl}}$.

There is one arbitrary positive separately closed background
and finitely many differentiated local-potential slots.
No product of two entirely arbitrary currents is defined.
The construction is an equality of original ordinary CLD
currents, with no promotion to all delta-test forms or to
the general nonclosed, nonpositive, off-diagonal GK calculus.

% END source: sections/signed_potential_current_products.tex


% BEGIN canonical R5 source: sections/induced_potential_background_quotient.tex
\section{Arbitrary backgrounds over extensions for induced potentials}
\label{sec:r5-induced-potential-background-quotient}

In this section the potentials originate over the fixed finite
extension $K/\mathbf Q_p$, but the background current need not.
Every scalar extension carries the extending absolute value.
The test category over each field is its full original CLD smooth
test sheaf, and positivity is the original positive cone.

\begin{corollary}[The quotient after an arbitrary extension and analytic map]
\label{cor:r5-induced-potential-background-quotient}
Let $U$ be an ordinary open in an allowed original-$K$ space.
Let $u,v\in\mathcal C_b(U)$ satisfy $D_K[u]=D_K[v]$ on all
original $K$-tests.  Let $L/K$ be any complete extension, let
$Z$ be a good Hausdorff boundaryless quasi-smooth $L$-analytic
space of pure dimension $m$, and let
\[
                            F:Z\longrightarrow U_L
\]
be analytic.  For every $0\leq p<m$ and every symmetric
original-positive $(p,p)$-current $S$ on $Z$ which is separately
$d'$- and $d''$-closed, one has
\begin{equation}
 D_L\bigl((u_L\circ F)S\bigr)
          =D_L\bigl((v_L\circ F)S\bigr).
 \label{eq:r5-induced-potential-background-quotient}
\end{equation}
The equality holds on every original compact smooth
$(m-p-1,m-p-1)$-test on $Z$.  Multiplication uses the actual
coefficient-Radon Borel functions.  The current $S$ is arbitrary
over $L$; it need not descend to $K$ or be the pullback of any
current.
\end{corollary}

\noindent\textbf{Status: \texttt{DIRECT\_COROLLARY}} of
Theorem~\ref{thm:r5-bounded-kernel-regularity} and the
field-independent prescribed-row first-order construction.

\begin{proof}
The original-field regularity theorem gives an actual finite
continuous function $H=u-v$ with $H,-H\in\mathcal C_b(U)$.
The authorized original-$K$ regularization supplies prescribed
smooth positive rows for both signs, locally on fixed ordinary
neighborhoods.  Pull each such row first to $U_L$ and then by
$F$.  Its members remain original smooth with positive
$(1,1)$-curvature forms and decrease to the corresponding
actual function $h=H_L\circ F$ or $-h$.  The rows are locally
uniformly bounded.  This uses specified rows from $K$, not a
regularization theorem for arbitrary new functions over $L$.

Proposition~\ref{prop:r5-prescribed-relative-first-order}
applies over $L$ to the fixed background $S$ and each of these
two rows.  Thus $D_L(hS)$ and $D_L((-h)S)$ are positive
symmetric currents.  They are opposites by linearity of the
actual Borel product and distributional differentiation.
Finite positive controls on all original moment refinements
therefore make them zero.  Since
$h=(u_L\circ F)-(v_L\circ F)$, this proves
\eqref{eq:r5-induced-potential-background-quotient}.
Ordinary locality glues the conclusions from the row
neighborhoods.  No operation on $S$ other than its fixed
coefficient-Radon multiplication was used.
\end{proof}

\begin{corollary}[Ordered products of induced local potentials]
\label{cor:r5-induced-ordered-slot-descent}
Keep $L,Z,F,S$ as above, allowing $0\leq p\leq m$, and let
$0\leq k\leq m-p$.  For
$u_1,\ldots,u_k\in\mathcal C_b(U)$ put
$a_i=u_{i,L}\circ F$.  The ordered recursion
\[
 \begin{split}
 T^{\rm ord}_{S,0}&=S,\\
 T^{\rm ord}_{S,k}(a_1,\ldots,a_k)
   &=D_L\bigl(a_1T^{\rm ord}_{S,k-1}(a_2,\ldots,a_k)\bigr)
 \end{split}
\]
defines a symmetric original-positive separately closed current
on all original tests and is independent of the chosen
prescribed rows.  Here \emph{symmetric} refers to the CLD
involution of the resulting current, not to permutation of
the ordered slots.

If also $v_1,\ldots,v_k\in\mathcal C_b(U)$ and
$D_K[u_i]=D_K[v_i]$ for every $i$, then
\begin{equation}
 T^{\rm ord}_{S,k}(u_{1,L}\circ F,\ldots,u_{k,L}\circ F)
   =T^{\rm ord}_{S,k}(v_{1,L}\circ F,\ldots,v_{k,L}\circ F).
 \label{eq:r5-induced-ordered-slot-descent}
\end{equation}
Consequently an ordered tuple of original-$K$ curvature currents
with bounded classical local potentials gives a well-defined
ordered operation with the arbitrary background $S$ on $Z$,
by choosing those local potentials on inverse-image ordinary
opens.  The operation respects ordinary restriction and gluing.
\end{corollary}

\noindent\textbf{Status: \texttt{DIRECT\_COROLLARY}}.  No
permutation symmetry or weighted joint convergence theorem over
an arbitrary new field is included in this assertion.

\begin{proof}
Each $a_i$ has a prescribed smooth positive row induced from
its original-$K$ row.  Starting at the innermost slot, apply
Proposition~\ref{prop:r5-prescribed-relative-first-order}.
At each step its background is one already constructed fixed
positive separately closed current.  The actual recursion is
defined using its coefficient-Radon measures, and gives the
claimed positivity, current symmetry, separate closedness,
and independence of row choices.  The same argument works
locally and commutes with ordinary restriction.

To replace slot $i$, keep all other slots fixed and put
\[
 R_i=T^{\rm ord}_{S,k-i}(a_{i+1},\ldots,a_k).
\]
This is one fixed positive separately closed background.
The preceding corollary gives
\[
 D_L\bigl((u_{i,L}\circ F)R_i\bigr)
       =D_L\bigl((v_{i,L}\circ F)R_i\bigr).
\]
Apply successively the unchanged outer operations
$R\mapsto D_L(a_{i-1}R),\ldots,R\mapsto D_L(a_1R)$
to this equality.  Both inputs at every stage are the same
positive current, so their actual Borel products and derivatives
are equal.  No exchange of adjacent slots is used.  Replacing
the finitely many slots proves
\eqref{eq:r5-induced-ordered-slot-descent}; for $k=0$ the
identity is $S=S$.

On an overlap of original local-potential charts, the two
original-$K$ potentials of each curvature current have the same
ambient original curvature.  Apply the proved slot descent
on the inverse-image overlap.  The local ordered currents
therefore agree on all original tests and glue uniquely.
\end{proof}

\paragraph{Exact boundary of the conclusion.}
The background over $L$ is unrestricted within the stated
positive separately closed class.  The potentials, however,
are induced from bounded classical potentials over the original
$K$.  The descent hypothesis is equality of the original
$K$-curvatures; no assertion is made for additional equalities
arising only after pullback over $L$.
Neither result asserts regularization, kernel regularity,
full slot symmetry, or joint monotone BT continuity for all
unknown bounded classical functions over $L$.  No pullback
of $S$, or of a mixed current with arbitrary background, is
defined.  All equalities remain ordinary original-current
equalities rather than equalities on every delta-test form.

% END source: sections/induced_potential_background_quotient.tex


% BEGIN canonical R5 source: sections/exact_kernel_stability.tex
\section{Finite evaluation and stability inside the exact kernel}
\label{sec:r5-exact-kernel-stability}

Write
\[
 \mathscr H_{\mathrm{cl}}(U)
   =\{H\in\mathscr{DSH}_{\mathrm{lb}}(U):D[H]=0\}
\]
for the space of actual bounded-classical local differences in
the original-current kernel.  Theorem
\ref{thm:r5-bounded-kernel-regularity} identifies this space with
$\mathscr{PH}(U)$.  The consequences below use that identification
and known properties of pluriharmonic functions.  Their status is
\texttt{DIRECT\_COROLLARY}; finite-dimensionality and the
pluriharmonic Shilov bound are not new claims of this project.

\begin{corollary}[Finite evaluations control an exact kernel on a compact stage]
\label{cor:r5-exact-kernel-finite-control}
Let $U$ be an allowed original-$K$ space and let $K_0\subset U$
be a nonempty compact subset.  There are finitely many compact
affinoid domains $V_1,\ldots,V_m\subset U$ whose ordinary
interiors cover $K_0$.  Put
\[
             V=\bigcup_{a=1}^m V_a,
       \qquad \Gamma=\bigcup_{a=1}^m\Gamma(V_a).
\]
Then $\Gamma$ is a finite subset of $U$, and every
$H\in\mathscr H_{\mathrm{cl}}(U)$ satisfies
\begin{equation}
       \sup_{x\in V}|H(x)|=\max_{z\in\Gamma}|H(z)|,
       \qquad
       \sup_{x\in K_0}|H(x)|\leq\max_{z\in\Gamma}|H(z)|.
 \label{eq:r5-exact-kernel-finite-control}
\end{equation}
The domains and the finite set can be fixed using only $U$ and
$K_0$; they do not depend on $H$, on a presentation $H=u-v$,
or on a background current.

The space of restrictions
$\mathscr H_{\mathrm{cl}}(U)|_V$ is a finite-dimensional
subspace of $C(V)$, of dimension at most $|\Gamma|$, and is
closed for the uniform norm.
\end{corollary}

\begin{proof}
Goodness supplies an affinoid neighborhood at every point of
$K_0$, chosen inside $U$.  Compactness supplies finitely many
whose ordinary interiors cover $K_0$.  Each Shilov boundary is
finite.  By Theorem~\ref{thm:r5-bounded-kernel-regularity}, the
restriction of $H$ to every $V_a$ is pluriharmonic.  The known
Shilov maximum and minimum principle
\cite[Proposition 7.3]{GR25} gives
\[
              \sup_{V_a}|H|=\max_{\Gamma(V_a)}|H|.
\]
Taking the maximum over the finite cover gives
\eqref{eq:r5-exact-kernel-finite-control}.  Evaluation on
$\Gamma$ is therefore an injective linear map from the space
of restrictions into $\mathbf R^\Gamma$, and is an isometry
for the displayed supremum norms.  Its image is a linear
subspace of a finite-dimensional vector space, hence is
finite-dimensional and closed.  This proves the last assertion.
All domain choices preceded the choice of $H$.
\end{proof}

\begin{corollary}[Finite pointwise limits in the exact kernel]
\label{cor:r5-exact-kernel-pointwise-stability}
Let $(H_i)_{i\in I}$ be a net in
$\mathscr H_{\mathrm{cl}}(U)$ which converges at every actual
point of $U$ to a finite real value $H(x)$.  Then the convergence
is uniform on every compact subset of $U$, and
\[
                  H\in\mathscr H_{\mathrm{cl}}(U)
                       =\mathscr{PH}(U).
\]
Neither continuity of the pointwise limit nor common bounds on
the individual psh presentations of $H_i$ are hypotheses.

In particular, on any quasicompact allowed $U$ the kernel space
itself is finite-dimensional, and the convergence is uniform
on $U$.
\end{corollary}

\begin{proof}
Fix a compact $K_0\subset U$ and the finite set $\Gamma$ of
Corollary~\ref{cor:r5-exact-kernel-finite-control}.  Apply its
inequality to $H_i-H_j$, which is again in the exact kernel.
For fixed $i$, pass to the pointwise limit in $j$.  The maximum
is over a finite set, so for every $x\in K_0$,
\[
 |H_i(x)-H(x)|
       \leq\max_{z\in\Gamma}|H_i(z)-H(z)|.
\]
The right-hand side tends to zero with $i$.  This proves uniform
convergence on $K_0$.  In particular $H$ is continuous locally.
Local uniform limits of pluriharmonic functions are
pluriharmonic by \cite[Proposition 7.2(4)]{GR25}; for a net,
one may choose a uniformly convergent sequence on each of the
fixed compact neighborhoods just constructed.  Hence
$H\in\mathscr{PH}(U)$, and the kernel identification gives the
displayed conclusion.

If $U$ is quasicompact, choose a finite affinoid cover of $U$
itself.  The same finite-evaluation argument applies to the
whole space.  This is the argument underlying
\cite[Theorem 7.4]{GR25}.
\end{proof}

\paragraph{Provenance and limits of the consequence.}
The preceding argument is the finite-dimensional Shilov argument
of \cite[Proposition 7.3 and Theorems 7.4--7.6]{GR25}, applied
after the new identification of the actual original-current
kernel.  The proof of its Theorem 7.6 already contains the
finite-evaluation mechanism; here it is written for nets and
with continuity of the finite pointwise limit derived explicitly.
No novelty is claimed for that elementary consequence.
The finite dimension is asserted on a fixed quasicompact domain
or for restrictions to a fixed compact affinoid cover.  It does
not assert a uniform bound for dimensions as neighborhoods
shrink, or a finite-dimensional stalk of the pluriharmonic sheaf.

Exact vanishing $D[H_i]=0$ is essential to this statement.
Weak convergence $D[H_i]\to0$ alone does not imply it, even
for uniformly bounded continuous classical difference
presentations: Example~\ref{ex:r5-universal-energy-polar-obstruction}
has a nonzero actual polar-carried limit and a persistent
relative curvature defect.  Thus this corollary does not
replace any capacitary or weighted convergence hypothesis by
weak curvature convergence.

% END source: sections/exact_kernel_stability.tex


% BEGIN canonical R5 source: sections/exact_kernel_energy.tex
\section{Cutoff energy and strong stability inside the exact kernel}
\label{sec:r5-exact-kernel-energy}

Let $\Omega$ be a norm-box domain on which the faithful degree-one
energy construction of Section~\textup{R4--84} applies.
Fix a symmetric original-positive, separately $d'$- and
$d''$-closed current $S$ of bidegree $(n-1,n-1)$.  Write
$(\mathcal E_S,\mathscr D_S)$ for its closed form and
$\lambda_S$ for the chosen faithful reference measure.  It is
locally finite; no finite total mass on $\Omega$ is assumed.

\begin{corollary}[Exact cutoff energy identity]
\label{cor:r5-exact-kernel-cutoff-energy}
Let $H\in\mathscr{DSH}_{\mathrm{lb}}(\Omega)$ be actual, with
$D[H]=0$ on all original compact tests.  For every real original
smooth compact function $\chi$ on $\Omega$, one has
$\chi H\in\mathscr D_S$ and
\begin{equation}
        \mathcal E_S(\chi H,\chi H)
                    =\int_\Omega H^2\,\Gamma_S(\chi).
 \label{eq:r5-exact-kernel-cutoff-energy}
\end{equation}
In particular the right-hand side is a fixed coefficient integral,
not an integral against a varying curvature measure.
\end{corollary}

\noindent\textbf{Status: \texttt{DIRECT\_COROLLARY}} of the
quotient theorem and the frozen bounded BT--energy Green and
product identities.

\begin{proof}
Theorem~\ref{thm:r5-bounded-kernel-regularity} makes $H$
continuous and pluriharmonic.  The arbitrary-background quotient
theorem gives $D(HS)=0$.  Proposition
\textup{R4--86.1}, applied to a local bounded
classical difference representing $H$, gives local form
membership and the bounded product identities.

Choose an ordinary relatively compact neighborhood $W$ of
$\operatorname{supp}\chi$ and an original smooth compact
$\theta$ equal to one on a neighborhood of $\overline W$.
For every original smooth compact $\varphi$ supported in $W$,
the Green identity gives
\[
       \mathcal E_S(\theta H,\varphi)
                         =-\langle D(HS),\varphi\rangle=0.
\]
The vector $\chi^2H$ belongs to the part space
$\mathscr D_{S,0}(W)$: the support-preserving approximations in
Proposition~\textup{R4--86.1}, multiplied by
$\chi^2$, are original core functions supported in $W$ and
converge in form norm.  Hence the displayed equality extends to
this test vector.  Since $\theta=1$ near its support, the local
Leibniz identity yields
\[
 0=\mathcal E_S(\theta H,\chi^2H)
   =\int\chi^2\Gamma_S(H)
                       +2\int\chi H\,\Gamma_S(H,\chi).
\]
All terms are finite by local boundedness and energy
Cauchy--Schwarz.  The product square identity gives
\[
 \mathcal E_S(\chi H,\chi H)
   =\int\chi^2\Gamma_S(H)
       +2\int\chi H\,\Gamma_S(H,\chi)
       +\int H^2\Gamma_S(\chi).
\]
Subtracting the preceding equality proves
\eqref{eq:r5-exact-kernel-cutoff-energy}.  Thus no nonsmooth
test was inserted directly into a distributional equation:
it was reached through the closed form and compact core density.
\end{proof}

\begin{corollary}[Finite evaluations control the local form norm]
\label{cor:r5-exact-kernel-energy-stability}
Fix $\chi$ as above.  There is a finite set
$\Gamma\subset\Omega$, depending on $\Omega$ and
$\operatorname{supp}\chi$ but not on $H$ or on its
classical-difference presentation, such that
\begin{equation}
 \begin{split}
  \mathcal E_S(\chi H,\chi H)
       &\leq\left(\int\Gamma_S(\chi)\right)
                              \max_{z\in\Gamma}|H(z)|^2,\\
  \|\chi H\|_{\mathscr D_S}^2
       &\leq\left(\int\Gamma_S(\chi)+\int\chi^2\,d\lambda_S\right)
                              \max_{z\in\Gamma}|H(z)|^2,
 \end{split}
 \label{eq:r5-exact-kernel-energy-control}
\end{equation}
where
$\|f\|_{\mathscr D_S}^2=\mathcal E_S(f,f)+\|f\|_{L^2(\lambda_S)}^2$.
Both constants are finite for the fixed current and cutoff.

Consequently, if a net $(H_i)$ in the actual exact kernel
converges at every point of $\Omega$ to a finite real function
$H$, then $H$ is pluriharmonic and
\[
                     \chi H_i\longrightarrow\chi H
                             \quad\hbox{strongly in }\mathscr D_S
\]
for every fixed $S$ and every such $\chi$.
\end{corollary}

\begin{proof}
If $\chi=0$ there is nothing to prove.  Otherwise apply
Corollary~\ref{cor:r5-exact-kernel-finite-control} to the
compact support of $\chi$.  Its finite set controls the
supremum of $|H|$ on that support.  The positive measure
$\Gamma_S(\chi)$ is carried by this compact set, so
\eqref{eq:r5-exact-kernel-cutoff-energy} gives the first
bound.  The fixed locally finite reference measure gives
the second bound.

For the net, Corollary
\ref{cor:r5-exact-kernel-pointwise-stability} first proves
that $H$ is pluriharmonic.  Apply
\eqref{eq:r5-exact-kernel-energy-control} to $H_i-H$.
The maximum over the fixed finite set tends to zero, proving
the asserted strong convergence.  This requires no common
bound on the separate psh presentations of the net.
\end{proof}

\paragraph{Remaining degree and scope.}
The scalar closed form above has a degree-one background.  Starting
with an arbitrary allowed $(p,p)$ background $S_0$ and
$q=n-p\geq1$, one may freeze any $q-1$ bounded classical slots
and use
\[
               S=T_{S_0,q-1}(u_1,\ldots,u_{q-1})
\]
as that fixed positive closed background.  The same identities
then apply to its faithful form.  They do not assert an
unfrozen all-degree Hilbert form, uniform constants as $S$
varies, or continuity for $S_i\to S$.  Exact membership in
the actual kernel is essential: weakly vanishing ambient
curvatures and uniformly bounded classical decompositions do
not give these conclusions for their assigned limits, as
Example~\ref{ex:r5-universal-energy-polar-obstruction} shows.

% END source: sections/exact_kernel_energy.tex


% BEGIN canonical R5 source: sections/relative_energy_convergence.tex
\section{A fixed-background energy criterion for curvature convergence}
\label{sec:r5-relative-energy-convergence}

Fix the original field $K$, an allowed ordinary open $U$ of
pure dimension $n$, and a symmetric original-positive separately
$d'$- and $d''$-closed current $S$ of bidegree $(p,p)$.
Put $q=n-p\geq1$. All tests below are original compact smooth
CLD forms. For a degree-one background $R$ on an algebraic
norm-box domain $\Omega$, let $(\mathcal E_R,\mathscr D_R)$
denote the faithful closed form of
Theorem~\textup{R4--84.3}. Its values on
bounded classical differences agree with the bounded BT energy
by Proposition~\textup{R4--86.1}.

\begin{proposition}[Finite-test energy control]
\label{prop:r5-finite-test-relative-energy}
For every original compact $(q-1,q-1)$-test $\eta$ on $U$,
there are finitely many ordinary open algebraic norm-box domains
$\Omega_b\Subset U$, original smooth compact real functions
$f_{ba}$ on $\Omega_b$, original smooth compact cutoffs $\chi_b$
equal to one near all the supports of $f_{ba}$, and constant
original-positive closed forms $\theta_{ba}$ of degree
$(q-1,q-1)$ in fixed unit-moment coordinates on $\Omega_b$,
such that the following statements hold.

The currents
\[
             R_{ba}=(S|_{\Omega_b})\wedge\theta_{ba}
\]
are symmetric original-positive separately closed currents of
bidegree $(n-1,n-1)$. For every actual
$H\in\mathscr{DSH}_{\mathrm{lb}}(U)$,
$\chi_bH\in\mathscr D_{R_{ba}}$ and
\begin{equation}
 \langle D(HS),\eta\rangle
       =-\sum_{b,a}\mathcal E_{R_{ba}}(\chi_bH,f_{ba}).
 \label{eq:r5-finite-test-energy-identity}
\end{equation}
Consequently
\begin{equation}
 \bigl|\langle D(HS),\eta\rangle\bigr|
 \leq C_{S,\eta}
       \left(\sum_{b,a}
       \mathcal E_{R_{ba}}(\chi_bH,\chi_bH)\right)^{1/2},
 \qquad
 C_{S,\eta}^2=\sum_{b,a}
                   \mathcal E_{R_{ba}}(f_{ba},f_{ba})<\infty.
 \label{eq:r5-finite-test-energy-bound}
\end{equation}
The boxes, cutoffs, forms, and constant are fixed by $S$ and
the original test; none depends on $H$ or on a regularization
of its classical-difference presentations.
\end{proposition}

\noindent\textbf{Status: \texttt{DIRECT\_COROLLARY}} of the
original positive-square test decomposition, the faithful
closed-form construction, and the bounded Green identity.

\begin{proof}
The current $D(HS)$ is symmetric, so replace $\eta$ by its
symmetric part. A finite original smooth partition localizes
it in moment presentations. Refine that finite localization
to relatively compact ordinary open algebraic norm boxes inside the
presentations. The local algebraic neighborhoods and these
boxes are available over the original $K$ by B, Lemma D.4,
and Lemma~\ref{lem:r5-algebraic-norm-boxes}. Each localized
test has compact support in one such box.
Use strict slack in that construction: $\Omega_b$ is defined
by strict norm inequalities and its closed norm box is compactly
contained in the moment presentation. The faithful energy form
is taken on the boundaryless ordinary open $\Omega_b$; the
closed affinoid box supplies only its compact containment.

Fix one box and write its moment coordinates as
$t_1,\ldots,t_r$. The symmetric coefficient matrix of the
localized test is a smooth real matrix indexed by
$(q-1)$-multi-indices. Constant positive rank-one matrices
span all such symmetric matrices: use the squares of the
basis vectors and the squares of sums of two basis vectors,
and absorb the signed scalar coefficients into $f_{ba}$.
Thus its exact expression is
\[
                         \sum_a f_{ba}\theta_{ba}.
\]
The scalar coefficients can be multiplied by one compact
cutoff equal to one near the localized support, so they
belong to $A_c^{0,0}(\Omega_b)$. Each $\theta_{ba}$ is the
positive signed square of a constant $(q-1,0)$-form in the
$d't_i$. This form need not be decomposable. The construction
is the finite square span of
\eqref{eq:r5-positive-square-test-span}; it uses the full
original positive cone, not just strongly positive squares.
For $q=1$ the square is simply the constant scalar $1$.

The forms $\theta_{ba}$ are separately closed. To check
positivity of $R_{ba}$, test it by an original positive
$(1,1)$-form on any common refinement. Such a form is a sum
of positive squares of one-covectors. Its wedge with
$\theta_{ba}$ is a sum of positive squares of the exterior
products of those one-covectors with the defining
$(q-1)$-covector. It is therefore an original positive test
for $S$. This proves positivity of $R_{ba}$, including all
original refinements. Symmetry and separate closedness follow from
the same properties of $S$ and the constant closed form.

By Proposition~\textup{R4--86.1}, a cutoff
times a bounded classical difference belongs to
$\mathscr D_{R_{ba}}$. For a locally presented difference,
use a finite partition on the compact cutoff support and
the bounded smooth multiplier property of the closed form.
This proves the asserted membership for $\chi_bH$.

The smooth closed product rule and the scalar Green identity
now give, with precisely the displayed sign,
\[
 \begin{split}
 \langle D(HS),f_{ba}\theta_{ba}\rangle
   &=\langle D(HR_{ba}),f_{ba}\rangle\\
   &=\int_{\Omega_b}H\,D(f_{ba}R_{ba})
     =-\mathcal E_{R_{ba}}(\chi_bH,f_{ba}).
 \end{split}
\]
The cutoff is one near the support of $f_{ba}$ and its
derivatives. The last identity is the strong-form extension
of compact smooth Stokes, not an insertion of a nonsmooth
function as an original smooth test. Summing proves
\eqref{eq:r5-finite-test-energy-identity}. Apply
Cauchy--Schwarz in each fixed closed form and then in the
finite index set to obtain
\eqref{eq:r5-finite-test-energy-bound}. Every energy of a
fixed smooth compact coefficient $f_{ba}$ is finite by local
coefficient-Radon control. This also establishes the claimed
independence of the constants.
\end{proof}

\begin{corollary}[Convergence modulo actual pluriharmonic corrections]
\label{cor:r5-energy-modulo-kernel-convergence}
Let $H_j\in\mathscr{DSH}_{\mathrm{lb}}(U)$ be a sequence
or net. Suppose that for every original test $\eta$ and the
finite family supplied by
Proposition~\ref{prop:r5-finite-test-relative-energy}, there
are $h_{b,j}\in\mathscr{PH}(\Omega_b)$ such that
\begin{equation}
 \sum_{b,a}\mathcal E_{R_{ba}}
       \bigl(\chi_b(H_j-h_{b,j}),
             \chi_b(H_j-h_{b,j})\bigr)\longrightarrow0.
 \label{eq:r5-modulo-kernel-energy-hypothesis}
\end{equation}
Then $D(H_jS)\longrightarrow0$ on every original compact
complementary test on $U$.

In particular, if $u_j,u\in\mathcal C_b(U)$ and the
hypothesis holds with $H_j=u_j-u$, then
\[
                           D(u_jS)\longrightarrow D(uS).
\]
No local uniform bound on the separate psh presentations is
required beyond what is needed for each displayed individual
energy; the convergence hypothesis itself is substantive.
\end{corollary}

\noindent\textbf{Status: \texttt{DIRECT\_COROLLARY}} of
the preceding estimate and the arbitrary-background quotient.

\begin{proof}
For each fixed box, actual pluriharmonicity and
Theorem~\ref{thm:r5-arbitrary-background-potential-quotient}
give $D(h_{b,j}S)=0$ there. Equivalently they give
$D(h_{b,j}R_{ba})=0$ for every member of its finite family.
Apply \eqref{eq:r5-finite-test-energy-identity} to each
localized term, subtract that zero current, and then apply
\eqref{eq:r5-finite-test-energy-bound} with the local
difference $H_j-h_{b,j}$. Its right-hand side tends to zero
by \eqref{eq:r5-modulo-kernel-energy-hypothesis}, while the
left-hand side is the pairing of $D(H_jS)$ with the original
test. This proves the statement for every test. Linearity
gives the last assertion.
\end{proof}

\paragraph{What the criterion does and does not require.}
All the degree-one backgrounds are fixed smooth wedges of
the one given $S$; there is no sequence of backgrounds.
The positive square spanning family accounts for every
original complementary test, including tests with new
moment coordinates. The criterion requires energy of the
error, after an actual pluriharmonic correction, to tend
to zero. Strong energy convergence to a different assigned
Borel representative does not satisfy this requirement.
In particular it does not rescue the weak-current or
completed-polar-selector failures recorded earlier.
No assertion about varying mixed BT tails or joint
higher mixed-current convergence follows from this
first-order estimate alone.

\begin{corollary}[No pluriharmonic correction removes the smooth relative defect]
\label{cor:r5-no-ph-energy-repair}
In Example~\ref{ex:r5-smooth-weak-current-defect}, one can
choose an ordinary algebraic norm box $\Omega_0\Subset\Omega$,
a nonnegative original smooth compact scalar $\varphi$ on
$\Omega_0$, and an original smooth compact cutoff $\chi$
equal to one near $\operatorname{supp}\varphi$, with the
following property. Put $R=[Y]|_{\Omega_0}$. There is a
constant $c_0>0$, independent of $N$ and of
$h\in\mathscr{PH}(\Omega_0)$, such that
\[
 \mathcal E_R\bigl(\chi(v_N-w-h),\chi(v_N-w-h)\bigr)
                              \geq c_0.
\]
\end{corollary}

\noindent\textbf{Status: \texttt{DIRECT\_COROLLARY}}.

\begin{proof}
The example supplies the fixed nonzero positive measure
$D(q[Y])$. Choose a nonnegative compact annular scalar
profile detecting it and multiply by a compact original
cutoff in the $z$-direction equal to one near $z=0$.
Its support can be placed inside a relatively compact
algebraic norm box $\Omega_0$ in the bidisc. This gives
$\varphi$ with
\[
                        c=\langle D(qR),\varphi\rangle>0.
\]
Choose $\chi$ as stated. For every $N$, the example gives
$D((v_N-w)R)=-D(qR)$, whereas the quotient theorem gives
$D(hR)=0$. Green and Cauchy--Schwarz in the one fixed
closed form therefore imply
\[
 c=\bigl|\mathcal E_R(\chi(v_N-w-h),\varphi)\bigr|
 \leq
 \mathcal E_R\bigl(\chi(v_N-w-h),\chi(v_N-w-h)\bigr)^{1/2}
 \mathcal E_R(\varphi,\varphi)^{1/2}.
\]
The last fixed energy is finite. It is positive, since a
zero value would contradict $c>0$ in this inequality.
Thus $c_0=c^2/\mathcal E_R(\varphi,\varphi)$ proves the
claim. The conclusion concerns this fixed relative form;
it does not define an intrinsic all-degree quotient norm.
\end{proof}

% END source: sections/relative_energy_convergence.tex


% BEGIN canonical R5 source: sections/normal_gauss_calibration_check.tex
\section{A finite-field check for normal Gauss coarea}
\label{sec:r5-normal-gauss-calibration-check}

This section checks a possible failure mode in the normal trace:
could a finite residue-field extension split after adjoining the
normal Gauss parameter, or create extra correlated points on the
calibrated support?  The answer for this particular product support
is no.  The following finite-dimensional argument requires neither
stability of the valued field nor a ramification--residue-degree
formula.

\begin{lemma}[Finite extensions commute with Gauss completion]
\label{lem:r5-finite-gauss-completion}
Let $F$ be a complete non-Archimedean field, let $E/F$ be a
finite valued extension with its extending absolute value, and let
$r>0$ be any real radius.  Denote by $F_r$ the completion of
$F(T)$ for the radius-$r$ Gauss norm and by $E_r$ the analogous
completion of $E(T)$.  Then the natural homomorphism
\begin{equation}
                    E\otimes_F F_r\longrightarrow E_r
 \label{eq:r5-finite-gauss-completion}
\end{equation}
is an isomorphism of finite $F_r$-algebras.  In particular, its
source is a field and
\[
                         [E_r:F_r]=[E:F].
\]
There is only one point in the Berkovich spectrum of this finite
$F_r$-algebra, with the full displayed algebraic degree.  No
rationality assumption on $r$ is imposed.
\end{lemma}

\noindent\textbf{Status: \texttt{DIRECT\_COROLLARY}} of
finite-dimensional norm equivalence and the Gauss norm.  The
proof is given because it is exactly the field assertion needed
in the product calibration, including in the presence of defect.

\begin{proof}
Choose an $F$-basis $e_1,\ldots,e_d$ of $E$.  There are
constants $0<c\leq C<\infty$ such that
\begin{equation}
 c\max_i|a_i|
       \leq\left|\sum_i e_i a_i\right|
       \leq C\max_i|a_i|\qquad(a_i\in F).
 \label{eq:r5-finite-basis-norm-equivalence}
\end{equation}
For completeness, the lower estimate needs only completeness
of $F$.  Induct on the number of basis vectors.  The span of
the first $d-1$ vectors is complete by the induction hypothesis,
hence closed in $E$.  Its distance from $e_d$ is positive.
Consequently the norm of $\sum_i e_i a_i$ controls $|a_d|$;
the ultrametric triangle inequality then controls the norm of
the first $d-1$ summands, and induction controls their
coefficients.  The upper estimate is the triangle inequality.
This proves \eqref{eq:r5-finite-basis-norm-equivalence} without
an orthogonal-basis assertion.

For polynomials $P_i(T)\in F[T]$, apply
\eqref{eq:r5-finite-basis-norm-equivalence} separately to
every coefficient and then take the weighted maximum over
the powers of $T$.  The result is
\begin{equation}
 c\max_i\|P_i\|_r
   \leq\left\|\sum_i e_iP_i\right\|_r
   \leq C\max_i\|P_i\|_r.
 \label{eq:r5-gauss-basis-norm-equivalence}
\end{equation}
The same inequalities hold for $P_i\in F(T)$ after clearing
one common denominator in $F[T]$ and dividing by its nonzero
Gauss norm.

Algebraically $E\otimes_F F(T)=E(T)$, and the same $e_i$
form an $F(T)$-basis.  To see explicitly that a denominator
in $E[T]$ causes no problem, multiplication by such a nonzero
polynomial is injective on the free $F[T]$-module $E[T]$.
After localizing $F[T]$, it is an injective endomorphism of a
finite-dimensional $F(T)$-space, hence is invertible.  Its
adjugate divided by its nonzero determinant gives its inverse
with coefficients in $F(T)$.  Thus every element of $E(T)$
has the required basis expression.

The Gauss norm on $E(T)$ is multiplicative.  Its completion
$E_r$ is a field: if a Cauchy sequence converges to a nonzero
element in that completion, the norms of its terms eventually
have a fixed positive value, and their inverses are Cauchy by
the identity $x^{-1}-y^{-1}=(y-x)/(xy)$.  Completing
\eqref{eq:r5-gauss-basis-norm-equivalence} identifies this
completion, as a normed vector space, with $F_r^d$ in the
basis $e_i$.  The lower bound proves injectivity of
\eqref{eq:r5-finite-gauss-completion}; the same estimate and
density prove its surjectivity.  Multiplication commutes with
completion, so it is the claimed algebra isomorphism.

Finally a finite field extension of the complete field $F_r$
has its unique extending absolute value.  Its Berkovich
spectrum therefore consists of one point.  The degree is the
vector-space dimension $d$ retained by the displayed basis;
it has not been computed from residue or ramification indices.
\end{proof}

\begin{corollary}[The special product top-form coarea]
\label{cor:r5-product-top-gauss-calibration}
Let $F$ be complete, nontrivially valued, and of characteristic
zero, and let $B$ be a good Hausdorff boundaryless
smooth $F$-analytic space of pure dimension $d$, and let
$\eta\in A_c^{d,d}(B)$ be an original form, with its signed
coefficient Radon measure $\mu_\eta$.  Let $\kappa$ be a
smooth real function of compact support in a radial interval
and put, on the corresponding normal annulus,
\[
 \beta_\kappa
    =\kappa(\log|T|)\,d'\log|T|\wedge d''\log|T|.
\]
Write $\sigma_r(x)$ for the image in
$B\times_F\mathbf A_F^{1,\mathrm{an}}$ of the radius-$r$
Gauss point over $\mathscr H(x)$.  For every bounded Borel
function $c$ on a neighborhood of the indicated compact
product support,
\begin{equation}
 \int c\,\beta_\kappa\wedge\pi^*\eta
  =\int_B\int_{\mathbf R}
        c(\sigma_{e^t}(x))\kappa(t)\,dt\,d\mu_\eta(x).
 \label{eq:r5-product-top-gauss-calibration}
\end{equation}
Thus this product coefficient measure is carried by the
canonical normal Gauss lifts of the base top-form measure.
\end{corollary}

\noindent\textbf{Status: \texttt{DIRECT\_COROLLARY}} of the
original product calibration, with the finite-field step made
explicit above.  This checks the particular formula used in
Lemma~\ref{lem:r5-arbitrary-field-normal-trace}; it is not an
independent-section formula for an arbitrary full analytic fiber.

\begin{proof}
Use finitely many original moment presentations for $\eta$.
On a calibrated top-dimensional cell, choose a $d$-coordinate
monomial projection of full rank.  At each of its relevant
Gauss points the finite source fibers have completed residue
fields $E$ finite over the moment Gauss field $F_0$.  Their
full finite algebra degrees are exactly the degrees retained
by original calibration.  Adjoining the normal coordinate
replaces $F_0$ by its radius-$r$ Gauss completion.  By
Lemma~\ref{lem:r5-finite-gauss-completion}, each finite field
$E$ yields precisely its canonical Gauss completion $E_r$,
with the same full degree.  Consequently there is one normal
Gauss lift for each such base point, and there is no additional
split factor or lost multiplicity on this calibrated support.

The product cell measure in the added coordinate is $dt$ with
multiplicity one.  The preceding field calculation therefore
gives \eqref{eq:r5-product-top-gauss-calibration} for original
smooth scalar coefficients.  Refining moments for such a
coefficient does not change the calibrated support or its
degree.  Signed linearity and the finite partition recover an
arbitrary original top form $\eta$.

The normal Gauss lift is continuous on each compact affinoid
stage and compact radial interval: for a polynomial
$\sum_j a_jT^j$ its norm is $\max_j|a_j(x)|r^j$, and uniform
analytic approximation gives the general assertion.  Hence
the right side defines a finite signed Radon measure on the
compact support.  Equality on original smooth scalar tests,
which are uniformly dense on compact stages, identifies that
measure with the left-hand coefficient measure.  Integration
of a fixed bounded Borel function proves the stated identity.
\end{proof}

\paragraph{The excluded inference.}
The full fiber of the projection to the normal parameter may
contain correlated points, such as those on which $|x-T|$ is
smaller than both $|x|$ and $|T|$.  It has not been replaced
by the image of $\sigma_r$.  The conclusion concerns the
coefficient measure of a \emph{pulled-back base top form}
wedged with a radial normal top form.  An arbitrary form on
the full fiber, involving mixed moments such as $x-T$, need
not have this support.  This precise distinction is what
makes the normal trace compatible with the correlated-point
counterexample recorded in R4.

% END source: sections/normal_gauss_calibration_check.tex


% BEGIN canonical R5 source: sections/finite_etale_collision_check.tex
\section{A finite-etale collision in the original piecewise test class}
\label{sec:r5-finite-etale-collision-check}

The trace construction must retain joint unit norms, allow genuine
coefficient jumps in normal directions, and retain full finite
degrees.  The following computations check these three features.
They are illustrations of the proved construction, not additional
kernel or regularization theorems.

\begin{example}[Two joint norm vectors which collide]
\label{ex:r5-finite-etale-joint-collision}
Let $F$ be a complete nontrivially valued field of characteristic
zero, with its given absolute value.  Choose $0<\epsilon<1$ and put
\[
 U=\{\epsilon<|T|<\epsilon^{-1},\ 4T+1\ne0\}
          \subset\mathbf A_F^{1,\mathrm{an}}.
\]
The inverse image $Z$ under the polynomial map
\[
                    p:Y\longmapsto T=Y(Y-1)
\]
is finite etale of degree two over $U$.  Both $Y$ and $Y-1$
are analytic units on $Z$.  Write
\[
 s=\log|T|,\qquad z_1=\log|Y|,\qquad z_2=\log|Y-1|.
\]
The unordered joint norm vectors, counted with their full geometric
multiplicities over each completed residue field, are
\begin{equation}
 \{(z_1,z_2)\}
 =
 \begin{cases}
   \{(s,0),(0,s)\},&s<0,\\
   \{(0,0),(0,0)\},&s=0,\\
   \{(s/2,s/2),(s/2,s/2)\},&s>0.
 \end{cases}
 \label{eq:r5-joint-collision-vectors}
\end{equation}
For the original smooth source form
\[
                    \alpha=d'z_1\wedge d''z_2,
\]
the ordinary pushforward current is represented by the finite
compatible piecewise form
\begin{equation}
 p_*[\alpha]=[\tau],\qquad
 \tau=
 \begin{cases}
  0,&s<0,\\[2pt]
  \tfrac12\,d's\wedge d''s,&s>0.
 \end{cases}
 \label{eq:r5-joint-collision-trace}
\end{equation}
The representative on the wall $s=0$ is immaterial.  The current
identity is on every compact original complementary test on $U$;
no all-delta-test assertion is being made.
\end{example}

\noindent\textbf{Status: \texttt{DIRECT\_COROLLARY}} of
Lemma~\ref{lem:r5-prescribed-finite-trace} and the elementary
Newton polygon calculation below.

\begin{proof}
The map is finite free of degree two, since its equation is monic
in $Y$.  Its derivative $2Y-1$ vanishes only over $T=-1/4$,
which was removed.  It is therefore finite etale.  The two source
functions are units because their product is the base unit $T$.

If $|T|<1$, the polynomial $Y^2-Y-T$ has root norms $|T|$ and
$1$.  For the small root $|Y-1|=1$; for the other root the
identity $Y(Y-1)=T$ gives $|Y-1|=|T|$.  If $|T|>1$, both
roots have norm $|T|^{1/2}>1$, and $|Y-1|=|Y|$.  If $|T|=1$,
both $Y$ and $Y-1$ have norm one.  These computations prove
\eqref{eq:r5-joint-collision-vectors}, over algebraic closures of
the completed residue fields and hence with exactly the finite
local multiplicities used in the trace theorem.

On $s<0$, each of the two vector maps has one constant coordinate,
so both pullbacks of $d'z_1\wedge d''z_2$ vanish.  On $s>0$,
each vector map is $(s/2,s/2)$ and contributes
$\frac14d's\wedge d''s$.  Their sum is the second line of
\eqref{eq:r5-joint-collision-trace}.  This is precisely the
joint-norm trace formula of the cited lemma.  Local compact
cutoffs make its ordinary-current identity applicable to every
compact test.  A cutoff may also use the original unit $4T+1$;
no radial cutoff which remains nonzero arbitrarily close to the
removed critical value is being declared compact on $U$.

Both tensor restrictions to the wall tangent space vanish, so the
piecewise form is tangentially compatible even though its ambient
normal coefficient jumps.  The unit-radius Gauss point belongs to
$U$: it is not the removed rigid point $T=-1/4$.  Thus this is a
genuine collision within an allowed ordinary neighborhood, not a
limiting event excluded from the space.
\end{proof}

\paragraph{Why marginal records are insufficient.}
For $s<0$, the two individual coordinate records are both
$\{s,0\}$.  Matching them independently as
$\{(s,s),(0,0)\}$ would incorrectly give
$d's\wedge d''s$ instead of zero.  The joint monomial record
for $Y(Y-1)=T$ rules out this matching: its two log norms are
$s,s$, whereas the false matching would give $2s,0$.
The full finite-grid tomography in the trace proof is the
general version of this elementary correction.

\paragraph{Why the coefficient jump is allowed.}
The same example may be multiplied by any additional smooth
analytic factor and the finite map replaced by $p\times\mathrm{id}$.
Then $\tau$ is an intermediate-degree piecewise form on the base.
For a symmetric original-positive separately closed current $R$
of complementary bidegree, the wall estimate in
Lemma~\ref{lem:r5-prescribed-piecewise-tests} gives, on each fixed
compact stage $H$,
\[
 (d's\wedge d''s\wedge R)
                  (H\cap\{|s|\leq\delta\})
                       \leq C\delta\sum_b\langle R,\Theta_b\rangle.
\]
In particular this coefficient measure gives the wall zero mass.
The value assigned to the discontinuous coefficient in
\eqref{eq:r5-joint-collision-trace} therefore does not affect
the pairing.  For a weakly convergent positive closed family,
the fixed-control masses are uniformly bounded, so smoothing
this particular jump has a uniform $O(\delta)$ error.  General
compatible jumps with only one normal factor use the proved
$O(\delta^{1/2})$ matrix Cauchy--Schwarz estimate.  There is
no claim that the trace is an original smooth form.

\paragraph{A degree check with a power map.}
For $T=Y^d$ on a torus annulus, let $s=\log|T|$ and
$z=\log|Y|$.  There are $d$ geometric root-norm records
$z=s/d$, even when they describe a single Berkovich fiber point.
Thus the trace of a scalar $\varphi(z)$ is
$d\varphi(s/d)$, while the trace of
$\varphi(z)d'z\wedge d''z$ is
\[
 d\left(\frac1d\right)^2\varphi(s/d)d's\wedge d''s
       =\frac1d\varphi(s/d)d's\wedge d''s.
\]
The latter agrees with the direct radial integration change of
variables.  Counting only Berkovich fiber points would lose this
factor.  The construction uses the full finite algebra degree;
it never replaces that degree by a ramification--residue product.

% END source: sections/finite_etale_collision_check.tex


% BEGIN canonical R5 source: sections/quotient_scope_checks.tex
\section{An exact closedness test for the potential quotient}
\label{sec:r5-quotient-scope-checks}

The following elementary example checks a hypothesis of the
arbitrary-background theorem.  It is not a counterexample under
that theorem's stated assumptions.

\begin{example}[Merely $D$-closed backgrounds do not suffice]
\label{ex:r5-D-closed-background-failure}
Let $\pi$ be a uniformizer of the original finite extension
$K/\mathbf Q_p$, put $R=|\pi|^{-1}$ and $\ell=\log R>0$, and
let
\[
 A=\{1<|T|<R\}\subset\mathbf A_K^{1,\mathrm{an}},
 \qquad t=\ell^{-1}\log|T|.
\]
Thus $0<t<1$ and $t$ is an original smooth classical
pluriharmonic function on the ordinary annulus $A$.
Set
\[
 u=t,\qquad v=0,\qquad S=(1+t)[1]_A.
\]
The symmetric $(0,0)$-current $S$ is original-positive and
$DS=0$.  Nevertheless
\begin{equation}
 D[u]=D[v]=0,
 \qquad D(uS)-D(vS)=2[d't\wedge d''t]\ne0.
 \label{eq:r5-D-closed-background-failure}
\end{equation}
In particular, the separate $d'$- and $d''$-closedness in
Theorem~\ref{thm:r5-arbitrary-background-potential-quotient}
cannot be replaced by $D$-closedness alone.
\end{example}

\noindent\textbf{Status: \texttt{FALSE}} for the weakened
statement; \texttt{DIRECT\_COROLLARY} of the original smooth
product rule for the displayed construction.

\begin{proof}
The analytic function $T$ is a unit on $A$, so $Dt=0$ and
both signs of $t$ are classical psh.  Since $1+t>0$, the
displayed coefficient multiplication makes $S$ positive.
Smooth Leibniz and closedness of $[1]_A$ give
\[
 DS=D[1+t]=0,
 \qquad D(uS)=D[t+t^2]=2[d't\wedge d''t].
\]
The last current is not zero.  Indeed, for a nonnegative
nonzero $\varphi\in C_c^\infty((0,1))$, original annulus
calibration with the coordinate $x=\log|T|=\ell t$ gives
\[
 \left\langle 2[d't\wedge d''t],\varphi(t)\right\rangle
       =\frac{2}{\ell}\int_0^1\varphi(s)\,ds>0.
\]
The support of $\varphi(t)$ is a compact closed subannulus,
so this is an allowed original test.  The same formulas show
that $d'S=[d't]$ and $d''S=[d''t]$ are not zero.  Hence the
example violates exactly the separate-closedness assumption,
while retaining positivity, symmetry, local boundedness,
the original field, and its prescribed absolute value.
\end{proof}

\begin{remark}[A nonconstant kernel is permitted]
\label{rem:r5-nonconstant-annulus-kernel}
\textbf{Status: \texttt{DIRECT\_COROLLARY}.}
On the same annulus, let $S$ instead be any background satisfying
the actual theorem's separate-closedness assumptions, on $A$ or
on its product with another allowed space.  The smooth product
rule then gives $D(tS)=Dt\wedge S=0$.  Thus the theorem recovers
the nonconstant logarithmic kernel; it does not imply that a
curvature-kernel function is locally constant.  It also allows
backgrounds carried by analytic subvarieties of that product,
since a restriction of the unit $T$ is still a unit there.
\end{remark}

\paragraph{Actual type-IV values.}
The endpoint argument has a separate check.  On a smooth curve
over an algebraically closed complete extension, let $x$ be of
type IV and choose nested closed discs with Gauss points
$\xi_j\to x$.  For either of the original bounded subharmonic
pullbacks $w$ separately, the disc maximum principle and upper
semicontinuity give
\[
       w(x)\leq w(\xi_j),\qquad
       \limsup_j w(\xi_j)\leq w(x).
\]
Thus both original functions have their correct actual endpoint
values, and so does their difference.  This checks the passage
in Lemma~\ref{lem:r5-actual-curve-harmonic-representative}; it
does not treat a type-IV point as a Gauss point or infer a
Borel-weighted limit from weak convergence of fiber measures.

\paragraph{No faithfulness condition on the background.}
The final argument for $S$ uses that a symmetric original current
and its negative cannot both be positive unless they are zero.
It does not require the coefficient measures of $S$ to detect
every point or every ordinary open set.  Thus $S=0$, backgrounds
with smaller support, and integration currents carried by
ambient pluripolar subvarieties all remain within scope.
At $p=n$ there is no further differentiated degree; the statement
is vacuous rather than an additional analytic assertion.

\paragraph{Finite extensions retain the original calibration.}
For a finite extension $F/K$ of degree $d$, the integration
current of an $F$-space viewed as a $K$-space is $d$ times its
$F$-calibrated current on common original tests.  By contrast,
scalar pushforward of the $F$-calibrated top-form measure has
factor one.  These are respectively the conventions of
\cite[Remark 14.4.3, p.~410, and Proposition 15.3.1, p.~425]{CLD25}.
They are compatible: the finite morphism viewed over $K$ has
degree $d$.  The kernel arguments divide a zero pairing by this
single positive factor when necessary; they neither change the
absolute value nor introduce a power of $d$ for differentiated
slots.  No ramification--residue-degree identity is required.

% END source: sections/quotient_scope_checks.tex


% BEGIN canonical R5 source: sections/curve_testing_scope_obstruction.tex
\section{Why same-field curve tests do not suffice}
\label{sec:r5-curve-testing-scope-obstruction}

\begin{example}[A point invisible to every same-field analytic curve]
\label{ex:r5-same-field-curve-tests-false}
Let $K/\mathbf Q_p$ carry its given absolute value, choose a
uniformizer $\pi$, and put $r=|\pi|<1$.  In
$U=(\mathbf A_K^2)^{\rm an}$, let $x$ be the two-variable
Gauss point of polyradius $(r,r)$ and set
\[
                            f=\mathbf1_{\{x\}}.
\]
The function is bounded, Borel and upper semicontinuous.  Its
pullback along every analytic map from every $K$-analytic
curve to $U$ is identically zero.  Nevertheless $f$ is not
continuous and fails the Shilov bound on the unit bidisc.
Thus replacing all complete extensions by same-field curve
tests in Theorems~\ref{thm:r5-curvewise-shilov-without-usc}
and~\ref{thm:r5-continuity-from-curves} is false, even with
these extra boundedness and semicontinuity hypotheses.
\end{example}

\noindent\textbf{Status: \texttt{FALSE}} for the weakened
testing assertion; \texttt{DIRECT\_COROLLARY} for this explicit
construction from the dimension invariant of analytic points.

\begin{proof}
Write the affine coordinates as $T_1,T_2$.  The norm defining
$x$ is
\[
 \left|\sum_{i,j}c_{ij}T_1^iT_2^j\right|(x)
                   =\max_{i,j}|c_{ij}|r^{i+j}.
\]
The residue field of $\mathscr H(x)$ contains the two
algebraically independent residues of $T_1/\pi$ and
$T_2/\pi$.  In fact it is
$\widetilde K(\overline{T_1/\pi},\overline{T_2/\pi})$,
and the value group has no new rational rank.  Hence
$d_K(x)=2$.

Every point $c$ of a $K$-analytic curve has $d_K(c)\leq1$;
this is the pointwise dimension bound
\cite[\S1.4.6]{DucrosFamilies}.  An analytic map induces an
isometric embedding of completed residue fields from its image
point to its source point.  Such an embedding injects residue
fields and value groups, so it cannot increase $d_K$ at the
image.  Consequently no point of a $K$-analytic curve can
map to $x$.  Every indicated pullback of $f$ is therefore
the zero harmonic function.  This argument includes analytic
maps whose images are not contained in an algebraic curve.

The singleton $\{x\}$ is closed in the Hausdorff space $U$,
so its indicator is upper semicontinuous.  As $\rho\uparrow r$
with $\rho<r$, the Gauss points of polyradius $(\rho,r)$
converge to $x$ and are distinct from it.  This proves
noncontinuity.  The unit bidisc has the single Shilov point
$\zeta_{(1,1)}$, different from $x$, while $x$ lies in its
ordinary interior.  Therefore
\[
                  f(x)=1>0=\max_{\Gamma(\mathbb B_K^2)}|f|,
\]
which is the asserted failure of the Shilov bound.
\end{proof}

\paragraph{An explicit detecting extension and curve.}
The all-extension hypothesis detects this example directly.
Let $L=\mathscr H(x)$ and let $b\in L$ be the value of
$T_2$ at $x$.  On the smooth algebraic $L$-line consider
\[
                    h(s)=(s,b):\mathbf A_L^1\longrightarrow
                                              \mathbf A_L^2.
\]
For its radius-$r$ Gauss point $\zeta_{0,r}$ and every
polynomial $P=\sum_{i,j}c_{ij}T_1^iT_2^j\in K[T_1,T_2]$,
\[
 \begin{split}
 |P(h(\zeta_{0,r}))|
   &=\max_i\left|\sum_jc_{ij}b^j\right|r^i\\
   &=\max_{i,j}|c_{ij}|r^{i+j}=|P(x)|.
 \end{split}
\]
Thus its image in the original $U$ is exactly $x$.  For
$0<\rho<r$, the image of $\zeta_{0,\rho}$ has
$|T_1|=\rho$ and hence is not $x$.  The actual pullback
$f_L\circ h$ takes value one at $\zeta_{0,r}$ and zero
along these approaching points.  It is not continuous, and
therefore not harmonic, on the ordinary open disc $|s|<1$
which contains them all.  This is a calculation of actual
points, not a replacement of a fiber integration current by
a selected section.

The same obstruction remains if one tests over all finite
extensions of $K$: those extensions have $d_K=0$, and an
analytic curve over any one of them still has $d_K\leq1$
at each point.

\paragraph{The assigned spike also has zero ambient distribution.}
For every original compact smooth $(2,2)$-form $\eta$, its
signed integration measure has no atom at $x$.  Indeed, the
calibrated integration construction represents its total
variation locally by densities with respect to two-dimensional
Lebesgue measure on skeletal cells; see
\cite[\S\S5.7.1, 5.8.4--5.8.5 and 16.8.1]{CLD25}.
A point outside the carrying skeleton has zero mass, and a
point on it is Lebesgue-null on each local cell, including
their lower-dimensional faces.  Finite compact localization
therefore gives $|\mu_\eta|(\{x\})=0$.  Consequently
\[
                            [f]=0,\qquad D[f]=0.
\]
This argument concerns the measures defined by original smooth
top forms.  It does not assert that arbitrary positive currents
are non-atomic, or that $\{x\}$ is pluripolar or has zero
Bedford--Taylor capacity.

The example concerns the testing hypotheses
for arbitrary actual functions.  It is not a counterexample
to the quotient theorem for actual bounded classical
differences, whose proof establishes the requisite complete
extension functoriality before applying the curve criterion.

% END source: sections/curve_testing_scope_obstruction.tex


% BEGIN canonical R5 source: sections/weak_cone_lineality.tex
\section{The weak-positive current cone is not pointed}
\label{sec:r5-weak-cone-lineality}

We use the terminology of CLD exactly.  A symmetric current is
\emph{weakly positive} when it is nonnegative on strongly positive
compact tests; it is \emph{positive} when it is nonnegative on
positive compact tests; it is \emph{strongly positive} when it is
nonnegative on weakly positive compact tests.  These are
\cite[Definition 5.4.1]{CLD12} and
\cite[Definition 7.2.1]{CLD25}.  R, Convention 2.2, chooses the
middle, original positive cone.  The distinction is essential to
the two-sign conclusion used in the quotient theorem.

\begin{example}[A nonzero smooth closed current in the weak lineality space]
\label{ex:r5-weak-cone-lineality}
On $X=(\mathbf G_{m,K}^4)^{\rm an}$ there is a nonzero symmetric
original smooth $(2,2)$-form $\gamma$ such that its integration
current $S=[\gamma]$ is separately $d'$- and $d''$-closed and
both $S$ and $-S$ are weakly positive on \emph{all} original
tests.  Neither is positive in R's cone.  In particular
\[
       S\text{ and }-S\text{ weakly positive}
                          \quad\not\Longrightarrow\quad S=0.
\]
\end{example}

\noindent\textbf{Status: \texttt{EXTERNAL\_KNOWN}} for the
linear-algebra phenomenon, explicitly recorded in
\cite[Proposition 1.2.13 and the following paragraph]{CLD25};
\texttt{DIRECT\_COROLLARY} for the following signed calculation
and original-torus realization.  No new positivity convention is
introduced.

\begin{proof}
Write $t_i=\log|T_i|$, let $e_i=d't_i$, and for $i<j$ put
$e_{ij}=e_i\wedge e_j$.  The CLD involution sends $e_i$ to
$d''t_i$.  The positive $(4,4)$-volume is
\[
  \Omega=e_1\wedge e_2\wedge e_3\wedge e_4
                  \wedge Je_1\wedge Je_2\wedge Je_3\wedge Je_4.
\]
Indeed the sign $\sigma_4=(-1)^{4\cdot3/2}$ equals one.
For a two-form $a=\sum_{i<j}z_{ij}e_{ij}$, the positive
square is $\theta_a=-a\wedge Ja$, since $\sigma_2=-1$.
It is strongly positive when $a$ is decomposable.  Define
\[
\begin{split}
 \gamma=-\tfrac12\bigl(&e_{12}\wedge Je_{34}
                       +e_{34}\wedge Je_{12}
                       -e_{13}\wedge Je_{24}
                       -e_{24}\wedge Je_{13}\\
                     &+e_{14}\wedge Je_{23}
                       +e_{23}\wedge Je_{14}\bigr).
\end{split}
\]
This form is symmetric.  Direct exterior multiplication gives
\begin{equation}
       \gamma\wedge\theta_a
        =(z_{12}z_{34}-z_{13}z_{24}+z_{14}z_{23})\,\Omega.
 \label{eq:r5-plucker-square-pairing}
\end{equation}
For example, the complementary pairings
$e_{12}\wedge e_{34}$, $e_{13}\wedge e_{24}$,
$e_{14}\wedge e_{23}$ have signs $+,-,+$ against
$e_{1234}$.  Moving a two-form past a two-form has sign
$+1$, so the two minus signs in $\gamma$ and $\theta_a$
give exactly the displayed coefficient.  Equivalently, the
complement matrix $C$ on the ordered basis
$(12,13,14,23,24,34)$ interchanges the pairs
$(12,34),(13,24),(14,23)$ with signs $+,-,+$; it satisfies
$C^2=1$, and the coefficient is $\tfrac12 z^{\mathsf T}Cz$.

If $a$ is decomposable then $a\wedge a=0$, while in general
\[
 a\wedge a
       =2(z_{12}z_{34}-z_{13}z_{24}+z_{14}z_{23})e_{1234}.
\]
Thus $\gamma$ and $-\gamma$ pair to zero with every strongly
positive $(2,2)$-form on $\mathbf R^4$.  They are both
weakly positive there by the exact definition
\cite[Definition 5.1.1(b)]{CLD12}, equivalently
\cite[Definition 1.2.3(a)]{CLD25}.

Pull this constant form to the analytic torus by its original
moment coordinates.  Weak positivity is preserved by pullback
to every common refinement.  Its wedge with a strongly positive
test is weakly positive; the same holds with its negative.
On every top-dimensional calibrated cell of the four-dimensional
analytic space, the resulting top form is therefore both
nonnegative and nonpositive.  Its integral is zero.  A finite
localization of any original compact strongly positive test
proves
\[
                         \langle S,\eta\rangle=0.
\]
This verifies both weak positivity assertions for the entire
original test sheaf, not only for tests expressed in the four
displayed coordinates.  The pullback and wedge facts used here
are \cite[Lemma 1.2.11 and Proposition 1.2.12]{CLD25}.

To see nonvanishing, take
$a_+=e_{12}+e_{34}$ and $a_-=e_{12}-e_{34}$.
These are nondecomposable, and their squares $\theta_{a_+}$,
$\theta_{a_-}$ are positive in the middle cone.  Formula
\eqref{eq:r5-plucker-square-pairing} gives respectively
$\Omega$ and $-\Omega$.  Choose a nonzero nonnegative
$\rho\in C_c^\infty(\mathbf R^4)$.  The function
$\rho(t_1,\ldots,t_4)$ has compact support on $X$, since
the inverse image of a compact tropical box is a compact
polyannulus.  The original torus calibration, of degree one,
gives
\[
 \langle S,\rho\theta_{a_+}\rangle
                   =\int_{\mathbf R^4}\rho(t)\,dt>0,
 \qquad
 \langle S,\rho\theta_{a_-}\rangle
                   =-\int_{\mathbf R^4}\rho(t)\,dt<0.
\]
Hence $S\ne0$, $S$ is not positive, and $-S$ is not positive.
All coefficients of $\gamma$ are constant in unit-log
coordinates, so $d'\gamma=d''\gamma=0$.  Compact Stokes on
the boundaryless torus proves separate closedness of $S$.
It is a smooth-form current and therefore of coefficient-Radon
type; no general order-zero claim for weakly positive currents
is needed.
\end{proof}

\paragraph{Consequence for the quotient argument.}
The original positive cone used by R is pointed on symmetric
currents: positive forms span the full symmetric test space,
and two opposite positive currents consequently vanish
\cite[Remark 7.2.3]{CLD25}.  Strongly positive tests do not
span that space in the example above.  Thus replacing
``positive'' by ``weakly positive'' in the last two-sign step
would be invalid.  The construction is not a counterexample
to potential independence for weak backgrounds: it supplies
neither a nonzero quotient defect nor any pair of potentials
with such a defect.  It only disproves that proposed proof
step and fixes the positivity scope of the present result.

% END source: sections/weak_cone_lineality.tex


% BEGIN canonical R5 source: sections/weak_order_zero_tropical_obstruction.tex
\section{A closed weakly positive tropical current need not have order zero}
\label{sec:r5-weak-order-zero-tropical}

This section concerns the standard real tropical space $\mathbf R^5$,
with its ordinary smooth superform test sheaf.  No distribution is
pushed from a skeleton into a Berkovich space.  We use the convention
of \cite[\S 1.1.2]{CLD25}: $Jd't_i=d''t_i$,
$Jd''t_i=-d't_i$, and a symmetric form is fixed by $J$.
A symmetric current is weakly positive if it is nonnegative on all
strongly positive tests, exactly as in
\cite[Definition 7.2.1]{CLD25}.

\begin{example}[Separate closedness does not repair strong-test control]
\label{ex:r5-closed-weak-nonradon-tropical}
There is a symmetric current $T$ of bidegree $(3,3)$ on
$\mathbf R^5$ such that
\[
 d'T=d''T=0,\qquad T\ne0,\qquad
 T\text{ and }-T\text{ are weakly positive},
\]
but $T$ does not have order zero.  In particular it is not a
coefficient-Radon current.
\end{example}

\noindent\textbf{Status: \texttt{DIRECT\_COROLLARY}} for the
construction from the explicit Pl\"ucker form in
Example~\ref{ex:r5-weak-cone-lineality};
\texttt{FALSE} for the assertion that nonnegativity on all strongly
positive tests implies order zero, even with separate closedness.

\begin{proof}
Write $t=(t_0,t')$, where $t'=(t_1,t_2,t_3,t_4)$, and
$e_i=d't_i$.  Let $\gamma_4$ be the constant symmetric
$(2,2)$-form from Example~\ref{ex:r5-weak-cone-lineality}, in the
four variables $t'$.  Put
\[
 \Omega_4=e_{1234}\wedge Je_{1234},\qquad
 \gamma_5=e_0\wedge Je_0\wedge\gamma_4,\qquad
 \Omega_5=e_0\wedge Je_0\wedge\Omega_4.
\]
The last form is the positive standard volume on $\mathbf R^5$:
moving $Je_0$ past the four $e_i$ has sign $+1$, and
$\sigma_5=(-1)^{5\cdot4/2}=1$.
All three displayed forms are symmetric in the stated convention.
For $\eta\in A_c^{2,2}(\mathbf R^5)$ define the smooth scalar
coefficient $c_\eta$ by
\[
                         \gamma_5\wedge\eta=c_\eta\Omega_5
\]
and set
\begin{equation}
       \langle T,\eta\rangle
       =-\int_{\mathbf R^4}
                   \frac{\partial c_\eta}{\partial t_0}(0,t')\,dt'.
 \label{eq:r5-tropical-derivative-current}
\end{equation}
This is a current of order at most one, since each $c_\eta$ is a
fixed linear combination of the coefficients of $\eta$ and the
integral is over a fixed compact set when the test supports are
fixed.  Algebraically it is the constant form $\gamma_5$ with
coefficient distribution $\delta'_0(t_0)\otimes dt'$.

Because $J\gamma_5=\gamma_5$ and $J\Omega_5=\Omega_5$,
we have $c_{J\eta}=c_\eta$; hence $T(J\eta)=T(\eta)$.
Suppose that $\eta$ is strongly positive.  In
$\gamma_5\wedge\eta$, every term of $\eta$ involving $e_0$
or $Je_0$ disappears.  Its remaining part is the restriction
of $\eta$ to the four-dimensional slice with $t_0$ fixed.
Restriction preserves strong positivity, and $\gamma_4$
annihilates every strongly positive $(2,2)$-form on that slice.
Consequently $c_\eta$ is identically zero, not merely zero at
$t_0=0$.  Formula~\eqref{eq:r5-tropical-derivative-current}
therefore gives $T(\eta)=0$.  Both $T$ and $-T$ are weakly
positive.

For separate closedness, differentiation in any $t_i$ with $i>0$
of $\delta'_0\otimes dt'$ vanishes, by ordinary compactly
supported integration by parts.  Its only possibly nonzero
coefficient derivatives are in $t_0$.  In $d'T$ these acquire
the factor $e_0$, already present in $\gamma_5$, and in $d''T$
they acquire $Je_0$, also already present.  Both products are
zero.  This proves $d'T=d''T=0$ as distributional supercurrents;
the convention-dependent overall differentiation sign does not
affect either zero identity.

Finally, take
\[
 a_+=e_{12}+e_{34},\qquad \theta_+=-a_+\wedge Ja_+.
\]
The signed calculation in
\eqref{eq:r5-plucker-square-pairing} gives
$\gamma_4\wedge\theta_+=\Omega_4$, and therefore
$\gamma_5\wedge\theta_+=\Omega_5$.
Choose $\chi\in C_c^\infty(\mathbf R)$ identically one near zero
and nonzero $\rho\in C_c^\infty(\mathbf R^4)$ with $\rho\ge0$.
For positive integers $m$ put
\[
 \eta_m=m^{-1}\sin(m^2t_0)\chi(t_0)\rho(t')\theta_+.
\]
These tests have a common compact support, and all their
coefficients converge uniformly to zero.  Nevertheless,
\[
                  \langle T,\eta_m\rangle
                     =-m\int_{\mathbf R^4}\rho(t')\,dt'.
\]
Thus $T\ne0$ and is not continuous for the coefficient $C^0$
topology on a fixed compact support.  It does not have order
zero, and its coefficients cannot all be Radon measures.
\end{proof}

\begin{corollary}[Nonunique decreasing smooth multiplier limits]
\label{cor:r5-tropical-multiplier-limit-false}
For the fixed separately closed weakly positive current $T$ above,
the limit of $u_jT$ need not be independent of the smooth convex
decreasing approximation to a continuous convex function $u$,
even when the convergence of the functions is uniform.  Here all
products $u_jT$ are the ordinary products of a smooth function
and a distribution.
\end{corollary}

\noindent\textbf{Status: \texttt{DIRECT\_COROLLARY}};
\texttt{FALSE} for the proposed approximation-independent
continuous-multiplier construction under only weak positivity.

\begin{proof}
Let $u(t)=\max(0,t_0)$.  For $c\ge0$ and $\varepsilon>0$ set
\[
 u_{\varepsilon,c}(t)
       =\varepsilon\log\bigl(1+\exp(t_0/\varepsilon+c)\bigr).
\]
These are smooth nonnegative convex functions; their only
nonzero Hessian entry is nonnegative in the $t_0$ direction.
Moreover
\[
    0\le u_{\varepsilon,c}-u\le\varepsilon(c+\log2)
                    \quad\hbox{on all of }\mathbf R^5.
\]
For fixed $t_0$, put $x=t_0/\varepsilon$ and
$s=\exp(x+c)/(1+\exp(x+c))$.  Direct differentiation gives
\[
 \frac{\partial u_{\varepsilon,c}}{\partial\varepsilon}
  =-s\log s-(1-s)\log(1-s)+cs\ge0.
\]
Thus every sequence $\varepsilon_j\downarrow0$ gives smooth
convex majorants decreasing uniformly to the same $u$.

Let $T_0$ be defined by
$\langle T_0,\eta\rangle=\int c_\eta(0,t')\,dt'$.
It is the nonzero order-zero current with coefficient
$\delta_0\otimes dt'$ and form factor $\gamma_5$.
The distributional identity
$a\delta'_0=a(0)\delta'_0-a'(0)\delta_0$ gives
\[
 u_{\varepsilon,c}T
  =\varepsilon\log(1+e^c)\,T
                  -\frac{e^c}{1+e^c}\,T_0
  \longrightarrow -\frac{e^c}{1+e^c}\,T_0.
\]
For instance $c=0$ and $c=\log3$ give the distinct limits
$-T_0/2$ and $-3T_0/4$.  This proves the multiplier assertion.
Every $u_{\varepsilon,c}T$ is separately closed: the only
derivative of the multiplier has the already present factor
$e_0$ or $Je_0$.  In particular
\[
                         D(u_{\varepsilon,c}T)=0
                         \quad\hbox{for every }\varepsilon,c.
\]
The example therefore does not give a nonunique BT curvature or
a quotient defect.  It isolates the failure of a multiplier
limit on a non-Radon weakly positive background, solely in the
standard real tropical calculus.
\end{proof}

\paragraph{Precise comparison with the printed statements.}
As written, \cite[Proposition 7.2.9, p.~162]{CLD25} assumes
symmetry and nonnegativity on strongly positive tests, and
concludes order zero on a locally differentially separated
tropical space.  The standard space $\mathbf R^5$ has that
separation property, so Example~\ref{ex:r5-closed-weak-nonradon-tropical}
contradicts the order-zero conclusion under its printed
hypotheses.  The displayed domination step in that proof would
require an arbitrary symmetric form $\alpha$ to satisfy
\[
             \alpha+C\|\alpha\|\beta^p
                          \quad\hbox{strongly positive}
\]
for a fixed positive definite $(1,1)$-form $\beta$.
The relative-interior statement in
\cite[Proposition 1.2.13, pp.~27--28]{CLD25} does not give this:
strongly positive forms span only a proper subspace in the
middle degrees used here.  The same proposition explicitly
records this proper-span phenomenon.

There is also a distinct printed strong-span assertion in
\cite[Lemma 5.3.3, p.~81]{CLD12}: every smooth $(p,p)$-form
with strongly paracompact support is asserted to be a difference
of strongly positive forms.  Example~\ref{ex:r5-weak-cone-lineality}
disproves that assertion even for symmetric compactly supported
\emph{original Berkovich} forms.  Indeed, on the four-dimensional
analytic torus take $\omega=\rho(t)\theta_{a_+}$ as in that
example.  Its compact support meets the stated support hypothesis,
and $\langle[\gamma_4],\omega\rangle>0$.
If $\omega=\omega_1-\omega_2$ with both $\omega_i$ strongly
positive, choose an original nonnegative smooth compact cutoff
$\xi$ equal to one near $\operatorname{supp}\omega$.
Then $\omega=\xi\omega_1-\xi\omega_2$, whereas
$[\gamma_4]$ annihilates both compact strongly positive terms.
This is a contradiction.  The replacement by the middle
positive span is the one stated in
\cite[Proposition 6.5.2]{CLD25}.  This analytic strong-span
counterexample uses only the smooth original torus current;
it does not transport the derivative distribution constructed
above.

The earlier \cite[Lemma 5.4.4, p.~82]{CLD12} has a different
hypothesis: its linear functional is nonnegative on every
\emph{weakly positive} test.  Our current does not satisfy that
stronger test requirement, so this construction does not refute
that lemma.  Explicitly, choose the preceding $\chi$ nonnegative
with support in $(-1/2,1/2)$.  The test
$(1+t_0)\chi(t_0)\rho(t')\theta_+$ is middle positive, hence
weakly positive, whereas its pairing with $T$ is
$-\int\rho<0$.  In particular one must not silently identify
the two printed hypotheses.

For the standard real test space, nonnegativity on the
\emph{middle positive} cone is already a sufficient replacement
for the order-zero conclusion.  Indeed $\beta^p$ lies in the
full interior of that cone.  On a fixed compact support,
finite-dimensional cone domination bounds a symmetric test
above and below by $C\|\eta\|_{C^0}\zeta\beta^p$, with a fixed
nonnegative smooth compact cutoff $\zeta$ equal to one on the
support.  Positivity then gives a coefficient $C^0$ bound;
symmetry annihilates the antisymmetric part.  This is precisely
the distinction relevant to R, Convention 2.2, which chooses
the middle cone of \cite[Definition 5.1.1(c)]{CLD12}.  Neither
this example nor its proof changes that internal input.

\paragraph{Scope of the obstruction.}
The derivative distribution above is defined on an ordinary real
tropical space.  It has not been transported to a Berkovich
skeleton.  Such a transport would require control of all
original refined tests, whose skeletal restrictions need not be
smooth in these fixed coordinates.  The original analytic
counterexample in Example~\ref{ex:r5-weak-cone-lineality} remains
a smooth, order-zero closed current.  The present stronger
non-Radon statement is confined to the real tropical setting,
and neither example asserts that the quotient theorem is false
for every possible weak-positive background class.

% END source: sections/weak_order_zero_tropical_obstruction.tex


% BEGIN canonical R5 source: sections/potential_quotient_energy.tex
\section{The degree-one energy reduction and its precise obstruction}
\label{sec:r5-quotient-energy}

Fix the original finite extension $K/\mathbf Q_p$ and its absolute
value. All forms, currents and positivity cones in this section are
the original CLD objects. Write $D=d'd''$. An allowed space and
the class $\mathcal C_b$ have their frozen meanings. In particular
the functions in $\mathcal C_b$ have specified actual values.

We study the implication
\begin{equation}\label{eq:r5-quotient-target}
 D[u]=D[v]\quad\Longrightarrow\quad D((u-v)S)=0
\end{equation}
for every fixed symmetric original-positive, separately closed
background $S$. The full implication is not proved here.

\subsection{An exact reduction to remaining degree one}

\begin{proposition}[Degree-one sufficiency for the arbitrary-background quotient]
\label{prop:r5-degree-one-quotient-reduction}
Suppose that, on every allowed ordinary open set of pure dimension
$n$, implication \eqref{eq:r5-quotient-target} holds for every
original-positive separately closed background of bidegree
$(n-1,n-1)$. Then it holds for every such background of bidegree
$(p,p)$, $0\leq p<n$, on the same class of spaces and over the
same original field.
\end{proposition}

\noindent\textbf{Status: \texttt{DIRECT\_COROLLARY}} of the original
positive-cone algebra and localization. This reduces the remaining
problem; it is not a proof of its degree-one hypothesis.

\begin{proof}
Put $q=n-p$ and $H=u-v$. The current $D(HS)$ is symmetric.
It therefore suffices to test its symmetric complementary forms.
Use a finite original smooth partition subordinate to moment
presentations of an arbitrary compact test. On one such
presentation let $t_1,\ldots,t_m$ be the real moment coordinates.
Every symmetric $(q-1,q-1)$-test is a finite sum
\begin{equation}\label{eq:r5-positive-square-test-span}
 \sum_\ell f_\ell\theta_\ell,
 \qquad f_\ell\in A_c^{0,0},
\end{equation}
where each $\theta_\ell$ is a constant positive
$(q-1,q-1)$-form in the $d't_i,d''t_i$. Indeed the symmetric
coefficient matrices are spanned by positive rank-one matrices:
the diagonal squares and the polarization squares of the finite
basis $d't_I$ give an explicit finite spanning family. Here a
rank-one form is the positive signed square of an arbitrary
constant $(q-1,0)$-form. That form need not be decomposable.
Thus this argument uses the full original positive cone, not
only the strongly positive decomposable squares.

Every $\theta_\ell$ is separately closed. The current
$Q_\ell=S\wedge\theta_\ell$, defined on this ordinary moment
neighborhood, has bidegree $(n-1,n-1)$ and is separately closed.
It is original-positive: wedging a positive $(1,1)$-test with
the positive square $\theta_\ell$ gives a positive test for $S$.
This last assertion follows by expanding that $(1,1)$-test as
a sum of positive covector squares; each product is the positive
square of the exterior product of the two covectors/forms.
The assertion survives pullback to a common moment refinement.

The assumed degree-one statement gives $D(HQ_\ell)=0$.
The smooth closed product rule, valid also for the coefficient
Radon current $HS$, gives
\[
 \langle D(HS),f_\ell\theta_\ell\rangle
       =\langle D(HQ_\ell),f_\ell\rangle=0.
\]
Sum \eqref{eq:r5-positive-square-test-span} and the finite
partition. Antisymmetric tests vanish by symmetry, so every
original compact complementary test has zero pairing.
\end{proof}

\begin{proposition}[The exact form-theoretic meaning of the degree-one gate]
\label{prop:r5-quotient-form-gate}
Let $S$ have bidegree $(n-1,n-1)$ on a norm box $\Omega$ to
which the faithful form construction applies. If
$H=u-v$ with $u,v\in\mathcal C_b(\Omega)$, then
$H\in\mathscr D_{S,\mathrm{loc}}$. For every original compact
smooth scalar $\phi$ and every original smooth compact cutoff
$\chi$ equal to one near $\operatorname{supp}\phi$,
\begin{equation}\label{eq:r5-defect-energy-duality}
 \langle D(HS),\phi\rangle
       =-\mathcal E_S(\chi H,\phi).
\end{equation}
In particular, \eqref{eq:r5-quotient-target} in remaining degree
one is precisely the assertion that the given actual difference
is locally harmonic for every such faithful closed form.
\end{proposition}

\noindent\textbf{Status: \texttt{DIRECT\_COROLLARY}} of
Proposition~\textup{R4--86.1}.

\begin{proof}
Apply that proposition to $u$ and $v$ and subtract. Its strong
local form convergence and Green identity prove
\eqref{eq:r5-defect-energy-duality}. The cutoff disappears from
the integral because it equals one near the support of the
test and its derivatives. Original compact scalar tests are
exactly the tests for a top-degree current, proving the last
assertion. Neither $H$ nor $-H$ is declared psh in this argument.
\end{proof}

The ambient equality $D[H]=0$ only tests $H$ against $D\eta$
for original smooth complementary forms $\eta$. The right side
of \eqref{eq:r5-defect-energy-duality} instead pairs it with the
coefficient measure of $D(\phi S)$. That measure-valued test
has not been replaced by an original smooth test. The faithful
measure's domination of fixed smooth coefficients does not make
this replacement: in particular it does not assert
$D(uS)\ll\lambda_S$.

\subsection{A failure of the unrestricted energy-completion route}

\begin{example}[Uniform bounded psh decompositions and universal strong energy do not suffice]
\label{ex:r5-universal-energy-polar-obstruction}
Choose a uniformizer $\pi$ of $K$, put $R=|\pi|^{-1}$ and
$\ell=\log R>0$, and work on the ordinary open bidisc
\[
 \Omega=\{|z|<R,\ |t|<R\}\subset(\mathbf A_K^2)^{\rm an}.
\]
Let $Y=\{z=0\}\cap\Omega$ and set
\[
 q(t)=\ell^{-1}\max\{0,\log|t|\},\qquad
 b_N(z)=\max\{0,1+N^{-1}\log(|z|/R)\}.
\]
The values at zeros are given by the constant branches. Thus
$0\leq q<1$, $0\leq b_N<1$, and $b_N$ increases pointwise to
$\mathbf1_{Y^c}$. Define
\begin{align*}
 u_N&=q+\tfrac12(q^2+b_N^2),&
 v_N&=\tfrac12(q+b_N)^2,\\
 f_N&=u_N-v_N=q(1-b_N),&
 f&=\mathbf1_Yq.
\end{align*}
Then the following assertions hold.
\begin{enumerate}
\item The two rows $u_N,v_N$ consist of actual continuous
classical psh functions on the whole bidisc, increase pointwise,
and lie between zero and two. The functions $f_N$ are continuous
bounded differences with $0\leq f_N\leq1$ and $f_N\downarrow f$.
\item On all original compact tests, $[f_N]\to0$ and
$D[f_N]\to0$ as ambient currents.
\item For the fixed background $S_0=[Y]$,
\[
 D(f_NS_0)=D(qS_0)
      =\ell^{-1}\delta_{(0,\zeta_{0,1})}\ne0
       \qquad\hbox{for every }N.
\]
Here $\zeta_{0,1}$ is the unit Gauss point of the $t$-disc,
and the measure has the original $K$ calibration.
\item For every fixed original-positive separately closed
background $S$ of bidegree $(1,1)$ and every original smooth
compact cutoff $\chi$, the vectors $\chi f_N$ converge strongly
in the faithful closed form of $S$ to the vector represented by
$\chi\mathbf1_Yq$. If $\mathbf1_YS=0$, this limit is zero.
For $S=S_0$ the relative local energy need not tend to zero.
\item The limit $f$ is not a difference of two actual locally
bounded classical psh functions.
\end{enumerate}
\end{example}

\noindent\textbf{Status: \texttt{DIRECT\_COROLLARY}} for this
sharpened scope test, from the frozen model formula, actual
faithfulness, and complete-pole energy splitting. It proves
\texttt{FALSE} an unrestricted passage from weak ambient
curvature vanishing to arbitrary-background curvature vanishing,
even with the displayed uniform decompositions and strong energy
convergence. It does not refute
\eqref{eq:r5-quotient-target} for a fixed actual difference.

\begin{proof}
The functions $q$ and $b_N$ are continuous maxima of constants
and positive scalar multiples of logarithms of analytic functions.
Their curve pullbacks are subharmonic over every complete extension,
so they are classical. No rationality of the real factor $\ell^{-1}$
or of the additive constants is required.
The functions $s\mapsto s^2$ on $[0,\infty)$ are convex and
nondecreasing. Their compositions with the displayed nonnegative
classical potentials are classical. This proves all psh,
continuity and monotonicity assertions. The numerical bounds
follow from $0\leq q,b_N\leq1$, and the stated identity for
$f_N$ is an algebraic expansion.

The proper analytic locus $Y$ has zero mass for every original
smooth top-degree coefficient measure on the bidisc. Consequently
$[f]=0$. Fixed-coefficient dominated convergence, using the
common bound of $f_N$, proves $[f_N]\to0$. Distributional
differentiation proves $D[f_N]\to0$; equivalently apply the
same bounded convergence to the fixed top form $D\eta$ for
each original compact complementary test $\eta$.

At every actual point of $Y$ one has $b_N=0$, so
$f_NS_0=qS_0$ as coefficient-Radon currents for every $N$.
The one-variable original model formula
$D\max(0,\log|t|)=\delta_{\zeta_{0,1}}$, followed by the
closed immersion, proves assertion 3.

For assertion 4 use the complete pole witness
$\log(|z|/R)$ and the splitting of
Theorem~\textup{R4--86.2}. In particular,
for every fixed compact nonnegative smooth cutoff $\rho$,
\begin{equation}\label{eq:r5-pole-selector-energy}
 \int\rho\,\Gamma_S(b_N)\longrightarrow0.
\end{equation}
Put $g=\chi q$, which is a bounded compactly supported form
vector by Proposition~\textup{R4--86.1} and
the product rule. The measure-valued product square inequality
gives, on a compact neighborhood of its support,
\[
 \mathcal E_S\bigl(g(b_N-b_M)\bigr)
 \leq
 2\int(b_N-b_M)^2\Gamma_S(g)
 +2\|g\|_\infty^2\int\rho\,\Gamma_S(b_N-b_M).
\]
In the second term choose $\rho=1$ on that neighborhood and
use $\Gamma_S(b_N-b_M)\leq
2\Gamma_S(b_N)+2\Gamma_S(b_M)$. Equation
\eqref{eq:r5-pole-selector-energy} makes it tend to zero.
The first term tends to zero by bounded convergence against
the fixed finite measure $\Gamma_S(g)$.
The $L^2(\lambda_S)$ limit is $g\mathbf1_{Y^c}$, so closedness
identifies the strong form limit. Subtract from $g$ to obtain
$\chi f_N\to\chi\mathbf1_Yq$. This proof applies separately
to every fixed $S$, with its own faithful form and controls.

If $\mathbf1_YS=0$, faithfulness of $\lambda_S$ makes the
limit vector zero. In contrast $f_N=q$ on $Y$, so the local
energy for $[Y]$ is constant in $N$. A nonzero cutoff supported
on a radial segment $1<|t|<R$ has positive energy pairing with
$\Gamma_{[Y]}(q)$; the slope there is $1/\ell$. Thus this
relative energy does not vanish.

Finally, if $f=a-b$ for $a,b\in\mathcal C_b(\Omega)$, the
ambient identity $[f]=0$ would imply $[a]=[b]$. Frozen actual
current-faithfulness would give $a=b$ everywhere, contrary to
$f=q>0$ on the part of $Y$ with $|t|>1$. This proves assertion 5.
\end{proof}

\paragraph{The excluded estimate.}
The example rules out any general estimate whose hypotheses
are only one uniform bound on the two classical decompositions,
weak vanishing of their ambient curvature difference, and strong
convergence in all fixed relative energy spaces, and whose
conclusion is vanishing of the relative curvature differences.
It also rules out obtaining the desired quotient by simply
replacing actual classical differences with the full linear
local energy domains. The function $\mathbf1_Yq$ is an energy
vector for every such background and annihilates all ambient
smooth current tests, but it does not annihilate
$D(\phi[Y])$. Therefore an approximation of those latter
functionals in the dual of the unrestricted linear energy
space cannot hold in general. Any successful energy argument
must preserve an additional property of the actual classical
differences; normal trace or a genuine regularity conclusion
cannot be discarded in completion.

\begin{remark}[Bounded upper semicontinuity does not repair the scope]
\label{rem:r5-bounded-usc-kernel-obstruction}
\textbf{Status: \texttt{DIRECT\_COROLLARY}.}
The limit $f=\mathbf1_Yq$ in the example is bounded and upper
semicontinuous. Indeed, outside the closed set $Y$ it is locally
zero. At a point of $Y$, the inequalities $0\leq f\leq q$ and
continuity of $q$ give
$\limsup_{x\to y}f(x)\leq q(y)=f(y)$.
Nevertheless
\[
                    [f]=0,\qquad D[f]=0,
             \qquad D(f[Y])=\ell^{-1}\delta_{(0,\zeta_{0,1})}
                                  \ne0.
\]
Thus the actual bounded-classical-difference hypothesis in the
kernel theorem cannot be replaced even by arbitrary bounded upper
semicontinuous functions with zero original curvature. In
particular, merely calling such a representative original
CLD-positive does not place it in $\mathcal C_b$.
This is not a counterexample to the quotient for actual
$\mathcal C_b-\mathcal C_b$ differences. Nor does it disprove
existence of a continuous representative for a zero distribution:
the distribution $[f]=0$ here already has the continuous
representative zero. The failure is precisely equality with the
chosen upper semicontinuous values on the analytic locus.
\end{remark}

\begin{remark}[A quantitative fixed-test ambient error]
\label{rem:r5-energy-fixed-test-rate}
\textbf{Status: \texttt{DIRECT\_COROLLARY}.}
For every fixed original compact smooth top form $\eta$ on
$\Omega$, there is a constant $C_\eta<\infty$ such that
\[
             |\langle[f_N],\eta\rangle|\leq C_\eta/N.
\]
Indeed the fixed original top-form integration is carried by a
compact set $K_\eta\subset\Omega\setminus Y$. This is the
finite moment-presentation support property used in the proof of
\cite[Lemma 4.6.1, pp.~72--73]{CLD12}: the calibrated
top-dimensional supports avoid every proper analytic locus.
Let $\mu_\eta$ be its signed coefficient measure and set
\[
 A_\eta=\sup_{K_\eta}\log(R/|z|)<\infty.
\]
The pointwise formula
\[
 0\leq f_N
   =q\min\{1,N^{-1}\log(R/|z|)\}
   \leq A_\eta/N\qquad\hbox{on }K_\eta
\]
proves the claim with $C_\eta=A_\eta|\mu_\eta|(K_\eta)$.
For every fixed original compact $(1,1)$-test $\alpha$,
apply the same statement to $D\alpha$ to obtain
\[
       |\langle D[f_N],\alpha\rangle|\leq C_{D\alpha}/N.
\]
These constants depend on the fixed tests. They are not uniform
over an unbounded class of tests or over tests approaching $Y$.
The nonzero relative defect in the example is unchanged.
\end{remark}

\begin{remark}[Failure on relative capacitary completions]
\label{rem:r5-completed-potential-quotient-false}
\textbf{Status: \texttt{DIRECT\_COROLLARY}; the proposed
completed-class extension of (Q) is \texttt{FALSE}.}
The two increasing classical rows in the example have actual limits
\[
 a=q+\tfrac12(q^2+\mathbf1_{Y^c}),\qquad
 b=\tfrac12(q+\mathbf1_{Y^c})^2.
\]
They are uniformly bounded by two on the entire bidisc. Thus
Corollary~\textup{R4--18.5} places both limits
in $\mathcal A_{S,b}(\Omega)$ for every
permitted fixed $S$ of positive remaining degree. The convergence
is the all-domain local outer-capacity convergence of that
corollary; no ambient upper regularization is substituted for
the displayed limits.

In particular take $S_0=[Y]$. The completed first-order recursion
of Theorem~\textup{R4--17.6} is still
$\mathcal T_{S_0,1}(a)=D(aS_0)$, with the actual Borel multiplier,
and similarly for $b$. Since $a-b=\mathbf1_Yq$, one has
\[
 [a]=[b],\qquad D[a]=D[b],\qquad
 \mathcal T_{S_0,1}(a)-\mathcal T_{S_0,1}(b)
      =\ell^{-1}\delta_{(0,\zeta_{0,1})}\ne0.
\]
Thus even a common global defining bound and increasing classical
defining rows do not extend the ambient-curvature quotient to
these relative completed classes. This does not contradict the
completed calculus's invariance under equality outside an
$S_0$-capacity-zero set: the two displayed functions are not
equal in that relative sense. It also leaves the quotient for
actual $\mathcal C_b$ potentials untouched.
\end{remark}

% END source: sections/potential_quotient_energy.tex


% BEGIN canonical R5 source: sections/universal_borel_kernel.tex
\section{The universal relative Borel kernel is larger than the classical kernel}
\label{sec:r5-universal-borel-kernel}

The quotient theorem concerns actual locally bounded differences of
classical psh functions.  Even vanishing against every allowed positive
closed background does not impose the same regularity on arbitrary
Borel functions.

Let $U$ be an allowed pure $n$-dimensional ordinary open, with
$n\geq1$.  Define
\[
 \mathscr K_{\mathrm{Bor}}(U)=
 \left\{H:U\to\mathbf R:\begin{array}{l}
 H\text{ is locally bounded and Borel, and }D(HS)=0\\
 \text{for every symmetric original-positive, separately}\\
 d'\text{- and }d''\text{-closed }(p,p)\text{-current }S,\quad
 0\leq p<n
 \end{array}\right\}.
\]
The product $HS$ means actual Borel multiplication of the
coefficient-Radon current.  The asserted distributional equality is on
all original compact smooth complementary forms.  No positive sign is
required of $H$ or of the product $HS$.

\begin{proposition}[Finite complete-pole selectors]
\label{prop:r5-universal-pole-selectors}
Let $P_1,\ldots,P_m$ be pairwise disjoint Borel subsets of $U$.
Assume each $P_a$ is locally the exact pole set of an extended
classical psh function.  Put
$P_0=U\setminus\bigcup_{a=1}^mP_a$.
If $h_0,\ldots,h_m\in\mathscr{PH}(U)$, then
\[
                 H=\sum_{a=0}^m\mathbf1_{P_a}h_a
                       \quad\hbox{belongs to }\mathscr K_{\mathrm{Bor}}(U).
\]
In particular both $\mathbf1_P$ and $\mathbf1_{P^c}$ belong to
this kernel for every locally complete classical pole set $P$.
\end{proposition}

\noindent\textbf{Status: \texttt{DIRECT\_COROLLARY}} of the
already proved complete-pole restriction theorem and first-order
positivity over a fixed background.

\begin{proof}
Fix one of the currents $S$ in the definition.  The R2 restriction
theorem, Theorem~\textup{R4--24.3}, applied locally to
the exact pole witness and then glued, makes
\[
                        S_a=\mathbf1_{P_a}S\qquad(1\leq a\leq m)
\]
symmetric, original-positive and separately closed.  The complement
\[
                     S_0=S-\sum_{a=1}^mS_a=\mathbf1_{P_0}S
\]
is separately closed by finite linearity and is positive because the
sets are disjoint.  This argument is a restriction of a current
already defined on $U$.  It does not assert an extension theorem
for a current initially defined off a pole set.

For each $a$, both $h_a$ and $-h_a$ are continuous locally bounded
classical psh functions.  First-order positivity over the fixed
positive closed background $S_a$ therefore makes both
$D(h_aS_a)$ and its negative positive.  The original positive cone
is pointed, so $D(h_aS_a)=0$.  Actual Borel multiplication and
finite linearity give
\[
                   D(HS)=\sum_{a=0}^mD(h_aS_a)=0.
\]
The finite sum is locally bounded and Borel.  This proves the
claim for every $S$ and every original test.
\end{proof}

\begin{corollary}[Strict enlargement and its precise boundary]
\label{cor:r5-universal-borel-kernel-strict}
On a bidisc with coordinates $(z,t)$ in $(\mathbf A_K^2)^{\rm an}$, with
$Y=\{z=0\}$, the function $\mathbf1_Y$ is bounded and upper
semicontinuous, and
\[
       [\mathbf1_Y]=0,\qquad
       D(\mathbf1_Y S)=0\quad\hbox{for every allowed }S.
\]
It is not continuous and hence is not pluriharmonic.  Thus
$\mathscr{PH}(U)\subset\mathscr K_{\mathrm{Bor}}(U)$ can be
a strict inclusion.  On every allowed $U$, however,
\[
 \mathscr K_{\mathrm{Bor}}(U)\cap\mathscr{DSH}_{\mathrm{lb}}(U)
                       =\mathscr{PH}(U).
\]
\end{corollary}

\begin{proof}
The divisor $Y$ is the exact pole set of $\log|z|$, so the
proposition applies.  The indicator of a closed set is upper
semicontinuous.  It is discontinuous, for example at the rigid
origin, which is approached by the points $(\pi^j,0)$ off $Y$.
Original smooth top forms have support disjoint from the
lower-dimensional analytic locus $Y$
\cite[Proposition 15.2.1]{CLD25}; consequently
$[\mathbf1_Y]=0$.  Its noncontinuity and
Theorem~\ref{thm:r5-bounded-kernel-regularity} therefore
exclude $\mathbf1_Y$ from $\mathscr{DSH}_{\mathrm{lb}}$.

Every pluriharmonic function lies in the universal kernel by
the two signs of first-order positivity.  Conversely, if
$H\in\mathscr K_{\mathrm{Bor}}(U)\cap
\mathscr{DSH}_{\mathrm{lb}}(U)$, the permitted choice $S=[1]$
gives $D[H]=0$.  Theorem~\ref{thm:r5-bounded-kernel-regularity}
then makes $H$ pluriharmonic, proving the equality.
\end{proof}

\paragraph{A necessary distinction.}
Being carried by a complete pole set does not by itself imply
membership in the universal kernel.  With the notation of
Example~\ref{ex:r5-universal-energy-polar-obstruction},
$f=\mathbf1_Yq$ still satisfies $[f]=0$, but
\[
                  D(f[Y])=D(q[Y])
                         =\ell^{-1}\delta_{(0,\zeta_{0,1})}\ne0.
\]
The finite selector proposition requires the displayed ambient
pluriharmonic weights; it does not replace them by arbitrary
bounded Borel or classical psh weights.  Nor does it assert a
classification of the whole Borel kernel, countable selector
closure, or invariance under arbitrary changes on ambient null
sets.  In particular the discontinuous selector does not
contradict existence of a continuous representative of the zero
distribution: its ambient current already has the representative
zero.  It shows that universal relative curvature does not
recover that continuous representative from arbitrary assigned
actual Borel values.

% END source: sections/universal_borel_kernel.tex


% BEGIN canonical R5 source: sections/weak_current_quotient_obstruction.tex
\section{The descended operation is not weakly continuous}
\label{sec:r5-weak-current-quotient-obstruction}

Potential independence is a statement about exact equality of
curvature currents. It gives no continuity assertion for the
operation when those currents vary weakly. The following example
keeps the background fixed and makes both the approximating
curvatures and their common ambient limit original smooth.

\begin{example}[Smooth weak convergence with a fixed relative defect]
\label{ex:r5-smooth-weak-current-defect}
Let $K/\mathbf Q_p$ and its absolute value be the prescribed ones.
Fix a uniformizer $\pi$, put $R=|\pi|^{-1}$ and
$\ell=\log R$, and let
\[
 \Omega=\{|z|<R,\ |t|<R\}\subset(\mathbf A_K^2)^{\rm an},
 \qquad Y=\{z=0\}\cap\Omega,\qquad S=[Y].
\]
There exist original smooth classical psh functions
$v_N,w$ on $\Omega$, all between zero and two, such that
\[
                  D[v_N]\longrightarrow D[w]
\]
on every original compact test, whereas
\[
        D(v_NS)\not\longrightarrow D(wS).
\]
More precisely, $D(v_NS)$ is independent of $N$, and
$D(wS)-D(v_NS)$ is a fixed nonzero positive current. The ambient
error is $O(N^{-1})$ on each fixed original test, with a
test-dependent constant.
\end{example}

\noindent\textbf{Status: \texttt{FALSE}} for weak-current
continuity of the descended operation; the construction and
proof are \texttt{DIRECT\_COROLLARY} of smooth chain rules,
fixed-coefficient dominated convergence, and the original
top-form support property. This is not a counterexample to
exact potential independence.

\begin{proof}
Choose a nonnegative even $\kappa\in C_c^\infty((-1,1))$
with integral one and not identically zero. Put
\[
       \psi(s)=\int_{-\infty}^s(s-a)\kappa(a)\,da.
\]
Then $\psi$ is smooth and convex, $0\leq\psi'\leq1$,
$\psi''=\kappa$, it is zero on $(-\infty,-1]$, and
it equals $s$ on $[1,\infty)$. Its derivative bounds give
\begin{equation}
 0\leq 1-\psi(1-A/N)\leq A/N\qquad(A\geq0).
 \label{eq:r5-smooth-selector-rate}
\end{equation}

Define
\[
 q(t)=\psi\!\left(\frac{\log|t|}{\ell}\right),
 \qquad
 b_N(z)=\psi\!\left(1+\frac{\log(|z|/R)}{N}\right),
\]
with value zero at the corresponding coordinate zero.
These are original smooth functions: the profiles are
identically zero on an ordinary neighborhood of each zero
locus, and elsewhere they are smooth functions of unit
logarithms. They are classical psh after every complete
extension, by the curve chain rule for convex nondecreasing
profiles of log norms. On the open bidisc,
\[
           0\leq q,b_N\leq1,
       \qquad b_N\uparrow\mathbf1_{Y^c}.
\]
Now set
\[
             v_N=\tfrac12(q+b_N)^2,
          \qquad w=\tfrac12(q+1)^2.
\]
The square is convex and nondecreasing on the nonnegative
range. Thus $v_N,w$ are original smooth classical psh and
are bounded by two. Their original curvatures are positive
smooth forms, by the smooth chain rule. The row $v_N$
increases, but its actual limit on $Y$ is $q^2/2$, not
$w|_Y$. No decreasing-row convergence theorem applies to
this non-upper-semicontinuous actual limit.

Every fixed original compact smooth top form $\eta$ has
compact actual support disjoint from $Y$, by
\cite[Proposition 15.2.1]{CLD25}. On that support the number
\[
                 A_\eta=\sup\log(R/|z|)
\]
is finite. The displayed formulas and
\eqref{eq:r5-smooth-selector-rate} give
\[
 \begin{aligned}
 0\leq w-v_N
  &=\tfrac12(1-b_N)(2q+1+b_N)\\
  &\leq2(1-b_N)\leq 2A_\eta/N.
 \end{aligned}
\]
Integration against the fixed signed measure of $\eta$
proves
\[
 |\langle[v_N]-[w],\eta\rangle|
       \leq 2A_\eta|\mu_\eta|(\Omega)/N.
\]
For a fixed original compact $(1,1)$-test $\alpha$, apply
this estimate to $D\alpha$. This proves the stated ambient
curvature convergence and its fixed-test rate.

At every actual point of $Y$, $b_N=0$. Consequently
\[
  D(v_NS)=D\!\left(\tfrac12q^2[Y]\right)
                      \quad\text{for every }N,
\]
while separate closedness of $[Y]$ gives
\[
 D(wS)=D\!\left(\tfrac12q^2[Y]\right)+D(q[Y]).
\]
The last summand is positive and nonzero. Indeed, on the
$t$-annulus in $Y$, with $s=\log|t|/\ell$,
\[
                     Dq=\kappa(s)d's\wedge d''s.
\]
For any nonnegative smooth compact scalar profile $\rho$
with $\int\rho(s)\kappa(s)\,ds>0$, the original annulus
calibration gives
\[
       \int_Y\rho(s)Dq
             =\ell^{-1}\int\rho(s)\kappa(s)\,ds>0.
\]
Multiply the scalar by an original compact cutoff in the
$z$-direction equal to one near $z=0$ to obtain a genuine
compact test on $\Omega$. Hence $D(q[Y])\ne0$, proving
the persistent relative defect.
\end{proof}

\paragraph{Consequences for the current-valued formulation.}
Let $\alpha_N=D[v_N]$ and $\alpha=D[w]$. By
Corollary~\ref{cor:r5-local-potential-current-products},
$\alpha_N\wedge[Y]$ and $\alpha\wedge[Y]$ are independent
of their bounded classical potentials. The example therefore
proves an intrinsic failure of weak continuity on this class
of positive currents, not an ambiguity caused by the choice
of potentials. In fact all the displayed currents already
come from smooth positive forms. The failure occurs when a
fixed singular background detects the unchanged trace on $Y$.

The example extends to every ambient dimension $n\geq2$
by taking a product with an open polydisc of dimension $n-2$
and pulling back the displayed functions. The same
hyperplane integration current remains the background; a
compact positive top form in the extra directions detects
the unchanged defect. No claim about arbitrary varying
backgrounds is needed.


% END source: sections/weak_current_quotient_obstruction.tex


% BEGIN canonical R5 source: sections/compact_capacity_finite_curvature.tex
\section{A finite-curvature sufficient condition for compact no-gain}
\label{sec:r5-compact-capacity-finite-curvature}

This section isolates one exact consequence of the new potential
quotient for the unresolved compact-capacity gate C1.  It does not
claim that an arbitrary normalized family satisfies the additional
finite-dimensional hypothesis, or that approximation by such
families suffices.

Fix the original finite extension $K/\mathbf Q_p$, an allowed
ordinary open $\Omega$, and a symmetric original-positive
separately closed $(p,p)$-current $S$.  Put $q=n-p\ge1$, and
write
\[
 \mathcal P(\Omega)
   =\{u\in\mathcal C_b(\Omega):0\le u\le1\text{ everywhere on }\Omega\}.
\]
The diagonal raw capacity is
\[
 C_S(E;\Omega)=\sup_{u\in\mathcal P(\Omega)}
                  T_{S,q}(u,\ldots,u)(E)
 \qquad(E\text{ Borel}).
\]
The capacity domain and this normalization are fixed throughout.

\begin{proposition}[Exact finite-dimensional curvature images]
\label{prop:r5-finite-curvature-tv}
Let $W\subset\Omega$ be an ordinary open, and let
$\mathcal F\subset\mathcal P(\Omega)^q$ be a family of tuples.
Suppose that for each $i=1,\ldots,q$ there is a fixed
finite-dimensional real vector space $E_i$ of original
$(1,1)$-currents on $W$ such that
\[
             D[u_i]|_W\in E_i
       \qquad((u_1,\ldots,u_q)\in\mathcal F).
\]
These are exact equalities in the ambient original current
space on the whole same open $W$.

Then the restrictions to $W$ of
\[
              \mu_{\mathbf u}=T_{S,q}(u_1,\ldots,u_q)
                     \qquad(\mathbf u\in\mathcal F)
\]
lie in a fixed finite-dimensional vector space of signed Radon
measures, of dimension at most $\prod_i\dim E_i$.
If all slots come from one space $E$ of dimension $m$, the
bound can be replaced by $\binom{m+q-1}{q}$, with the
zero-space case interpreted as the zero image.

For every net $(\mathbf u_a)$ in this family, local weak
convergence $\mu_{\mathbf u_a}|_W\to\nu$ on all original
compact smooth scalar tests implies
\begin{equation}
       |\mu_{\mathbf u_a}-\nu|(L)\longrightarrow0
                      \qquad(L\Subset W\text{ compact}).
 \label{eq:r5-finite-curvature-tv}
\end{equation}
No finite total mass on $W$ is required.
\end{proposition}

\noindent\textbf{Status: \texttt{DIRECT\_COROLLARY}} of
the arbitrary-background quotient, signed potential-current
polarization, and finite-dimensional linear algebra.

\begin{proof}
Replace $E_i$ by the span of the curvatures actually occurring
in its $i$th slot.  Choose a basis
\[
       \alpha_{i,l}=D[b_{i,l}]|_W,\qquad1\le l\le m_i,
\]
from those occurring curvatures, retaining their original
potentials $b_{i,l}\in\mathcal P(\Omega)$.  Write
$D[u_i]|_W=\sum_l c_{i,l}\alpha_{i,l}$.
The current-valued operation and its real multilinear signed
extension from Section~\ref{sec:r5-signed-potential-current-products}
give, on $W$,
\begin{equation}
 \mu_{\mathbf u}
   =\sum_{l_1,\ldots,l_q}
        \left(\prod_{i=1}^q c_{i,l_i}\right)
        T_{S,q}(b_{1,l_1},\ldots,b_{q,l_q}).
 \label{eq:r5-finite-curvature-expansion}
\end{equation}
All the currents on the right are fixed positive Radon
measures on $W$.  To make the dependence on the quotient
explicit, a linear relation among potentials with zero
ambient curvature is a locally bounded classical difference;
the quotient theorem makes its insertion against each fixed
positive expanded tail vanish.  Finite polarization therefore
respects precisely the displayed ambient current relation.
There is no assumption that a signed linear combination of
the potentials is itself psh.

The expansion proves the product bound on dimension.  If all
slots use one $E$, choose one basis from the union of the
occurring curvatures.  Symmetry identifies the products whose
indices differ only by permutation; there are
$\binom{m+q-1}{q}$ possible multisets of $q$ indices.
If one occurring curvature space is zero, that slot has the
same curvature as the zero potential and the product vanishes
by the quotient theorem.

Let $M$ be the resulting finite-dimensional measure space,
and choose a basis $\tau_1,\ldots,\tau_d$ of $M$.
Original smooth compact scalar tests separate these signed
Radon measures.  Consequently there exist
$\phi_1,\ldots,\phi_d\in A_c^{0,0}(W)$ such that the matrix
\[
                   A_{st}=\langle\tau_t,\phi_s\rangle
\]
is invertible.  Indeed, the common annihilator in $M$ of all
original tests is zero, and a finite subcollection of those
linear functionals already has zero annihilator.
Writing $\mu_{\mathbf u_a}=\sum_t x_{a,t}\tau_t$, weak
convergence on these particular tests gives
\[
 (x_{a,1},\ldots,x_{a,d})^{\mathsf T}
   =A^{-1}\bigl(\langle\mu_{\mathbf u_a},\phi_s\rangle\bigr)_s
          \longrightarrow (x_1,\ldots,x_d)^{\mathsf T}.
\]
The measure $\sum_t x_t\tau_t$ has the same original scalar
pairings as $\nu$, so it equals $\nu$.  For every compact
$L\Subset W$,
\[
 |\mu_{\mathbf u_a}-\nu|(L)
       \le\sum_{t=1}^d |x_{a,t}-x_t|\,|\tau_t|(L)
                   \longrightarrow0.
\]
This proves the assertion, including its local-mass scope.
\end{proof}

\begin{corollary}[Sharp diagonal compact no-gain under the exact hypothesis]
\label{cor:r5-finite-curvature-compact-no-gain}
Let $K_0\Subset W\subset\Omega$, with $K_0$ compact and $W$
ordinary open.  Let $u_a\in\mathcal P(\Omega)$, and assume
that all $D[u_a]|_W$ lie in one fixed finite-dimensional space
of original ambient currents.  If
$T_{S,q}(u_a,\ldots,u_a)|_W\to\nu$ locally weakly, then
\[
 \nu(K_0)=\lim_aT_{S,q}(u_a,\ldots,u_a)(K_0)
                               \le C_S(K_0;\Omega).
\]
In particular, in remaining degree one, a normalized completed
potential with a global normalized defining sequence satisfying
this exact hypothesis has no completed compact-mass gain on
$K_0$.
\end{corollary}

\noindent\textbf{Status: \texttt{DIRECT\_COROLLARY}}.

\begin{proof}
The proposition gives convergence in total variation on $K_0$.
Every term before passage to the limit is a permitted diagonal
tester on the original whole domain $\Omega$; hence its mass
is at most $C_S(K_0;\Omega)$.  This proves the exact constant
one.  For a normalized completed defining sequence, the
existing completed calculus supplies the required weak
convergence.  This last invocation does not assert a
completed representation for every higher-degree weak limit.
\end{proof}

The same measure-space argument applies to mixed families,
but its constant-one capacity bound is then with respect to
that mixed family's own raw supremum.  It does not assert a
constant-one mixed-to-diagonal capacity inequality.
For the cutoff version of C1, choose $W$ inside the ordinary
neighborhood on which the fixed cutoff equals one.  The
restricted weak limit is unchanged there, so the corollary
applies without changing the capacity domain.

\begin{example}[Strong approximation still permits compact-mass gain]
\label{ex:r5-annulus-compact-mass-jump}
Let $\lambda=-\log|\pi|>0$, and on the original ordinary
annulus
\[
            \Omega=\{|\pi|<|T|<|\pi|^{-1}\}
\]
take $S=[1]$.  Write $t=\log|T|$, let $\zeta_s$ be its
radius-$e^s$ Gauss point, and put
\[
 s_j=\frac{\lambda}{j+2},\qquad
 u_j=\frac{\max(0,t-s_j)}{\lambda},\qquad
 u=\frac{\max(0,t)}{\lambda}.
\]
These are actual continuous classical psh functions on the
same $\Omega$, with $0\le u_j,u\le1$, and
\[
        u_j\uparrow u,\qquad
        \|u_j-u\|_{L^\infty(\Omega)}\le s_j/\lambda.
\]
The original curve Laplacian and degree-one torus calibration
give the equalities of Radon measures, hence of currents on
all original tests,
\[
           D[u_j]=\lambda^{-1}\delta_{\zeta_{s_j}},
           \qquad D[u]=\lambda^{-1}\delta_{\zeta_0}.
\]
Thus for the fixed compact $K_0=\{\zeta_0\}$,
\begin{equation}
               D[u_j](K_0)=0,
               \qquad D[u](K_0)=\lambda^{-1}.
 \label{eq:r5-annulus-compact-jump}
\end{equation}

Even the normalized psh-dual distance tends to zero.  Indeed,
for every $a\in\mathcal P(\Omega)$ the function
$s\mapsto a(\zeta_s)$ is convex on $(-\lambda,\lambda)$.
This follows by the one-dimensional annulus maximum principle,
or by Lemma~\textup{R4--51.1}.  Its range
lies in $[0,1]$, so its Lipschitz constant on
$[-\lambda/2,\lambda/2]$ is at most $2/\lambda$, by the
two endpoint secant inequalities.  Therefore
\[
 \sup_{a\in\mathcal P(\Omega)}
       \left|\int a\,(D[u_j]-D[u])\right|
                       \le 2s_j/\lambda^2\longrightarrow0.
\]
By contrast the total variation of $D[u_j]-D[u]$ on the
fixed compact closed subannulus
$\{|\pi|^{1/2}\le|T|\le|\pi|^{-1/2}\}$ is $2/\lambda$
for every $j$.

There is also local strong convergence for each fixed faithful
degree-one energy form.  Set $f_j=u_j-u$, and take any original
compact smooth $\chi$.  The bounded BT Green and product
identities, valid in the closed form by
Proposition~\textup{R4--86.1}, give
\[
 \begin{split}
 \mathcal E(\chi f_j,\chi f_j)
   &=\int f_j^2\,\Gamma_{[1]}(\chi)
                     -\int\chi^2 f_j\,D[f_j]\\
   &=\int f_j^2\,\Gamma_{[1]}(\chi)
                       +\chi(\zeta_{s_j})^2s_j/\lambda^2\\
   &\le \frac{s_j^2}{\lambda^2}
                       \Gamma_{[1]}(\chi)(\Omega)
           +\frac{s_j}{\lambda^2}\|\chi\|_\infty^2
                       \longrightarrow0.
 \end{split}
\]
Here $f_j(\zeta_0)=0$ and
$f_j(\zeta_{s_j})=-s_j/\lambda$ fix the sign of the atomic
term.  The accompanying local $L^2$ convergence follows from
the uniform bound and local finiteness of the reference
measure.  Uniform convergence also implies convergence in
every local outer-capacity pair, since each fixed positive-gap
set is eventually empty.
\end{example}

\noindent\textbf{Status: \texttt{DIRECT\_COROLLARY}} of
the explicit annulus calculation; \texttt{FALSE} for an inference
from these strong potential or psh-dual convergences to
continuity of mass on an arbitrary fixed compact.
This is not a counterexample to C1: the limit potential $u$
itself is a permitted tester and realizes the limiting mass.
The point masses at the distinct $\zeta_{s_j}$ are linearly
independent on every fixed neighborhood of $\zeta_0$ after
discarding finitely many terms, so the exact finite-dimensional
curvature hypothesis does not hold in this example.

\paragraph{The remaining gate.}
For arbitrary $S$ in remaining degree one, the normalized
completed class still supplies only the outer compact value
\[
 C_{\mathcal A}(K_0;\Omega)=C_S^*(K_0;\Omega).
\]
The required general assertion is still that for every
normalized completed $a$ and every $\varepsilon>0$ there
exists $v\in\mathcal P(\Omega)$ with
\[
                 D(vS)(K_0)\ge D(aS)(K_0)-\varepsilon.
\]
The proposition proves this for exact finite-curvature defining
families.  No such hypothesis or approximation theorem is
known here for an arbitrary family.  The new potential quotient
controls exact relations, and finite Shilov evaluations control
the exact harmonic kernel.  Neither statement turns approximate
curvature relations into exact ones.  The annulus example shows
why uniform, outer-capacity, local form, or psh-dual density
alone does not fill that logical gap.  No compact or Borel
outer regularity for the full arbitrary-background capacity is
claimed by this section.

% END source: sections/compact_capacity_finite_curvature.tex


% BEGIN canonical R5 source: sections/global_finite_curvature_compactness.tex
\section{Global finite-curvature families retain normalized classical limits}
\label{sec:r5-global-finite-curvature-compactness}

\noindent\textbf{Status: \texttt{DIRECT\_COROLLARY}} of
the actual kernel theorem, finite Shilov control, and the
existing bounded BT calculus.  No countability hypothesis on
the domain is imposed.  The global finite-dimensional
hypothesis is stronger than the neighborhood hypothesis in
Proposition~\ref{prop:r5-finite-curvature-tv}.

\begin{proposition}
\label{prop:r5-global-finite-curvature-compactness}
Let $\Omega$ be an ordinary open in an allowed original-$K$
space of pure dimension $n\ge1$, where $K/\mathbf Q_p$ is
finite with its original absolute value.  Let $E$ be a
finite-dimensional real vector space of original $(1,1)$-currents
on the \emph{whole} $\Omega$, and put
\[
 \mathcal P_E=
 \{u\in\mathcal C_b(\Omega):0\le u\le1,\ D[u]\in E\}.
\]
Then $\mathcal P_E$ is compact for local uniform convergence
of the actual functions.

For every fixed symmetric original-positive separately closed
$(p,p)$-current $S$, with $q=n-p\ge1$, the image
\[
       \{T_{S,q}(u,\ldots,u):u\in\mathcal P_E\}
\]
is locally weakly compact and closed.  Every one of its weak
closure points has the same-domain normalized classical
representative displayed in this formula.  On this image
weak convergence of nets implies convergence in total
variation on each compact subset.

The corresponding assertions hold for the mixed image of
$\mathcal P_{E_1}\times\cdots\times\mathcal P_{E_q}$ when
each $E_i$ is finite-dimensional on the whole $\Omega$.
\end{proposition}

\begin{proof}
We first record a compactness property of the exact kernel.
For $M\ge0$, the set
\[
 \mathscr H_M
       =\{H\in\mathscr{PH}(\Omega): |H(x)|\le M
                                      \text{ for every }x\in\Omega\}
\]
is a closed subset of the pointwise product
$[-M,M]^\Omega$.  Indeed a pointwise limit of a net in this
set is finite at every actual point, and
Corollary~\ref{cor:r5-exact-kernel-pointwise-stability} makes
it pluriharmonic and the convergence locally uniform.
Equivalently, the finite-evaluation estimate on every compact
stage shows directly that pointwise and local uniform
convergence have the same nets on $\mathscr H_M$.
Tychonoff's theorem therefore makes $\mathscr H_M$ compact
also for local uniform convergence.  This argument uses no
countable exhaustion or sequential compactness assertion.

Replace $E$ by the span of the curvatures actually occurring
in $\mathcal P_E$, and choose a basis
$\alpha_l=D[b_l]$, $1\le l\le m$, with
$b_l\in\mathcal P_E$.
For $u\in\mathcal P_E$ write
\[
                         D[u]=\sum_l c_l(u)\alpha_l.
\]
Finitely many original complementary compact smooth tests
extract the coordinates $c_l(u)$, because such tests
separate the independent currents $\alpha_l$.
Each coordinate-extracting test $\eta$ has a fixed compactly
supported top form $D\eta$.  Its integration measure is a
finite signed Radon measure, and $0\le u\le1$ gives directly
\[
       |\langle D[u],\eta\rangle|
          =\left|\int_\Omega u\,D\eta\right|
                         \le |\mu_{D\eta}|(\Omega).
\]
Thus these finitely many pairings are uniformly bounded;
no CLN estimate is needed for this particular step.
Consequently all coefficient vectors lie in a fixed compact
box $B\subset\mathbf R^m$.

The actual function
\[
                        H_u=u-\sum_l c_l(u)b_l
\]
is locally a difference of locally bounded classical psh
functions and has zero ambient curvature on all of $\Omega$.
Theorem~\ref{thm:r5-bounded-kernel-regularity} gives
$H_u\in\mathscr{PH}(\Omega)$.
Because the basis potentials also satisfy $0\le b_l\le1$
on the whole domain,
\[
                |H_u|\le1+\sum_l|c_l(u)|\le M
                          \quad\text{everywhere on }\Omega
\]
for one finite $M$.  Thus $\mathcal P_E$ is contained in
the image of the compact space $B\times\mathscr H_M$ under
the map
\[
                     (c,H)\longmapsto H+\sum_l c_l b_l.
\]
This map is continuous for local uniform convergence, since
the finitely many $b_l$ are bounded.  Its image is therefore
a compact set of actual locally bounded functions.

The subset $\mathcal P_E$ is closed in that compact image.
To check this for nets, a local uniform limit of its members
still has values in $[0,1]$ and is upper semicontinuous.
After every complete extension and analytic curve pullback,
the convergence is uniform on each compact curve neighborhood:
its image in $\Omega$ is compact.  The subharmonic class is
closed under such uniform limits, so the limit is classical
psh.  If desired, this last local assertion for a net reduces
on each compact neighborhood to a uniformly convergent
sequence chosen from the net.  Distributional
differentiation is continuous for local uniform convergence,
and $E$ is a finite-dimensional, hence closed, subspace of
the Hausdorff original-current space.  The limit curvature
therefore remains in $E$.  This proves compactness of
$\mathcal P_E$.  When $m=0$ the sum is empty and the same
argument simply uses the compact kernel ball.

Fix $S$.  On tuples of normalized bounded classical
potentials the mixed top operator is continuous for local
uniform convergence, including for nets.  One may use the
established capacity continuity, or verify the assertion
directly by telescoping the slots.  Against any fixed
original compact smooth scalar test $\phi$, each telescoping
term has the form
\[
              \langle (u_i-v_i)T_i,D\phi\rangle,
\]
where $T_i$ is a positive mixed tail with the fixed background
$S$.  Coefficient-Radon control and relative CLN bound its
absolute value by a fixed constant times
$\|u_i-v_i\|$ on a fixed compact neighborhood, uniformly
over the other normalized slots.  This tends to zero under
local uniform convergence.  The same CLN mass bounds and
uniform approximation of compact continuous scalar tests
give continuity for local vague convergence of Radon measures.

The diagonal image is therefore a compact subset of that
Hausdorff measure space, hence is closed, and every closure
point has a representative in $\mathcal P_E$.
For mixed images apply the same argument to the finite
product of the compact spaces $\mathcal P_{E_i}$.
Finally Proposition~\ref{prop:r5-finite-curvature-tv}, with
$W=\Omega$, gives the asserted local total-variation
convergence inside these images.
\end{proof}

This conclusion requires exact finite-dimensional control on
the entire capacity domain and retains the original pointwise
normalization there.  Finite-dimensional curvature span only
near one compact core gives the earlier local-TV result, but
does not supply this global normalized representative.
No finite dimension of the pluriharmonic stalk, compactness
of the whole unrestricted normalized family, or
finite-dimensional approximation theorem for general C1 is
asserted.

% END source: sections/global_finite_curvature_compactness.tex


% BEGIN canonical R5 source: sections/two_intersection_image_obstruction.tex
\section{The classical relative image need not be TV-closed without boundary conditions}
\label{sec:r5-two-intersection-image}

The obstruction in this section concerns exact classical
representatives of relative curvature measures.  No boundary
value is prescribed.  The limiting measure nevertheless satisfies
the desired sharp capacity bound, so this is not a counterexample
to compact no-gain C1.

\begin{lemma}[A necessary carrier-incidence condition]
\label{lem:r5-carrier-incidence}
Let $\Omega$ be an allowed original-$K$ space of pure dimension
$n$, and let $\iota_i:\Delta_i\hookrightarrow\Omega$,
$1\le i\le m$, be closed analytic embeddings, where each
$\Delta_i$ is an ordinary open unit disc.  Suppose that the
images $C_i$ have finite pairwise intersections, and let $Z$
be the union of these intersection sets.  Take $h_i\in\mathcal C_b(\Delta_i)$
such that $D[h_i]$ gives no mass to $\iota_i^{-1}(Z)$, and set
\[
 S=\sum_i[C_i],\qquad \nu=\sum_i(\iota_i)_*D[h_i].
\]
If $u\in\mathcal C_b(\Omega)$ satisfies $D(uS)=\nu$, then
there are constants $c_i\in\mathbf R$ with
\[
                 u\circ\iota_i=h_i+c_i\quad\text{on }\Delta_i.
\]
In particular every intersection point $P\in C_i\cap C_l$
imposes the actual-value relation
\begin{equation}
 c_l-c_i=h_i(\iota_i^{-1}P)-h_l(\iota_l^{-1}P).
 \label{eq:r5-carrier-incidence}
\end{equation}
The sum of the right sides along every cycle of the incidence
multigraph must be zero.  This is only a necessary condition;
no ambient extension theorem is asserted.
\end{lemma}

\noindent\textbf{Status: \texttt{DIRECT\_COROLLARY}} of
the actual curve kernel, positivity, and integration-current
compatibility.

\begin{proof}
Write $\mu_i=D[u\circ\iota_i]$.  These are positive measures,
and first-order integration compatibility gives
$D(uS)=\sum_i(\iota_i)_*\mu_i$.  On $C_i\setminus Z$ no
other carrier contributes.  Thus $\mu_i$ agrees there with
$D[h_i]$.  The target measure gives zero mass to $Z$, so
positivity of all the summands forces each $\mu_i$ to give
zero mass to its preimage as well.  Consequently
$\mu_i=D[h_i]$ on the whole disc.

The actual curve kernel theorem makes $u\circ\iota_i-h_i$
pluriharmonic.  Every finite pluriharmonic function on an
ordinary open unit disc is constant: exhaust it by nested
closed discs, apply the maximum and minimum principles on
each disc with its singleton Shilov boundary, and note that
the constants agree on the overlaps.  This proves the
displayed representation.  The single actual value $u(P)$
at an intersection gives \eqref{eq:r5-carrier-incidence}.
Summing it along a cycle makes the constants telescope to
zero.
\end{proof}

\begin{proposition}[A normalized boundary-free image obstruction in every $n\ge2$]
\label{prop:r5-two-intersection-tv-obstruction}
For every $n\ge2$ there are an ordinary unit polydisc
$\Omega$ over the original finite extension $K/\mathbf Q_p$,
a symmetric original-positive separately closed current $S$
of bidegree $(n-1,n-1)$, normalized continuous classical
potentials $a_j$ on $\Omega$, and a finite positive measure
$\nu$ such that
\[
 0\le a_j\le1,\qquad a_j\uparrow a,\qquad
          \|D(a_jS)-\nu\|_{\rm TV}\longrightarrow0,
\]
but there is no actual $u\in\mathcal C_b(\Omega)$ with
$D(uS)=\nu$, even without any normalization or boundary
condition on $u$.

The actual limit $a$ belongs to the normalized relative
completion and satisfies $D(aS)=\nu$.  Moreover the
approximating measures can be chosen to dominate $\nu$:
\begin{equation}
 D(a_jS)\ge\nu\quad\text{for every sufficiently large }j,
 \qquad
 \nu(E)\le C_S(E;\Omega)\quad(E\text{ Borel}).
 \label{eq:r5-two-intersection-cap-bound}
\end{equation}
Here $C_S$ is the diagonal raw capacity in remaining degree
one on the same whole $\Omega$.
\end{proposition}

\noindent\textbf{Status: \texttt{NEW\_PROVED}} for this
boundary-free strengthening.  The assertion that the actual
classical relative image is always TV-closed even without
boundary conditions is \texttt{FALSE}.  The proof uses
previously established curve and completed-calculus results;
it does not modify them.

\paragraph{Provenance of the strengthening.}
Corollary~\textup{R4--65.2}
already gives two coordinate carriers, a fixed zero-boundary
condition, an exact TV error, and domination of the limiting
measure by classical approximating measures.  The completed
carrier construction and the distinction between capacity
regularity and exact representability are therefore prior
inputs, not new claims here.  The new point is the use of
\emph{two} actual intersections: their incidence cycle
eliminates the otherwise free additive constants and proves
nonexistence after all boundary conditions have been removed.

\begin{proof}
Put $\lambda=-\log|\pi|>0$ and fix $0<\delta<1/2$.
Write the coordinates of the ordinary unit polydisc as
$(x,y,z_3,\ldots,z_n)$, omitting the last coordinates for
$n=2$.  Define two closed analytic embedded unit discs by
\[
 \iota_1(t)=(t,0,0,\ldots,0),\qquad
 \iota_2(t)=(t,t(t-\pi),0,\ldots,0).
\]
They meet precisely at the two original rigid points with
$t=0$ and $t=\pi$.  Set $S=[C_1]+[C_2]$ for their images.
These integration currents have exactly the stated positivity,
closedness, and bidegree on the boundaryless ordinary domain.

On the open unit disc put
\[
 h(t)=\delta+(1-2\delta)
                  \max\{0,1+(2\lambda)^{-1}\log|t|\},
 \qquad c=\frac{1-2\delta}{2\lambda}>0.
\]
This is continuous classical psh, with
$\delta\le h\le1-\delta$, and the original curve Laplacian
calibration gives
\[
     D[h]=c\delta_{\zeta_{0,|\pi|^2}},\qquad
     h(0)=\delta,\qquad h(\pi)=1/2.
\]
Define $z_* =\iota_2(\zeta_{0,|\pi|^2})$ and
$\nu=c\delta_{z_*}$.  Its support is disjoint from $C_1$:
at that Gauss point the norm of $t(t-\pi)$ is
$|\pi|^3>0$.  In particular the target gives zero mass to
both intersection points.

Define actual potentials on the whole polydisc by
\begin{equation}
 a_j(x,y,z)=\max\{0,h(x)+j^{-1}\log|y|\},
 \qquad \log|0|=-\infty.
 \label{eq:r5-two-intersection-row}
\end{equation}
They are continuous classical psh and satisfy
$0\le a_j\le1-\delta$.  For completeness, away from the
zero coordinates the expression is a convex coordinatewise
nondecreasing function of the logarithmic unit norms.  Near
$x=0$ the inner maximum in $h$ is constant, and near $y=0$
the outer maximum in \eqref{eq:r5-two-intersection-row} is
identically zero, using the common upper bound for $h$.
Thus these expressions extend continuously across both zero
loci.  Truncating the analytic logarithms and taking the same
finite maxima verifies classicality after every complete
extension.  Since $\log|y|\le0$, the sequence increases
pointwise to
\[
 a(x,y,z)=
   \begin{cases}
      h(x),&y\ne0,\\
      0,&y=0.
   \end{cases}
\]
Corollary~\textup{R4--18.5} puts this actual Borel
limit in the relative completion, using the same globally
normalized defining row on $\Omega$.

We compute every relative measure for
$j>\max\{3\lambda/\delta,4\lambda\}$.
The restriction of $a_j$ to $C_1$ is zero.  On $C_2$ it is
\[
       b_j(t)=\max\{0,h(t)+j^{-1}\log|t(t-\pi)|\}.
\]
Put
\[
             r_{0,j}=\exp(\lambda-j\delta),\qquad
             r_{\pi,j}=\exp(\lambda-j/2).
\]
The chosen inequalities give $r_{0,j}<|\pi|^2$ and
$r_{\pi,j}<|\pi|$.  On the disc $|t|<|\pi|^2$ one has
$h(t)=\delta$ and $|t-\pi|=|\pi|$, so
\[
 b_j(t)=\max\{0,\delta+j^{-1}(\log|t|-\lambda)\}.
\]
It is zero on the disc of radius $r_{0,j}$ and gains outward
logarithmic slope $1/j$ at its Gauss boundary.
On $|t-\pi|<|\pi|$ one has $|t|=|\pi|$ and $h(t)=1/2$,
so likewise
\[
 b_j(t)=\max\{0,1/2+j^{-1}(\log|t-\pi|-\lambda)\},
\]
with the same slope jump at radius $r_{\pi,j}$.

Outside these two smaller discs the nonzero branch is
nonnegative.  To check that no additional contact locus was
omitted, for $|\pi|^2\le|t|<|\pi|$ it is at least
$\delta-3\lambda/j>0$; for $|t|>|\pi|$ it is at least
$1/2-2\lambda/j>0$; and the remaining case $|t|=|\pi|$
is exactly the second displayed radial expression or has
$|t-\pi|=|\pi|$, where the same positive bound applies.
At the circle $|t|=|\pi|^2$ the branch is strictly active.
Its curvature there is therefore $D[h]$.  On the complement
of the two zeros, the additional polynomial logarithm is
harmonic.  The one-dimensional slope formula, or original
Poincar\'e--Lelong and ordinary maximum locality on the
displayed annuli, now gives on all original curve tests
\[
 D[b_j]
   =D[h]+j^{-1}\delta_{\zeta_{0,r_{0,j}}}
                       +j^{-1}\delta_{\zeta_{\pi,r_{\pi,j}}}.
\]
Both centers are $K$-rational, so each annular slope has
degree-one calibration.  The radii need not belong to
$|K^\times|$; the Gauss points are points of the original
$K$-analytic disc and the same slope formula applies.

First-order integration compatibility and linearity in $S$
therefore give the exact original-current identity
\begin{equation}
 D(a_jS)=\nu+j^{-1}\delta_{\iota_2(\zeta_{0,r_{0,j}})}
                    +j^{-1}\delta_{\iota_2(\zeta_{\pi,r_{\pi,j}})}.
 \label{eq:r5-two-intersection-exact-measure}
\end{equation}
In particular its total-variation distance to $\nu$ on the
whole $\Omega$ is exactly $2/j$.  The completed calculus
identifies this limit with $D(aS)$, by
Corollary~\textup{R4--18.6}.

Suppose that some actual $u\in\mathcal C_b(\Omega)$ had
$D(uS)=\nu$.  The two intrinsic curve Laplacians of its
restrictions are positive.  Their pushforwards sum to
$\nu$, whose support is disjoint from $C_1$.  The first
Laplacian is therefore zero, and the second is $D[h]$.
As in Lemma~\ref{lem:r5-carrier-incidence}, the actual
restrictions have the forms
\[
                    u\circ\iota_1=c_1,\qquad
                    u\circ\iota_2=h+c_2.
\]
At the two common actual points this would give simultaneously
$c_1=\delta+c_2$ and $c_1=1/2+c_2$, contrary to
$\delta<1/2$.  This proves nonexistence without imposing
boundary values or any global bound on $u$ beyond its
membership in the locally bounded classical class.

Finally \eqref{eq:r5-two-intersection-exact-measure} gives
$\nu\le D(a_jS)$ for every sufficiently large $j$.
Since each $a_j$ is an admissible normalized tester on the
same whole $\Omega$, for every Borel $E$ one has
\[
                  \nu(E)\le D(a_jS)(E)\le C_S(E;\Omega).
\]
This proves \eqref{eq:r5-two-intersection-cap-bound} with the
exact constant one.
\end{proof}

\paragraph{What is, and is not, excluded.}
The convergence in \eqref{eq:r5-two-intersection-exact-measure}
concerns the relative measures $D(a_jS)$, not the ambient
curvatures $D[a_j]$.  The counterexample shows that finding an
exact actual classical representative of every completed
relative measure is a strictly stronger task than the C1
mass bound: for this specific limit no such representative
exists, while every sufficiently late classical approximant
already dominates it.  It makes no assertion that general
C1 fails for this background or for arbitrary backgrounds.
The arbitrary-background quotient remains valid; its
hypothesis concerns two actual classical potentials, and the
completed $a$ here cannot be replaced by any such potential
with the same relative curvature.  Finite pluriharmonic
adjustments cannot remove the two incompatible actual
intersection values.

% END source: sections/two_intersection_image_obstruction.tex


% BEGIN canonical R5 source: sections/gauss_pushforward_obstruction.tex
\section{A wall obstruction to a naive Gauss-current regularization}
\label{sec:r5-gauss-pushforward-obstruction}

The following calculation rules out one proposed construction of a
regularized background.  It is not a counterexample to potential
independence for actual classical psh functions.

\begin{proposition}[The piecewise Jacobian rule does not preserve closedness]
\label{prop:r5-gauss-pushforward-not-closed}
Let $K$ retain its original absolute value, and let
\[
 A=\{a<\log|T|<b\}\subset\mathbf G_{m,K}^{\rm an},
 \qquad a<c<b.
\]
Write $s=\log|T|$ and $r=e^c$.  The radius-$r$ Gauss inflation of
a point, formed by adjoining a Gauss variable of radius $r$ to its
completed residue field and replacing $T$ by $T+Z$, satisfies
\[
 s\circ\tau_r=\max(s,c).
\]
Consider the proposed operation on $[1]_A$ which pulls back a
radial top test by the ordinary Jacobian of this piecewise affine
map, with no wall term.  On every $\phi\in C_c^\infty((a,b))$ this
rule is
\begin{equation}
 \left\langle R_r[1],\phi(s)d's\wedge d''s\right\rangle
       =\int_c^b\phi(t)\,dt.
 \label{eq:r5-radial-jacobian-rule}
\end{equation}
No separately closed original current can satisfy
\eqref{eq:r5-radial-jacobian-rule} on all these tests.
\end{proposition}

\noindent\textbf{Status: \texttt{FALSE}} for this proposed
closed-current regularization; the following elementary obstruction
is \texttt{NEW\_PROVED}.  The ordinary annular integration
calibration used in the proof is \texttt{EXTERNAL\_KNOWN}.

\begin{proof}
For the Gauss seminorm over $\mathscr H(x)$ one has
$|T(x)+Z|=\max(|T(x)|,r)$, which proves the displayed formula for
$s\circ\tau_r$.  The derivative of $t\mapsto\max(t,c)$ is zero
on $t<c$ and one on $t>c$.  Pulling back both differential factors
therefore gives zero on the first interval and the original top
form on the second.  The ordinary annulus has tropical weight one,
so integration against $[1]_A$ gives
\eqref{eq:r5-radial-jacobian-rule}.  The value assigned to the
derivative at $t=c$ does not change this integral.

Choose $\psi\in C_c^\infty((a,b))$ with $\psi(c)=1$, and set
$\alpha=\psi(s)d''s$.  This is an original compact smooth
$(0,1)$-test, and
\[
 d'\alpha=\psi'(s)d's\wedge d''s.
\]
Consequently the proposed rule gives
\[
 \langle R_r[1],d'\alpha\rangle
       =\int_c^b\psi'(t)\,dt=-1.
\]
This is incompatible with $d'R_r[1]=0$.  Interchanging primes
gives the corresponding obstruction to separate $d''$-closedness.
\end{proof}

\paragraph{What the calculation excludes.}
The finite-wall pairing theorem permits compatible piecewise smooth
tests to be evaluated against closed positive currents.  It does
not state that piecewise differentiation commutes with $D$ or that
its wall terms vanish.  Indeed $Ds=0$, whereas
$D\max(s,c)$ is the unit Dirac curvature at the radius-$e^c$
Gauss point.  Thus compatible piecewise test approximation alone
cannot turn Gauss inflation into a closed-current smoothing
operator.  A construction using integration over the full analytic
product fiber would require its own pullback, calibration, and
closedness argument; it is not supplied by this failed Jacobian
rule.

% END source: sections/gauss_pushforward_obstruction.tex


% BEGIN canonical R5 source: sections/type_four_perturbation_obstruction.tex
\section{A failed Gauss-perturbation shortcut}
\label{sec:r5-gauss-perturbation-obstruction}

One possible route to scalar kernel functoriality was to perturb a
finite parameter tuple by arbitrarily small independent Gauss
variables, hoping that the resulting parameter point would always
be Abhyankar.  Exact norm-profile stability would then avoid any
limit with moving coefficient measures.  The premise is false,
even for one parameter over the original finite field $K$.

\begin{example}[A non-Abhyankar point fixed by every sufficiently small perturbation]
\label{ex:r5-gauss-perturbation-fixed-type-four}
Let $\pi$ be a uniformizer of $K$, and normalize its additive
valuation by $v(\pi)=1$.  There are a complete extension $M/K$, an
element $a\in M$, and a constant $r_0>0$ such that the point
$x\in\mathbf A_K^{1,\rm an}$ represented by $a$ is Zariski generic
and non-Abhyankar, but for every $0<\epsilon<r_0$ the point
represented by $a+U$ in the radius-$\epsilon$ Gauss extension
$M_\epsilon$ is exactly $x$.
\end{example}

\noindent\textbf{Status: \texttt{FALSE}} for the proposed universal
Abhyankar-perturbation premise.  The following explicit construction
also records why the obstruction occurs over the original $K$,
not only over a changed ground field.

\begin{proof}
Let $C=\widehat{\overline K}$.  Choose increasing primes
$\ell_j$ so rapidly that
\[
 \ell_j>j\prod_{i<j}\ell_i.
\]
Put $q_j=1-1/\ell_j$, choose
$\theta_j\in\overline K$ with $\theta_j^{\ell_j}=\pi^{\ell_j-1}$,
and set
\[
 c_j=\sum_{i=1}^j\theta_i,\qquad
 r_j=|\pi|^{q_{j+1}},\qquad r_0=|\pi|.
\]
The closed discs $D_j=D(c_j,r_j)$ in $C$ are nested, since
$|c_{j+1}-c_j|=r_j$, and $r_j\downarrow r_0>0$.

They have empty intersection in $C$.  To check this directly,
first fix $c\in\overline K$ and put $d=[K(c):K]$.
For $j>d$, the finite extension
\[
 H_j=K(c,\theta_1,\ldots,\theta_{j-1})
\]
has degree at most $d\prod_{i<j}\ell_i<\ell_j$.  Its ramification
index is therefore smaller than $\ell_j$, so its value group
does not contain $q_j=(\ell_j-1)/\ell_j$.
Since $c-c_{j-1}\in H_j$, the two summands in
$c-c_j=(c-c_{j-1})-\theta_j$ cannot have equal finite valuation.
Thus
\[
 v(c-c_j)=\min\{v(c-c_{j-1}),q_j\}\leq q_j<q_{j+1},
\]
also when $c=c_{j-1}$.  Hence $c\notin D_j$.
If $b\in C$ belonged to every $D_j$, approximate it by
$c\in\overline K$ with $|b-c|<r_0$.  Since every $r_j>r_0$,
the ultrametric inequality would put $c$ in every $D_j$, a
contradiction.

Let $y$ be the multiplicative seminorm on $C[T]$ obtained as the
decreasing limit of the Gauss norms of these discs.  It is the
type-IV point represented by the nested discs.  Put
$M=\mathscr H(y)$ and let $a$ be the image of $T$ in $M$.
For every $c\in C$, some $D_j$ omits $c$, and nesting makes
$|c_j-c|$ constant from that stage onward.  In particular
\begin{equation}\label{eq:r5-type-four-strict-distance}
 |a-c|>r_0\qquad(c\in C).
\end{equation}
Every nonzero polynomial over $C$ splits into such linear factors;
its value at $a$ is therefore nonzero.  This also shows that the
restricted point $x$ on $\mathbf A_K^{1,\rm an}$ is Zariski generic.

For completeness, $M/C$ is immediate.  Given a nonzero rational
function $R\in C(T)$, choose $j$ so large that $D_j$ contains none
of its finitely many zeros or poles.  Its norm on every later disc
equals $|R(c_j)|$.  Moreover all its linear factors satisfy
\[
 \left|\frac{a-c}{c_j-c}-1\right|
       \leq\frac{r_j}{|c_j-c|}<1.
\]
Taking products and quotients gives
$|R(a)/R(c_j)-1|<1$.  Thus both the value group and every initial
residue of $C(a)$ are already supplied by $C$.  Completion does not
change the value group or residue field, so the same is true of
$M$.  The value group of $C$ is $|K^\times|^{\mathbf Q}$ and its
residue field is algebraic over $\widetilde K$.  Any intermediate
completed residue field, in particular $\mathscr H(x)\subset M$,
therefore has rational value-group rank zero and residue
transcendence degree zero over $K$.  Consequently $d_K(x)=0$,
whereas its Zariski closure has dimension one.  It is not
Abhyankar.

Finally take $0<\epsilon<r_0$ and give $M(U)$ the
radius-$\epsilon$ Gauss norm.  By
\eqref{eq:r5-type-four-strict-distance}, every $c\in C$ satisfies
\[
 |a+U-c|=\max\{|a-c|,\epsilon\}=|a-c|.
\]
Factor any polynomial in $K[T]$ over $C$ and use multiplicativity.
Its norm at $a+U$ is exactly its norm at $a$.  The two points of
$\mathbf A_K^{1,\rm an}$ are therefore identical for every such
$\epsilon$.  Arbitrarily small independent Gauss perturbations
do not remove this transcendence defect.
\end{proof}

This is a failure of the proposed perturbation mechanism, not of
scalar kernel functoriality.  In the one-variable proof, the
positive limiting radius supplies a different tool: the continuous
potential obtained from the uniformly Cauchy truncated logarithms
of the nested discs.  Its full-fiber curvature is treated in
Proposition~\ref{prop:r5-type-four-full-fiber-potential}.

% END source: sections/type_four_perturbation_obstruction.tex


% BEGIN canonical R5 source: sections/round5_status.tex
\section{Status change from R4 and remaining proof boundaries}
\label{sec:r5-status}

The general question (Q) in the last subsection of R4 was exactly
$D[u]=D[v]\Rightarrow D((u-v)S)=0$ for every allowed fixed $S$.
Its two stated unresolved gates were all-complete-extension
curve testing and continuity of the actual bounded classical
difference. Sections~\ref{sec:r5-scalar-extension-kernel}--
\ref{sec:r5-algebraic-continuity-geometry} prove those gates;
Theorem~\ref{thm:r5-arbitrary-background-potential-quotient}
therefore changes that precise question from \texttt{OPEN} to
\texttt{PROVED}, relative to the declared inputs.

\subsection{Focused theorem comparison}

\begingroup
\setlength{\tabcolsep}{3pt}
\begin{longtable}{@{}p{.23\textwidth}p{.26\textwidth}p{.43\textwidth}@{}}
\toprule
Statement & R4 & R5\\
\midrule
\endhead
Seven bounded BT properties for fixed arbitrary $S$ &
\texttt{PROVED} in the earlier records &
Unchanged input. Existence, positivity, separate closedness,
symmetry, potential multilinearity, ordinary restriction/gluing,
and weighted joint decreasing convergence are not recounted as
new results.\\[4pt]
Original-$K$ projection, finite-etale and analytic kernel &
\texttt{PROVED} &
Unchanged starting point for the simultaneous product induction.\\[4pt]
Every complete scalar extension $L/K$ &
Only finite-over-Gauss extensions proved; general parameter
residue fields remained \texttt{OPEN} &
\texttt{PROVED}, Theorem~\ref{thm:r5-all-complete-scalar-kernel}.
Every finite tuple of coefficients is treated by successive
one-parameter fields, including non-Gauss type-IV parameters.\\[4pt]
Every analytic map over every such $L$ &
Only the original-$K$ analytic theorem; no unrestricted
new-field assertion &
\texttt{PROVED}, Theorem~\ref{thm:r5-all-extension-analytic-kernel},
for the actual induced original potentials.\\[4pt]
Actual curve harmonicity with all extensions &
The full extension quantifier remained \texttt{OPEN} &
\texttt{PROVED}, Theorem~\ref{thm:r5-all-extension-curve-harmonicity}.
The type-I and type-IV values are recovered separately for
each of the two original subharmonic functions.\\[4pt]
Curvewise harmonicity without ambient semicontinuity &
No continuity mechanism established &
\texttt{PROVED} in the stated original-field scope by the new
Shilov bound and semistable logarithmic interpolation.
Even local boundedness and Borel measurability need not be
assumed for the finite real-valued curvewise harmonic function.\\[4pt]
$\ker D$ on actual $\mathcal C_b-\mathcal C_b$ &
Regularity \texttt{OPEN} &
\texttt{PROVED}: it is exactly the sheaf of continuous
classical pluriharmonic functions.\\[4pt]
Potential quotient for every fixed $S$ &
\texttt{OPEN} &
\texttt{PROVED}, Theorem~\ref{thm:r5-arbitrary-background-potential-quotient}.
There is no new restriction on the background.\\[4pt]
Mixed operation on bounded-potential curvature currents &
The general quotient obstruction remained &
\texttt{DIRECT\_COROLLARY}: local potentials and differentiated
slots descend and glue by
Corollaries~\ref{cor:r5-differentiated-slot-quotient} and
\ref{cor:r5-local-potential-current-products}.\\[4pt]
Signed differentiated local-potential slots &
No general potential-independent current operation &
\texttt{DIRECT\_COROLLARY}: real multilinearity on the
specified bounded classical difference-potential class,
Theorem~\ref{thm:r5-signed-potential-current-products}.\\[4pt]
Arbitrary new-field backgrounds and original induced potentials &
No unrestricted complete-extension transfer &
\texttt{DIRECT\_COROLLARY}: the background over $L$ is
arbitrary, while relations between local potentials must
originate over $K$; Section~\ref{sec:r5-induced-potential-background-quotient}.\\[4pt]
Finite pointwise limits of exact kernel functions &
Known PH finite-evaluation theory, but no general identification
of the actual current kernel &
\texttt{DIRECT\_COROLLARY}: locally uniform convergence,
including nets, with the finite-evaluation bound of
Section~\ref{sec:r5-exact-kernel-stability}.\\[4pt]
Energy stability inside the actual exact kernel &
The faithful form was available &
\texttt{DIRECT\_COROLLARY}: the exact cutoff identity and
strong local form convergence for pointwise finite kernel nets,
Section~\ref{sec:r5-exact-kernel-energy}.\\[4pt]
Relative curvature convergence modulo PH corrections &
Green and closed-form tools available &
\texttt{DIRECT\_COROLLARY}: a finite-test fixed-background
energy bound supplies a sufficient convergence condition;
Section~\ref{sec:r5-relative-energy-convergence}.\\[4pt]
Weak continuity of the descended current operation &
Not a consequence of the potential calculus &
\texttt{FALSE}: Example~\ref{ex:r5-smooth-weak-current-defect}
has smooth positive curvatures, a smooth positive limit,
uniformly bounded smooth potentials, and fixed $S=[Y]$.\\[4pt]
Q for arbitrary bounded USC representatives or the full
relative capacitary completion &
No justified extension &
\texttt{FALSE}, by the polar trace example and
Remarks~\ref{rem:r5-bounded-usc-kernel-obstruction} and
\ref{rem:r5-completed-potential-quotient-false}.\\
Only same-field or finite-extension curve tests &
Insufficient to close the complete-extension gate &
\texttt{FALSE} as a substitute: the full Gauss-point spike
passes every such curve test and has zero ambient distribution,
but is not continuous; Section~\ref{sec:r5-curve-testing-scope-obstruction}.\\
Regularity from universal vanishing on arbitrary Borel functions &
Not an asserted replacement for classical potentials &
\texttt{FALSE}: complete-pole selectors can vanish against
every allowed relative curvature and still be discontinuous,
Section~\ref{sec:r5-universal-borel-kernel}.\\
Replacing the middle cone by weak positivity &
Not permitted by the original input convention &
The two-sign vanishing step is \texttt{FALSE} in the weak cone.
The printed general strong-test order-zero inference also has
a closed real-tropical counterexample; Sections
\ref{sec:r5-weak-cone-lineality}--\ref{sec:r5-weak-order-zero-tropical}.\\
Exact finite-dimensional ambient curvature families &
The general current quotient was unavailable &
\texttt{DIRECT\_COROLLARY}: all-degree mixed images are
finite-dimensional, weak convergence is locally total variation,
and diagonal compact capacity control has constant one;
Section~\ref{sec:r5-compact-capacity-finite-curvature}.
The additional exact hypothesis is not automatic.\\
Normalized families with exact finite curvature on the whole domain &
No general actual-current kernel identification &
\texttt{DIRECT\_COROLLARY}: local uniform compactness and
retention of normalized classical representatives for every
fixed-$S$ mixed weak limit in the specified family;
Proposition~\ref{prop:r5-global-finite-curvature-compactness}.
The proof uses nets and imposes no countability assumption.\\
Classical relative image without prescribed boundary values &
Failure of TV-closedness for a prescribed zero-boundary
problem was already proved &
\texttt{FALSE} even with all boundary requirements removed:
two intersections impose incompatible actual values.
The boundary-free strengthening in
Proposition~\ref{prop:r5-two-intersection-tv-obstruction} is
\texttt{NEW\_PROVED}. Its limit still obeys the sharp
capacity bound.\\
\bottomrule
\end{longtable}
\endgroup

\subsection{What the proof adds}

The main new arguments are the all-extension kernel propagation
and the passage from actual curvewise harmonicity to ambient
continuity without assuming either sign semicontinuous. The
former preserves all original tests through full analytic fibers.
The latter supplies the regularity statement needed by the final
arbitrary-background positivity argument. The coefficient controls,
standard finite-flat Shilov facts, polynomial normalization,
vertical-divisor interpolation, and the subsequent current
linearizations retain their stated direct or external provenance.

The complete-field steps have deliberately narrower function
hypotheses than an unrestricted new-field regularization theorem:
the functions come with prescribed positive smooth rows, and the
rows used in the main proof are pullbacks of original-$K$ rows.
No statement of B over an arbitrary new field is inferred from
the existence of those particular rows.

\subsection{The eleven-direction project remains larger than Q}

The exact original-literature theorem comparison remains in the
frozen R4 record, Section 13, pp.~41--50. This continuation changes its
positive-current/mixed-product row and adds the bounded classical
kernel theorem. It does not replace the individual hypotheses
of the remaining capacity, comparison, or Dirichlet results.

\begin{longtable}{p{.28\textwidth}p{.65\textwidth}}
\toprule
Direction & Current boundary\\
\midrule
\endhead
1. Locally bounded psh & B's class and regularization remain
inputs. The actual bounded-difference curvature kernel now
has a continuous PH representative equal to its actual values.\\
2. Positive closed currents and mixed products & The seven
fixed-background properties and the general differentiated
potential quotient are now both proved.\\
3. CLN inequalities & The earlier finite-control relative
CLN theorem is unchanged. The prescribed-row auxiliary-field
estimate is used in this proof.\\
4. Symmetry and weighted monotone convergence & The fixed-$S$
bounded theorem is unchanged. Weak current convergence is
strictly weaker than its hypotheses, as the new smooth
counterexample demonstrates.\\
5. Locality and maxima & The earlier ordinary and Borel
maximum-locality results are unchanged. A strict gap set of
two discontinuous psh functions is still not automatically
an ordinary open.\\
6. BT capacities & The qualified earlier results remain.
General sharp compact control C1 and Borel outer regularity
C2 remain \texttt{OPEN}.\\
7. Quasicontinuity and convergence in capacity & The earlier
relative bounded and energy-representative conclusions retain
their stated domains and capacities. Q adds no automatic
compactness theorem for the normalized psh family.\\
8. Pluripolar and negligible sets & Relative measures can
still charge ambient pluripolar sets. The examples here use
that possibility and do not remove it by changing representatives.\\
9. Comparison and domination & The earlier compact-gap,
relative-q.e., and qualified ambient principles remain.
No new unrestricted pointwise uniqueness statement follows
from the quotient.\\
10. Relative extremal functions & The earlier constrained
variational and qualified ambient results remain. The false
unrestricted relative obstacle-contact formula is not repaired
by potential independence.\\
11. Monge--Ampere and Dirichlet problems & The qualified
existence results remain unchanged. The new two-intersection
example strengthens the earlier relative nonexistence result
by removing every boundary condition. A useful general relative
solvability criterion remains \texttt{OPEN}.\\
\bottomrule
\end{longtable}

\subsection{Minimal next gates}

There is no remaining unproved lemma in the stated proof of Q.
For the broader project the next gates remain those isolated
in R4:
\begin{enumerate}
\item \textbf{C1, sharp compact approximation.}
On one fixed normalized capacity domain, a measure $\nu$ in
the local vague closure of the classical BT measure family
should satisfy
\[
                    \nu(K_0)\leq C_S(K_0;\Omega)
\]
for every compact core $K_0$. In the remaining-degree-one
completed image, the precise missing step is lower
approximation of the completed compact mass by a normalized
classical tester on that same domain. Neither weak density
nor Q proves this lower compact-mass approximation.
Section~\ref{sec:r5-compact-capacity-finite-curvature} proves
the sharp bound for defining families whose ambient curvatures
lie exactly in one fixed finite-dimensional space near the
compact set. If this exact hypothesis holds on the whole domain,
Proposition~\ref{prop:r5-global-finite-curvature-compactness}
also supplies normalized actual classical representatives of
limits in that specified family. Its local hypothesis alone
does not supply this global conclusion.
The annulus example shows that uniform potential,
strong energy, and psh-dual convergence by themselves do not
give convergence of compact masses. The example does not
refute C1, since its normalized classical limit is admissible.
Moreover, Section~\ref{sec:r5-two-intersection-image} rules
out the stronger strategy of always representing a completed
relative measure by one exact actual classical potential:
its explicit measure satisfies the sharp capacity bound and
is dominated by normalized classical approximants, but has
no such representative even without boundary conditions.

\item \textbf{C2, ordinary-open outer regularity.}
Even sharp compact control would still need the appropriate
increasing-set/outer-capacity argument. The proved finite
tester-pasting lemma supplies the final combination step;
the required simultaneous compact extraction with vanishing
capacity and pasting losses is not supplied by this round.

\item \textbf{D, a verifiable relative solvability criterion.}
The completed constrained images and their relative-q.e.
uniqueness are available. A useful independent criterion for
a datum to lie in the actual image remains to be proved.
The unrestricted bounded-subsolution version in the global
bounded-row completion is already false and must not be
reinstated as a conjectural automatic principle.
For actual classical representatives, the new carrier-incidence
condition is an additional necessary constraint: harmonic
adjustments on unit-disc carriers are constants, and their
values must match at every actual intersection point.
\end{enumerate}

Outside the main field/function scope, a general regularity theorem
for all locally integrable original-current potentials, and a
new-field theorem for all new classical functions without
prescribed rows, are not established here. The two questions
raised in CLD2025, Remark 7.8.3, have a broader tropical or
distributional scope. The present actual bounded-classical
difference theorem must not be advertised as resolving all of
those questions.

\paragraph{Comparison with GK and arithmetic applications.}
The conclusions are original CLD-current identities. Equality
on every GK $\delta$-test needs a separate comparison theorem.
No wedge of two arbitrary positive currents and no nonclosed,
off-diagonal GK calculus is asserted. A2 remains a completed
arithmetic input. These boundaries are unchanged by the new
potential quotient.


% END source: sections/round5_status.tex

\clearpage

% BEGIN canonical R5 source: references.tex
\begin{thebibliography}{Duc18}
\bibitem[Ber93]{Ber93}
V.~G.~Berkovich,
\emph{\'Etale cohomology for non-Archimedean analytic spaces},
Publ. Math. IHES \textbf{78} (1993), 5--161.
\url{https://www.numdam.org/item/PMIHES_1993__78__5_0.pdf}.

\bibitem[Bosch08]{Bosch08}
S.~Bosch, \emph{Lectures on Formal and Rigid Geometry},
author manuscript, April 2008. The numbering and printed pages
used here are those of this manuscript, not the 2014 book.
\url{https://www.uni-muenster.de/IVV5WS/WebHop/user/weckerm/heft378.pdf}.

\bibitem[CLD12]{CLD12}
A.~Chambert-Loir and A.~Ducros,
\emph{Formes diff\'erentielles r\'eelles et courants sur les espaces
de Berkovich}, arXiv:1204.6277v1 (2012).
\url{https://arxiv.org/pdf/1204.6277v1}.

\bibitem[CLD25]{CLD25}
A.~Chambert-Loir and A.~Ducros,
\emph{Formes diff\'erentielles r\'eelles et courants sur les espaces
de Berkovich}, arXiv:1204.6277v2 (2025), 608 pp.
\url{https://arxiv.org/pdf/1204.6277v2}.

\bibitem[dJ96]{dJ96}
A.~J.~de Jong, \emph{Smoothness, semi-stability and alterations},
Publ. Math. IHES \textbf{83} (1996), 51--93.
\url{https://www.numdam.org/item/PMIHES_1996__83__51_0.pdf}.

\bibitem[Duc12]{Duc12}
A.~Ducros, \emph{Espaces de Berkovich, polytopes, squelettes
et th\'eorie des mod\`eles}, Confluentes Math.
\textbf{4} (2012), no.~4, 1250007, 57 pp.
\url{https://www.numdam.org/item/10.1142/S1793744212500077.pdf}.

\bibitem[Duc18]{DucrosFamilies}
A.~Ducros, \emph{Families of Berkovich Spaces},
Ast\'erisque \textbf{400} (2018), vii+262 pp.
\url{https://webusers.imj-prg.fr/~antoine.ducros/Families-Berko.pdf}.

\bibitem[GR25]{GR25}
W.~Gubler and J.~Rabinoff,
\emph{Classically psh and pluriharmonic functions on Berkovich spaces},
arXiv:2506.13548v2, September 17, 2025.
\url{https://gubler.app.uni-regensburg.de/Publikationen/2506.13548.pdf}.

\bibitem[Mih24]{Mih24}
A.~Mihatsch, \emph{$\delta$-Forms on Lubin--Tate Spaces},
arXiv:2112.10018v4, September 7, 2024.
\url{https://arxiv.org/pdf/2112.10018v4}.

\bibitem[PT20]{PT20}
J.~Poineau and D.~Turchetti,
\emph{Berkovich curves and Schottky uniformization},
arXiv:2010.09043v1 (2020).
\url{https://arxiv.org/pdf/2010.09043}.

\bibitem[Stacks]{Stacks}
The Stacks Project Authors, \emph{The Stacks Project},
Tags 00R4, 0AGX, 07NY, and 0AVB, accessed November 10, 2026.
\url{https://stacks.math.columbia.edu}.

\bibitem[Thu05]{Thu05}
A.~Thuillier, \emph{Th\'eorie du potentiel sur les courbes en
g\'eom\'etrie analytique non archim\'edienne. Applications
\`a la th\'eorie d'Arakelov}, Ph.D. thesis, Universit\'e de
Rennes~1, 2005.
\url{https://theses.hal.science/tel-00010990}.
The source-access ledger distinguishes the directly checked
restatements in Wanner from the original thesis locations.

\bibitem[Wan19]{Wan19}
V.~Wanner, \emph{Comparison of two notions of subharmonicity
on non-archimedean curves}, Math. Z. \textbf{293} (2019),
443--474. The checked preprint is arXiv:1801.04713v3,
November 8, 2018.
\url{https://arxiv.org/pdf/1801.04713}.

\bibitem[Wan22]{Wan22}
V.~Wanner, \emph{Energy Minimization Principle for
non-archimedean curves}, J. Th\'eor. Nombres Bordeaux
\textbf{34} (2022), no.~1, 1--39.
\url{https://www.numdam.org/item/JTNB_2022__34_1_1_0.pdf}.
\end{thebibliography}

% END source: references.tex

\end{document}
